\chapter{外部シースの物理モデル}
\label{ch:sheath_model}

第\ref{ch:sheath}章の closure は、計算箱の外にある平面・定常のシースを解き、その解から上端面の条件を作る。
本章では、この外部シースのモデルを BEACH との結合から切り離して説明する。
対象は、太陽風のイオンと電子、表面から放出される光電子がつくる 1 次元・無衝突のシースである。
最初に仮定と粒子の軌道を定め、各粒子集団の密度と流束を電位の関数として導く。
次に Poisson 方程式の第一積分から電位構造を Type A、B、C に分け、零電流の閉じ方、物理解の条件、数値アルゴリズムを示す。
最後に、太陽風と月面の典型条件で解いた数値例により、電位構造と粒子集団を可視化する。

\section{対象と仮定}
\label{sec:sm_assumptions}
座標 $z$ を表面（壁、$z=0$）から上流へ向けてとり、上流 $z\to\infty$ の電位を 0、壁の電位を $\phi_0$ とする。
電場は $E=E_z=-\dd\phi/\dd z$ で、$E>0$ は壁から上流へ向かう。温度 $T_e$、$T_i$、$T_{\mathrm{pe}}$ は eV 単位で表し、
電位を温度で割った量（例えば $\phi/T_e$）は無次元である。$e>0$ は素電荷である。モデルは次を仮定する。

\begin{enumerate}
\item \textbf{平面・1 次元：}壁の曲率半径と凹凸の尺度が Debye 長より十分大きく、量は $z$ だけに依存する。
  接線方向の速度は保存するので、法線速度 $v_z$ の分布だけを扱う。
\item \textbf{定常：}照射とプラズマ条件の変化は、粒子がシースを横切る時間（電子で数 $\mu$s）より十分遅い。
\item \textbf{無衝突・非磁化：}平均自由行程と Larmor 半径は Debye 長より十分長い。粒子はシース内でエネルギーを保存する。
\item \textbf{半無限の上流：}$z\to\infty$ でプラズマは準中性になり、電場は 0 に近づく。上流の入射分布は固定した物理条件であり、
  解の途中で入射速度や電子の drift を変更しない。
\item \textbf{吸収する壁：}壁に達した太陽風の電子とイオン、および壁へ戻った光電子はすべて吸収される。
  二次電子放出、後方散乱、表面伝導は含めない。
\item \textbf{三つの粒子集団：}単一価のイオンは流体として、太陽風電子と光電子は運動論的に扱う。
\end{enumerate}

表\ref{tab:sm_scales}に、本章の数値例（表\ref{tab:sm_example_inputs}）での尺度を示す。
シースの厚さは Debye 長の数倍（数十 m）であり、Larmor 半径と平均自由行程はそれより桁違いに長い。
イオンの入射エネルギーは電子温度の約 70 倍で、イオンは壁の電位でほとんど減速されない。

\begin{table}[htbp]
\centering
\caption{数値例の条件での尺度。$n_{i\infty}=8.7\,\mathrm{cm^{-3}}$、$T_e=12$ eV、$T_{\mathrm{pe}}=2.2$ eV、磁場 5 nT とした。}
\label{tab:sm_scales}
\begin{tabular}{lr}
\toprule
量 & 値 \\
\midrule
太陽風電子の Debye 長 $\lambda_{D,e}=\sqrt{\eps T_e/(e\,n_{i\infty})}$ & 8.7 m \\
光電子の Debye 長 $\sqrt{\eps T_{\mathrm{pe}}/(e\,n_{\mathrm{pe,ref}})}$（$n_{\mathrm{pe,ref}}=64\,\mathrm{cm^{-3}}$） & 1.4 m \\
電子の Larmor 半径（12 eV、5 nT） & 約 2 km \\
光電子の Larmor 半径（2.2 eV、5 nT） & 約 1 km \\
イオンの入射エネルギー $K=m_iu_i^2/(2e)$（$u_i=405$ km/s） & 0.86 keV \\
電子が 10 m を横切る時間 & 約 5 $\mu$s \\
\bottomrule
\end{tabular}
\end{table}

この形のシースは、光を受ける板と希薄プラズマの運動論的な解析\citep{GuernseyFu1970}、宇宙の表面上の光電子シースとダスト浮上の解析\citep{Nitter1998}、
月面の PIC 計算\citep{PoppeHoranyi2010}で調べられてきた。本モデルの粒子集団と三つの型は Zhao らの定式化に従い\citep{Zhao2020}、
イオンの有限温度は Sana と Mishra の流体モデルに従う\citep{SanaMishra2026}。

\section{軌道と粒子集団}
\label{sec:sm_orbits}

\subsection{エネルギー保存と速度領域}
定常・無衝突の 1 次元 Vlasov 方程式
\begin{equation}
 v_z\frac{\partial f}{\partial z}+\frac{q}{m}E\frac{\partial f}{\partial v_z}=0
\end{equation}
の解は、法線方向の全エネルギー $W=\tfrac12 mv_z^2+q\phi(z)$ と軌道の種類だけの関数である（Liouville の定理）。
したがって、上流で入射する分布と壁で放出する分布を与えれば、各位置の分布は境界の分布をエネルギー保存で写したものになる。
写せない速度、つまり境界から到達できない速度の分布は 0 である。壁は吸収するので、壁から上流へ向かう太陽風電子はなく、
上流から入射する光電子もない。

電子（$q=-e$）は $W=-e\phi(z)$ の点で折り返す。電子にとっての障壁は経路上の最小電位
\begin{equation}
 \phi_{\min}=\min_{z\ge0}\phi(z)=
 \begin{cases}\phi_m&\text{(Type A、内部の極小)}\\0&\text{(Type B)}\\\phi_0&\text{(Type C)}\end{cases}
\end{equation}
である。上流での法線エネルギーが $eT_ea^2$（$a$ は後述の規格化速さ）の太陽風電子は、$a\ge b\equiv\sqrt{-\phi_{\min}/T_e}$ なら壁に達して吸収され、
$a<b$ なら障壁の手前で反射して上流へ戻る。壁での法線エネルギーが $eK$ の光電子は、$K\ge B_{\mathrm{pe}}\equiv\phi_0-\phi_{\min}$ なら上流へ脱出し、
それ以下なら折り返して壁へ戻る。反射する太陽風電子は極小より上流側だけを、捕捉される光電子は極小より壁側だけを往復する。
このため、Type A の電荷密度は同じ電位でも極小の両側で異なる（表\ref{tab:sm_populations}）。
図\ref{fig:sm_populations}に、Type A の数値解で電子のエネルギー図と軌道の分類を示す。

\begin{table}[htbp]
\centering
\caption{区間ごとに存在する電子の集団。「通過」は壁へ向かう一方向、「反射」「捕捉」は往復する集団である。}
\label{tab:sm_populations}
\begin{tabular}{lll}
\toprule
区間 & 太陽風電子 & 光電子 \\
\midrule
Type A の下側（壁から極小まで） & 通過 & 脱出、捕捉 \\
Type A の上側（極小から上流まで） & 通過、反射 & 脱出 \\
Type B（全域） & 通過 & 脱出、捕捉 \\
Type C（全域） & 通過、反射 & 脱出（障壁 0） \\
\bottomrule
\end{tabular}
\end{table}

% Generated by tools/sheath_model_figures.py from profile_A.csv; edit the generator, not this file.
\begin{figure}[htbp]
\centering
\begin{tikzpicture}[beach diagram]
\fill[black!4] (0,0) rectangle (3.801,7.6);
\draw[line width=0.6pt] (0,0) rectangle (11.5,7.6);
\draw[line width=0.6pt] (0,0)--++(0,0.1);
\node[below,font=\scriptsize\sffamily] at (0,0) {$0$};
\draw[line width=0.6pt] (1.278,0)--++(0,0.1);
\node[below,font=\scriptsize\sffamily] at (1.278,0) {$5$};
\draw[line width=0.6pt] (2.556,0)--++(0,0.1);
\node[below,font=\scriptsize\sffamily] at (2.556,0) {$10$};
\draw[line width=0.6pt] (3.833,0)--++(0,0.1);
\node[below,font=\scriptsize\sffamily] at (3.833,0) {$15$};
\draw[line width=0.6pt] (5.111,0)--++(0,0.1);
\node[below,font=\scriptsize\sffamily] at (5.111,0) {$20$};
\draw[line width=0.6pt] (6.389,0)--++(0,0.1);
\node[below,font=\scriptsize\sffamily] at (6.389,0) {$25$};
\draw[line width=0.6pt] (7.667,0)--++(0,0.1);
\node[below,font=\scriptsize\sffamily] at (7.667,0) {$30$};
\draw[line width=0.6pt] (8.944,0)--++(0,0.1);
\node[below,font=\scriptsize\sffamily] at (8.944,0) {$35$};
\draw[line width=0.6pt] (10.222,0)--++(0,0.1);
\node[below,font=\scriptsize\sffamily] at (10.222,0) {$40$};
\draw[line width=0.6pt] (11.5,0)--++(0,0.1);
\node[below,font=\scriptsize\sffamily] at (11.5,0) {$45$};
\draw[line width=0.6pt] (0,0.386)--++(0.1,0);
\node[left,font=\scriptsize\sffamily] at (0,0.386) {$-4$};
\draw[line width=0.6pt] (0,1.675)--++(0.1,0);
\node[left,font=\scriptsize\sffamily] at (0,1.675) {$-3$};
\draw[line width=0.6pt] (0,2.963)--++(0.1,0);
\node[left,font=\scriptsize\sffamily] at (0,2.963) {$-2$};
\draw[line width=0.6pt] (0,4.251)--++(0.1,0);
\node[left,font=\scriptsize\sffamily] at (0,4.251) {$-1$};
\draw[line width=0.6pt] (0,5.539)--++(0.1,0);
\node[left,font=\scriptsize\sffamily] at (0,5.539) {$0$};
\draw[line width=0.6pt] (0,6.827)--++(0.1,0);
\node[left,font=\scriptsize\sffamily] at (0,6.827) {$1$};
\node[font=\small\sffamily] at (5.75,-0.75) {壁からの距離 $z$ [m]};
\node[rotate=90,font=\small\sffamily] at (-1.05,3.8) {電子のポテンシャルエネルギー $U=-e\phi$ [eV]};
\begin{scope}
\clip (0,0) rectangle (11.5,7.6);
\draw[black!45,dashed,line width=0.5pt] (0,5.539)--(11.5,5.539);
\draw[BeachInk,line width=1.4pt] plot coordinates {(0,0.592) (0.002,0.608) (0.004,0.625) (0.006,0.641) (0.008,0.657) (0.01,0.673) (0.012,0.689) (0.039,0.897) (0.064,1.086) (0.09,1.266) (0.116,1.438) (0.141,1.601) (0.167,1.757) (0.193,1.908) (0.219,2.052) (0.245,2.188) (0.27,2.32) (0.296,2.445) (0.322,2.566) (0.348,2.682) (0.374,2.795) (0.399,2.901) (0.425,3.004) (0.451,3.103) (0.476,3.198) (0.502,3.29) (0.528,3.378) (0.553,3.464) (0.579,3.547) (0.605,3.628) (0.631,3.705) (0.657,3.779) (0.683,3.851) (0.709,3.921) (0.735,3.988) (0.761,4.053) (0.786,4.115) (0.812,4.176) (0.838,4.235) (0.864,4.292) (0.89,4.347) (0.916,4.4) (0.942,4.452) (0.968,4.501) (0.994,4.55) (1.019,4.596) (1.046,4.642) (1.071,4.686) (1.097,4.729) (1.123,4.77) (1.149,4.81) (1.175,4.848) (1.201,4.886) (1.227,4.923) (1.253,4.958) (1.279,4.993) (1.305,5.026) (1.331,5.058) (1.358,5.09) (1.384,5.12) (1.41,5.149) (1.436,5.177) (1.462,5.205) (1.488,5.232) (1.514,5.258) (1.54,5.283) (1.566,5.308) (1.592,5.331) (1.618,5.354) (1.644,5.376) (1.67,5.397) (1.695,5.418) (1.722,5.439) (1.748,5.459) (1.774,5.477) (1.8,5.496) (1.826,5.514) (1.852,5.531) (1.878,5.548) (1.904,5.564) (1.931,5.58) (1.956,5.595) (1.982,5.61) (2.008,5.624) (2.035,5.638) (2.061,5.652) (2.087,5.665) (2.113,5.677) (2.14,5.69) (2.167,5.702) (2.192,5.713) (2.218,5.724) (2.245,5.735) (2.271,5.745) (2.298,5.755) (2.324,5.765) (2.35,5.774) (2.376,5.783) (2.403,5.792) (2.43,5.801) (2.455,5.808) (2.481,5.816) (2.508,5.824) (2.534,5.831) (2.561,5.838) (2.588,5.845) (2.614,5.852) (2.64,5.858) (2.666,5.864) (2.693,5.87) (2.72,5.875) (2.747,5.881) (2.773,5.886) (2.799,5.891) (2.826,5.895) (2.852,5.9) (2.879,5.904) (2.907,5.909) (2.932,5.912) (2.958,5.916) (2.984,5.92) (3.011,5.923) (3.038,5.926) (3.065,5.93) (3.092,5.933) (3.118,5.935) (3.144,5.938) (3.17,5.94) (3.196,5.943) (3.223,5.945) (3.25,5.947) (3.277,5.949) (3.305,5.951) (3.33,5.953) (3.356,5.954) (3.382,5.956) (3.409,5.957) (3.436,5.958) (3.463,5.96) (3.49,5.961) (3.517,5.962) (3.545,5.963) (3.571,5.963) (3.597,5.964) (3.623,5.964) (3.649,5.965) (3.676,5.965) (3.702,5.966) (3.728,5.966) (3.756,5.966) (3.777,5.966) (3.824,5.966) (3.851,5.966) (3.883,5.966) (3.915,5.966) (3.941,5.965) (3.97,5.965) (4.001,5.964) (4.028,5.964) (4.057,5.963) (4.084,5.962) (4.112,5.962) (4.139,5.961) (4.167,5.96) (4.194,5.959) (4.222,5.958) (4.249,5.957) (4.275,5.956) (4.303,5.955) (4.329,5.954) (4.357,5.953) (4.384,5.952) (4.41,5.951) (4.437,5.95) (4.464,5.948) (4.49,5.947) (4.518,5.946) (4.545,5.944) (4.571,5.943) (4.596,5.942) (4.623,5.94) (4.65,5.939) (4.676,5.937) (4.703,5.936) (4.73,5.934) (4.757,5.933) (4.783,5.931) (4.808,5.93) (4.835,5.928) (4.862,5.927) (4.888,5.925) (4.914,5.923) (4.94,5.922) (4.967,5.92) (4.993,5.918) (5.02,5.917) (5.046,5.915) (5.071,5.913) (5.097,5.912) (5.124,5.91) (5.15,5.908) (5.177,5.906) (5.203,5.905) (5.229,5.903) (5.255,5.901) (5.28,5.899) (5.307,5.897) (5.334,5.896) (5.36,5.894) (5.386,5.892) (5.413,5.89) (5.439,5.888) (5.465,5.886) (5.49,5.885) (5.516,5.883) (5.543,5.881) (5.57,5.879) (5.596,5.877) (5.623,5.875) (5.649,5.873) (5.675,5.871) (5.701,5.87) (5.728,5.868) (5.754,5.866) (5.78,5.864) (5.805,5.862) (5.831,5.86) (5.857,5.858) (5.883,5.857) (5.908,5.855) (5.935,5.853) (5.962,5.851) (5.988,5.849) (6.015,5.847) (6.042,5.845) (6.068,5.843) (6.095,5.842) (6.121,5.84) (6.147,5.838) (6.174,5.836) (6.2,5.834) (6.226,5.832) (6.253,5.83) (6.279,5.829) (6.305,5.827) (6.331,5.825) (6.358,5.823) (6.384,5.821) (6.41,5.82) (6.436,5.818) (6.462,5.816) (6.488,5.814) (6.514,5.812) (6.54,5.811) (6.567,5.809) (6.593,5.807) (6.619,5.805) (6.645,5.804) (6.671,5.802) (6.697,5.8) (6.723,5.798) (6.749,5.797) (6.775,5.795) (6.802,5.793) (6.828,5.791) (6.854,5.79) (6.88,5.788) (6.906,5.786) (6.933,5.785) (6.959,5.783) (6.985,5.781) (7.011,5.78) (7.038,5.778) (7.064,5.776) (7.09,5.775) (7.117,5.773) (7.143,5.771) (7.17,5.77) (7.196,5.768) (7.223,5.766) (7.249,5.765) (7.276,5.763) (7.302,5.762) (7.329,5.76) (7.356,5.758) (7.382,5.757) (7.409,5.755) (7.436,5.754) (7.461,5.752) (7.487,5.751) (7.513,5.749) (7.538,5.748) (7.564,5.746) (7.59,5.745) (7.616,5.743) (7.642,5.742) (7.668,5.74) (7.693,5.739) (7.72,5.738) (7.746,5.736) (7.772,5.735) (7.798,5.733) (7.824,5.732) (7.85,5.73) (7.877,5.729) (7.903,5.728) (7.929,5.726) (7.956,5.725) (7.982,5.723) (8.009,5.722) (8.036,5.721) (8.062,5.719) (8.089,5.718) (8.116,5.716) (8.143,5.715) (8.168,5.714) (8.194,5.713) (8.22,5.711) (8.246,5.71) (8.272,5.709) (8.297,5.707) (8.323,5.706) (8.349,5.705) (8.376,5.704) (8.402,5.702) (8.428,5.701) (8.454,5.7) (8.481,5.699) (8.507,5.698) (8.534,5.696) (8.56,5.695) (8.587,5.694) (8.613,5.693) (8.64,5.692) (8.667,5.69) (8.694,5.689) (8.72,5.688) (8.745,5.687) (8.771,5.686) (8.797,5.685) (8.823,5.684) (8.849,5.682) (8.875,5.681) (8.901,5.68) (8.927,5.679) (8.953,5.678) (8.979,5.677) (9.006,5.676) (9.032,5.675) (9.059,5.674) (9.085,5.673) (9.112,5.672) (9.139,5.671) (9.166,5.67) (9.193,5.669) (9.22,5.667) (9.245,5.667) (9.271,5.666) (9.297,5.665) (9.322,5.664) (9.348,5.663) (9.374,5.662) (9.4,5.661) (9.426,5.66) (9.453,5.659) (9.479,5.658) (9.505,5.657) (9.532,5.656) (9.558,5.655) (9.585,5.654) (9.612,5.653) (9.638,5.652) (9.665,5.651) (9.692,5.65) (9.719,5.65) (9.746,5.649) (9.772,5.648) (9.798,5.647) (9.824,5.646) (9.849,5.645) (9.875,5.644) (9.901,5.644) (9.928,5.643) (9.954,5.642) (9.98,5.641) (10.006,5.64) (10.033,5.639) (10.059,5.639) (10.086,5.638) (10.113,5.637) (10.14,5.636) (10.167,5.635) (10.194,5.635) (10.221,5.634) (10.248,5.633) (10.273,5.632) (10.299,5.632) (10.325,5.631) (10.351,5.63) (10.377,5.629) (10.403,5.629) (10.429,5.628) (10.455,5.627) (10.481,5.627) (10.508,5.626) (10.534,5.625) (10.561,5.624) (10.587,5.624) (10.614,5.623) (10.641,5.622) (10.668,5.622) (10.695,5.621) (10.722,5.62) (10.749,5.62) (10.776,5.619) (10.802,5.618) (10.827,5.618) (10.853,5.617) (10.879,5.616) (10.905,5.616) (10.931,5.615) (10.957,5.615) (10.983,5.614) (11.01,5.613) (11.036,5.613) (11.062,5.612) (11.089,5.611) (11.116,5.611) (11.142,5.61) (11.169,5.61) (11.196,5.609) (11.223,5.609) (11.25,5.608) (11.278,5.607) (11.305,5.607) (11.333,5.606) (11.358,5.606) (11.384,5.605) (11.41,5.605) (11.435,5.604) (11.461,5.604) (11.487,5.603) (11.514,5.602) (11.54,5.602) (11.566,5.601) (11.593,5.601) (11.619,5.6) (11.646,5.6) (11.672,5.599) (11.699,5.599) (11.726,5.598) (11.753,5.598) (11.78,5.597) (11.807,5.597) (11.835,5.596) (11.862,5.596) (11.89,5.595) (11.917,5.595) (11.943,5.594) (11.968,5.594) (11.994,5.594) (12.02,5.593) (12.046,5.593) (12.072,5.592) (12.098,5.592) (12.124,5.591) (12.151,5.591) (12.177,5.59) (12.204,5.59) (12.23,5.59) (12.257,5.589) (12.284,5.589) (12.311,5.588) (12.338,5.588) (12.365,5.588) (12.392,5.587) (12.42,5.587) (12.447,5.586) (12.474,5.586) (12.502,5.585) (12.53,5.585) (12.558,5.585) (12.583,5.584) (12.609,5.584) (12.635,5.584) (12.661,5.583) (12.687,5.583) (12.713,5.582) (12.739,5.582) (12.765,5.582) (12.792,5.581) (12.818,5.581) (12.845,5.581) (12.872,5.58) (12.898,5.58) (12.925,5.58) (12.952,5.579) (12.979,5.579) (13.006,5.579) (13.034,5.578) (13.061,5.578) (13.089,5.578) (13.116,5.577) (13.144,5.577) (13.172,5.577) (13.2,5.576) (13.228,5.576) (13.253,5.576) (13.279,5.575) (13.305,5.575) (13.331,5.575) (13.356,5.574) (13.383,5.574) (13.409,5.574) (13.435,5.574) (13.461,5.573) (13.488,5.573) (13.514,5.573) (13.541,5.572) (13.567,5.572) (13.594,5.572) (13.621,5.572) (13.648,5.571) (13.675,5.571) (13.703,5.571) (13.73,5.57) (13.757,5.57) (13.785,5.57) (13.812,5.57) (13.84,5.569) (13.868,5.569) (13.896,5.569) (13.924,5.569) (13.952,5.568) (13.981,5.568) (14.009,5.568) (14.035,5.567) (14.06,5.567) (14.086,5.567) (14.112,5.567) (14.138,5.567) (14.164,5.566) (14.19,5.566) (14.217,5.566) (14.243,5.566) (14.27,5.565) (14.296,5.565) (14.323,5.565) (14.35,5.565) (14.376,5.565) (14.403,5.564) (14.43,5.564) (14.458,5.564) (14.485,5.564) (14.512,5.563) (14.54,5.563) (14.567,5.563) (14.595,5.563) (14.623,5.563) (14.651,5.562) (14.679,5.562) (14.707,5.562) (14.735,5.562) (14.763,5.562) (14.792,5.561) (14.82,5.561) (14.849,5.561) (14.877,5.561) (14.903,5.561) (14.929,5.56) (14.955,5.56) (14.98,5.56) (15.006,5.56) (15.032,5.56) (15.059,5.559) (15.085,5.559) (15.111,5.559) (15.138,5.559) (15.164,5.559) (15.191,5.559) (15.217,5.558) (15.244,5.558) (15.271,5.558) (15.298,5.558) (15.325,5.558) (15.352,5.558) (15.38,5.557) (15.407,5.557) (15.434,5.557) (15.462,5.557) (15.49,5.557) (15.517,5.557) (15.545,5.557) (15.573,5.556) (15.601,5.556) (15.629,5.556) (15.658,5.556) (15.686,5.556) (15.714,5.556) (15.743,5.555) (15.772,5.555) (15.8,5.555) (15.829,5.555) (15.858,5.555) (15.887,5.555) (15.916,5.555) (15.942,5.554) (15.968,5.554) (15.993,5.554) (16.019,5.554) (16.045,5.554) (16.071,5.554) (16.097,5.554) (16.123,5.554) (16.15,5.553) (16.176,5.553) (16.202,5.553) (16.229,5.553) (16.255,5.553) (16.282,5.553) (16.309,5.553) (16.336,5.553) (16.362,5.552) (16.389,5.552) (16.416,5.552) (16.444,5.552) (16.471,5.552) (16.498,5.552) (16.526,5.552) (16.553,5.552) (16.581,5.552) (16.608,5.551) (16.636,5.551) (16.664,5.551) (16.692,5.551) (16.72,5.551) (16.748,5.551) (16.776,5.551) (16.804,5.551) (16.833,5.551) (16.861,5.55) (16.889,5.55) (16.918,5.55) (16.947,5.55) (16.975,5.55) (17.004,5.55) (17.033,5.55) (17.062,5.55) (17.091,5.55) (17.12,5.549) (17.15,5.549) (17.179,5.549) (17.208,5.549) (17.238,5.549) (17.268,5.549) (17.297,5.549) (17.327,5.549) (17.353,5.549) (17.378,5.549) (17.404,5.549) (17.43,5.548) (17.455,5.548) (17.481,5.548) (17.507,5.548) (17.533,5.548) (17.559,5.548) (17.585,5.548) (17.611,5.548) (17.638,5.548) (17.664,5.548) (17.69,5.548) (17.717,5.548) (17.743,5.548) (17.77,5.547) (17.796,5.547) (17.823,5.547) (17.849,5.547) (17.876,5.547) (17.903,5.547) (17.929,5.547) (17.956,5.547) (17.983,5.547) (18.01,5.547) (18.037,5.547) (18.064,5.547) (18.091,5.547) (18.118,5.547) (18.145,5.547) (18.173,5.546) (18.2,5.546) (18.227,5.546) (18.255,5.546) (18.282,5.546) (18.309,5.546) (18.337,5.546) (18.364,5.546) (18.392,5.546) (18.419,5.546) (18.447,5.546) (18.474,5.546) (18.502,5.546) (18.53,5.546) (18.557,5.546) (18.585,5.545) (18.613,5.545) (18.641,5.545) (18.668,5.545) (18.696,5.545) (18.724,5.545) (18.752,5.545) (18.78,5.545) (18.807,5.545) (18.835,5.545) (18.863,5.545) (18.891,5.545) (18.919,5.545) (18.947,5.545) (18.975,5.545) (19.002,5.545) (19.03,5.545) (19.058,5.545) (19.086,5.545) (19.114,5.544) (19.141,5.544) (19.169,5.544) (19.197,5.544) (19.225,5.544) (19.252,5.544) (19.28,5.544) (19.308,5.544) (19.335,5.544) (19.363,5.544) (19.39,5.544) (19.418,5.544) (19.445,5.544) (19.472,5.544) (19.5,5.544) (19.527,5.544) (19.554,5.544) (19.581,5.544) (19.608,5.544) (19.635,5.544) (19.662,5.544) (19.688,5.544) (19.715,5.543) (19.742,5.543) (19.768,5.543) (19.794,5.543) (19.821,5.543) (19.847,5.543) (19.873,5.543) (19.899,5.543) (19.924,5.543) (19.95,5.543) (19.98,5.543) (20.009,5.543) (20.039,5.543) (20.068,5.543) (20.097,5.543) (20.125,5.543) (20.154,5.543) (20.182,5.543) (20.21,5.543) (20.237,5.543) (20.265,5.543) (20.292,5.543) (20.318,5.543) (20.345,5.543) (20.371,5.543) (20.397,5.543) (20.426,5.543) (20.455,5.542) (20.483,5.542) (20.511,5.542) (20.538,5.542) (20.565,5.542) (20.591,5.542) (20.616,5.542) (20.645,5.542) (20.672,5.542) (20.699,5.542) (20.725,5.542) (20.753,5.542) (20.78,5.542) (20.805,5.542) (20.833,5.542) (20.859,5.542) (20.885,5.542) (20.913,5.542) (20.94,5.542) (20.966,5.542) (20.993,5.542) (21.019,5.542) (21.045,5.542) (21.071,5.542) (21.074,5.542)};
\end{scope}
\draw[BeachBlue,-{Latex[length=2mm]},line width=1.0pt] (11.5,7.02)--(0.038,7.02);
\node[font=\scriptsize\sffamily,text=BeachBlue] at (7.667,7.304) {太陽風電子：障壁を越えて壁で吸収};
\draw[BeachBlue,dashed,-{Latex[length=2mm]},line width=1.0pt] (11.5,5.797)--(6.749,5.797);
\node[beach dot,fill=BeachBlue] at (6.749,5.797) {};
\node[anchor=west,font=\scriptsize\sffamily,text=BeachBlue] at (7.005,5.088) {太陽風電子：上側で反射};
\draw[BeachOrange,-{Latex[length=2mm]},line width=1.0pt] (0.038,6.338)--(11.5,6.338);
\node[font=\scriptsize\sffamily,text=BeachOrange] at (7.922,6.621) {光電子：障壁を越えて上流へ脱出};
\draw[BeachRed,-{Latex[length=2mm]},line width=1.0pt] (0.038,4.38)--(0.906,4.38);
\node[beach dot,fill=BeachRed] at (0.906,4.38) {};
\node[anchor=west,font=\scriptsize\sffamily,text=BeachRed] at (1.059,4.38) {光電子：下側で折り返し、壁へ戻る};
\draw[BeachInk,dotted,line width=0.6pt] (0,5.966)--(11.5,5.966);
\node[beach dot] at (3.801,5.966) {};
\node[anchor=west,font=\scriptsize\sffamily] at (4.056,6.152) {障壁 $-e\phi_m=0.33$ eV};
\node[beach dot] at (0,0.592) {};
\node[anchor=west,font=\scriptsize\sffamily] at (0.153,0.657) {壁 $-e\phi_0=-3.84$ eV};
\node[font=\small\sffamily] at (2.428,1.288) {下側（壁〜極小）};
\node[font=\small\sffamily] at (7.667,1.288) {上側（極小〜上流）};
\end{tikzpicture}
\caption[Type A の電子のエネルギー図と軌道の分類]{Type A の電子のポテンシャルエネルギー $U(z)=-e\phi(z)$ と軌道の分類。黒の曲線は図\ref{fig:sm_profiles}の Type A の数値解で、色の水平線は保存する全エネルギー $W$ の例である（$W$ の値は説明用）。電子は $U(z)=W$ の点で折り返す。太陽風電子は $W>-e\phi_m$ なら壁に達し、そうでなければ上側で反射する。光電子は壁での法線エネルギーが $e(\phi_0-\phi_m)$ を超えれば脱出し、それ以下なら下側で折り返して壁に戻る。この二つの閾値が、各位置に存在する速度領域を決める。}
\label{fig:sm_populations}
\end{figure}


\subsection{太陽風電子}
上流で壁へ向かう太陽風電子は、法線方向に速さ $v_d$ で drift した Maxwell 分布である。
$v_{te}=\sqrt{2eT_e/m_e}$、$u=v_d/v_{te}$、上流での法線の速さを $a=|v_z|/v_{te}$ とすると、上流の入射分布は
\begin{equation}
 f_\infty(v_z)=\frac{N_e}{\sqrt\pi\,v_{te}}\exp\left[-(a-u)^2\right],\qquad v_z<0
 \label{eq:sm_e_upstream}
\end{equation}
である。$N_e$ は入射分布の規格化であり、上流の実際の電子密度ではない。上流の太陽風電子は入射電子と反射して戻る電子の和であり、壁で吸収された高速電子は戻らないからである。
位置 $z$ での速さを $w=|v_z|/v_{te}$、$\psi=\phi/T_e$、$\psi_{\min}=\phi_{\min}/T_e$ とすると、エネルギー保存 $w^2=a^2+\psi$ と Liouville の定理から
\begin{equation}
 f(z,v_z)=\frac{N_e}{\sqrt\pi\,v_{te}}\exp\left[-\left(\sqrt{w^2-\psi}-u\right)^2\right]
\end{equation}
を得る。壁へ到達する通過電子は $w\ge w_c\equiv\sqrt{\psi-\psi_{\min}}$ の内向きの速度を、反射電子は $0\le w<w_c$ の両方向を占める。
密度はこの分布の速度積分
\begin{align}
 n_{e,\mathrm{pass}}&=\frac{N_e}{\sqrt\pi}\int_{w_c}^{\infty}e^{-(\sqrt{w^2-\psi}-u)^2}\dd w,&
 n_{e,\mathrm{refl}}&=\frac{2N_e}{\sqrt\pi}\int_{0}^{w_c}e^{-(\sqrt{w^2-\psi}-u)^2}\dd w
 \label{eq:sm_e_density}
\end{align}
であり、反射電子は Type A の上側と Type C だけにある。drift が 0 のとき、積分は閉じた形
\begin{equation}
 n_{e,\mathrm{pass}}=\frac{N_e}2e^{\psi}\operatorname{erfc}\sqrt{\psi-\psi_{\min}},\qquad
 n_{e,\mathrm{refl}}=N_e\,e^{\psi}\operatorname{erf}\sqrt{\psi-\psi_{\min}}
 \label{eq:sm_e_density_zero}
\end{equation}
になる。障壁ですべての電子が反射されるなら密度は Boltzmann 分布 $N_ee^{\psi}$ になるが、
吸収壁では障壁を越えた電子が戻らないので、外向きの高速成分が欠けて密度はそれより小さい。

流束は通過電子だけが担う。$w\,\dd w=a\,\dd a$ と変数変換すると
\begin{equation}
 \Gamma_e=\frac{N_ev_{te}}{2\sqrt\pi}\left[e^{-(b-u)^2}+\sqrt\pi\,u\operatorname{erfc}(b-u)\right]
 \label{eq:sm_e_flux}
\end{equation}
となり、位置によらない（連続の式）。上流 $\psi=0$ での密度は
\begin{equation}
 n_e(0)=\frac{N_e}2\left[1+2\operatorname{erf}u+\operatorname{erf}(b-u)\right]\quad\text{(A、C)},\qquad
 n_e(0)=\frac{N_e}2\left(1+\operatorname{erf}u\right)\quad\text{(B)}
 \label{eq:sm_e_upstream_density}
\end{equation}
である。光電子がなく drift が 0 の Type C では、上流の準中性から
\begin{equation}
 N_e=\frac{n_{i\infty}}{1-\tfrac12\operatorname{erfc}\sqrt{-\phi_0/T_e}}
 \label{eq:upstream_density}
\end{equation}
となり、$N_e$ は上流のイオン密度 $n_{i\infty}$ より大きい。

\subsection{光電子}
壁は温度 $T_{\mathrm{pe}}$ の半 Maxwell 分布で光電子を放出する。$v_{tp}=\sqrt{2eT_{\mathrm{pe}}/m_e}$、規格化を $N_{\mathrm{pe}}$ とすると、
壁での外向き分布は $f_w(v_z)=N_{\mathrm{pe}}\exp(-v_z^2/v_{tp}^2)/(\sqrt\pi v_{tp})$（$v_z>0$）、放出束は
$\Gamma_{\mathrm{emit}}=N_{\mathrm{pe}}v_{tp}/(2\sqrt\pi)$ である。障壁 $B_{\mathrm{pe}}$ を越える割合が脱出し、
\begin{equation}
 \Gamma_{\mathrm{esc}}=\Gamma_{\mathrm{emit}}\,e^{-B_{\mathrm{pe}}/T_{\mathrm{pe}}},\qquad
 \Gamma_{\mathrm{ret}}=\Gamma_{\mathrm{emit}}-\Gamma_{\mathrm{esc}}
 \label{eq:sm_pe_flux}
\end{equation}
となる。$\chi=(\phi-\phi_{\min})/T_{\mathrm{pe}}$ とすると、脱出する光電子と捕捉される光電子（往路と復路の両方）の密度は
\begin{equation}
 n_{\mathrm{pe,esc}}=\frac{N_{\mathrm{pe}}}2e^{(\phi-\phi_0)/T_{\mathrm{pe}}}\operatorname{erfc}\sqrt\chi,\qquad
 n_{\mathrm{pe,capt}}=N_{\mathrm{pe}}\,e^{(\phi-\phi_0)/T_{\mathrm{pe}}}\operatorname{erf}\sqrt\chi
 \label{eq:sm_pe_density}
\end{equation}
であり、捕捉された光電子は Type A の下側と Type B だけにある。Type C は障壁が 0 なので、放出した光電子はすべて脱出する。

光電子源は、法線エネルギー $K$ の bin $[K_j,K_{j+1})$ ごとの外向き数流束 $\Gamma_j$ でも与えられる。
bin の中では $g_j=\dd\Gamma/\dd K=\Gamma_j/(K_{j+1}-K_j)$ を一定とする。電位 $\phi$ の点での法線エネルギーは
$K-q$（$q=\phi_0-\phi$）、速さは $v_*\sqrt{K-q}$（$v_*=\sqrt{2e/m_e}$）なので、$\dd n=g\,\dd K/(v_*\sqrt{K-q})$ を積分して
\begin{align}
 n_{\mathrm{pe,esc}}&=\sum_j\frac{2g_j}{v_*}\left[\sqrt{K_{j+1}-q}-\sqrt{\max(K_j,B_{\mathrm{pe}},q)-q}\right]_+,\nonumber\\
 n_{\mathrm{pe,capt}}&=\sum_j\frac{4g_j}{v_*}\left[\sqrt{\min(K_{j+1},B_{\mathrm{pe}})-q}-\sqrt{\max(K_j,q)-q}\right]_+
 \label{eq:sm_pe_bins}
\end{align}
を得る。$[\,\cdot\,]_+$ は上限が下限以下のとき 0 とする意味である。脱出束は $\Gamma_{\mathrm{esc}}=\sum_jg_j\,[K_{j+1}-\max(K_j,B_{\mathrm{pe}})]_+$ で、
同じ $g_j$ を密度、流束、Poisson 方程式の積分のすべてに使う。bin の分布を Maxwell 分布へ当てはめる操作はない。
第\ref{ch:sheath}章の外部根の更新は、この形で上端面を出た光電子のスペクトルを渡す。

\subsection{太陽風イオン}
イオンは上流で密度 $n_{i\infty}$、法線速度 $u_i$（壁向きを正）、温度 $T_i$ を持つ流体である。
連続の式 $n_iv=n_{i\infty}u_i$ と、圧力 $p=e fT_in_i$（係数 $f$ は一定、1 は等温、3 は 1 次元断熱）を含む運動方程式
$m_in_iv\,\dd v/\dd z=-\dd p/\dd z-en_i\,\dd\phi/\dd z$ を積分すると、入射エネルギー $K=m_iu_i^2/(2e)$、$a=fT_i$、
密度比 $r=n_i/n_{i\infty}$ に対して
\begin{equation}
 K\left(\frac1{r^2}-1\right)+a\log r+\phi=0
 \label{eq:ion_fluid}
\end{equation}
が成り立つ\citep{SanaMishra2026}。$f$ は同論文の断熱指数 $\psi$ に当たる。上流につながる枝は超音速の根 $2K>a$ から連続する根であり、
$a=0$ では冷たいビーム $r=(1-\phi/K)^{-1/2}$ に戻る。イオンが到達できる最大電位は
\begin{equation}
 \phi_{\mathrm{crit}}=K-\frac a2-\frac a2\log\frac{2K}{a}\qquad(a=0\text{ では }K)
 \label{eq:sm_ion_crit}
\end{equation}
で、経路上の電位がこれを超えるとイオンの流れは途中で止まる。流束は $\Gamma_i=n_{i\infty}u_i$ で一定である。
太陽風では $K$（405 km/s の陽子で約 0.86 keV）が $T_i$ よりはるかに大きく、$T_i\le12$ eV、$f=1,3$ で根の変化は 0.1 mV 以下である。
このモデルはイオンの速度分布そのものは定めない。

\section{Poisson 方程式と電位構造}
\label{sec:sm_poisson}

\subsection{第一積分}
電荷密度を $\rho=e(n_i-n_{e,\mathrm{pass}}-n_{e,\mathrm{refl}}-n_{\mathrm{pe,esc}}-n_{\mathrm{pe,capt}})$ とする。
電位が単調な区間では存在する粒子集団が決まっているので、$\rho$ は $\phi$ だけの関数である。
Poisson 方程式 $\dd^2\phi/\dd z^2=-\rho/\eps$ に $\dd\phi/\dd z$ を掛けると $\tfrac12\,\dd(E^2)/\dd z=-(\rho/\eps)\,\dd\phi/\dd z$ となり、
電位について積分できる。上流で $E=0$ の単調な区間（Type B、C）では
\begin{equation}
 E^2(\phi)=\frac2\eps\int_\phi^0\rho(\psi)\,\dd\psi
 \label{eq:sagdeev}
\end{equation}
である。これは $\phi$ を位置、$z$ を時間とみなした質点のエネルギー保存（Sagdeev の擬ポテンシャル）と同じ形をしている。
Type A では極小 $\phi_m$ で $E=0$ なので、下側と上側の区間 $s$ ごとに
\begin{equation}
 E_s^2(\phi)=-\frac2\eps\int_{\phi_m}^{\phi}\rho_s(\psi)\,\dd\psi,\qquad s\in\{\text{下側},\text{上側}\}
 \label{eq:sm_typeA_first}
\end{equation}
となる。上側は上流でも $E=0$ でなければならないので、接続条件
\begin{equation}
 \int_{\phi_m}^{0}\rho_{\text{上側}}(\psi)\,\dd\psi=0
 \label{eq:sm_connection}
\end{equation}
が加わる。極小が存在するには $\dd^2\phi/\dd z^2>0$、つまり極小で $\rho(\phi_m)<0$ が必要である。
壁の電場は $E_0=\pm\sqrt{E^2(\phi_0)}$ で、符号は Type A と B で正、Type C で負である。

図\ref{fig:sm_first_integral}に、Type A の数値解での $\rho(\phi)$ と $E^2(\phi)$ を示す。
下側では捕捉された光電子が負の電荷を担い、$E^2$ は壁へ向かって単調に増える。
上側の電荷密度は極小の近くで負、上流の近くで正になり、二つの面積が打ち消すことが式\eqref{eq:sm_connection}である。

% Generated by tools/sheath_model_figures.py from first_integral_A.csv; edit the generator, not this file.
\begin{figure}[htbp]
\centering
\begin{tikzpicture}[beach diagram]
\draw[black!10,line width=0.4pt] (0.493,4.6)--(0.493,7.9);
\draw[black!10,line width=0.4pt] (1.979,4.6)--(1.979,7.9);
\draw[black!10,line width=0.4pt] (3.465,4.6)--(3.465,7.9);
\draw[black!10,line width=0.4pt] (4.951,4.6)--(4.951,7.9);
\draw[black!10,line width=0.4pt] (0,5.2)--(6.2,5.2);
\draw[black!10,line width=0.4pt] (0,6.4)--(6.2,6.4);
\draw[black!10,line width=0.4pt] (0,7.6)--(6.2,7.6);
\draw[line width=0.6pt] (0,4.6) rectangle (6.2,7.9);
\draw[line width=0.6pt] (0.493,4.6)--++(0,0.1);
\node[below,font=\scriptsize\sffamily] at (0.493,4.6) {$0$};
\draw[line width=0.6pt] (1.979,4.6)--++(0,0.1);
\node[below,font=\scriptsize\sffamily] at (1.979,4.6) {$1$};
\draw[line width=0.6pt] (3.465,4.6)--++(0,0.1);
\node[below,font=\scriptsize\sffamily] at (3.465,4.6) {$2$};
\draw[line width=0.6pt] (4.951,4.6)--++(0,0.1);
\node[below,font=\scriptsize\sffamily] at (4.951,4.6) {$3$};
\draw[line width=0.6pt] (0,5.2)--++(0.1,0);
\node[left,font=\scriptsize\sffamily] at (0,5.2) {$-40$};
\draw[line width=0.6pt] (0,6.4)--++(0.1,0);
\node[left,font=\scriptsize\sffamily] at (0,6.4) {$-20$};
\draw[line width=0.6pt] (0,7.6)--++(0.1,0);
\node[left,font=\scriptsize\sffamily] at (0,7.6) {$0$};
\node[font=\small\sffamily] at (3.1,4) {$\phi$ [V]};
\node[rotate=90,font=\small\sffamily] at (-1,6.25) {$\rho/e$ [cm$^{-3}$]};
\begin{scope}
\clip (0,4.6) rectangle (6.2,7.9);
\draw[black!45,dashed,line width=0.5pt] (0,7.6)--(6.2,7.6);
\draw[BeachRed,line width=1.2pt] plot coordinates {(0,7.583) (0.013,7.573) (0.028,7.567) (0.043,7.562) (0.062,7.557) (0.078,7.553) (0.097,7.549) (0.117,7.544) (0.14,7.539) (0.155,7.535) (0.172,7.532) (0.19,7.528) (0.208,7.524) (0.228,7.521) (0.248,7.516) (0.269,7.512) (0.291,7.508) (0.314,7.503) (0.338,7.499) (0.362,7.494) (0.387,7.489) (0.414,7.484) (0.441,7.479) (0.469,7.473) (0.498,7.468) (0.512,7.465) (0.527,7.462) (0.543,7.459) (0.558,7.456) (0.574,7.453) (0.589,7.45) (0.605,7.446) (0.622,7.443) (0.638,7.44) (0.655,7.437) (0.672,7.433) (0.689,7.43) (0.706,7.426) (0.724,7.423) (0.742,7.419) (0.76,7.416) (0.778,7.412) (0.796,7.408) (0.815,7.404) (0.834,7.401) (0.853,7.397) (0.872,7.393) (0.891,7.389) (0.911,7.384) (0.931,7.38) (0.951,7.376) (0.971,7.372) (0.992,7.367) (1.013,7.363) (1.034,7.358) (1.055,7.354) (1.076,7.349) (1.098,7.344) (1.12,7.34) (1.142,7.335) (1.164,7.33) (1.187,7.325) (1.209,7.32) (1.232,7.314) (1.256,7.309) (1.279,7.304) (1.302,7.298) (1.326,7.293) (1.35,7.287) (1.374,7.281) (1.399,7.276) (1.424,7.27) (1.448,7.264) (1.473,7.258) (1.499,7.252) (1.524,7.245) (1.55,7.239) (1.576,7.232) (1.602,7.226) (1.628,7.219) (1.655,7.212) (1.682,7.206) (1.709,7.199) (1.736,7.192) (1.764,7.184) (1.791,7.177) (1.819,7.17) (1.847,7.162) (1.875,7.154) (1.904,7.146) (1.933,7.139) (1.962,7.13) (1.991,7.122) (2.02,7.114) (2.05,7.106) (2.08,7.097) (2.11,7.088) (2.14,7.079) (2.17,7.07) (2.201,7.061) (2.232,7.052) (2.263,7.042) (2.294,7.033) (2.326,7.023) (2.358,7.013) (2.39,7.003) (2.422,6.993) (2.454,6.982) (2.487,6.972) (2.52,6.961) (2.553,6.95) (2.586,6.939) (2.619,6.927) (2.653,6.916) (2.687,6.904) (2.721,6.892) (2.756,6.88) (2.79,6.868) (2.825,6.855) (2.86,6.842) (2.895,6.83) (2.93,6.816) (2.966,6.803) (3.002,6.789) (3.038,6.775) (3.074,6.761) (3.111,6.747) (3.147,6.732) (3.184,6.718) (3.222,6.703) (3.259,6.687) (3.296,6.672) (3.334,6.656) (3.372,6.64) (3.41,6.623) (3.449,6.606) (3.488,6.589) (3.526,6.572) (3.565,6.554) (3.605,6.537) (3.644,6.518) (3.684,6.5) (3.724,6.481) (3.764,6.462) (3.804,6.442) (3.845,6.422) (3.886,6.402) (3.927,6.381) (3.968,6.36) (4.009,6.339) (4.051,6.317) (4.093,6.295) (4.135,6.272) (4.177,6.249) (4.22,6.226) (4.263,6.202) (4.306,6.178) (4.349,6.153) (4.392,6.128) (4.436,6.102) (4.48,6.076) (4.524,6.05) (4.568,6.022) (4.612,5.995) (4.657,5.967) (4.702,5.938) (4.747,5.909) (4.792,5.879) (4.838,5.849) (4.883,5.818) (4.929,5.787) (4.976,5.755) (5.022,5.722) (5.069,5.689) (5.115,5.655) (5.162,5.62) (5.21,5.585) (5.257,5.549) (5.305,5.513) (5.353,5.476) (5.401,5.438) (5.449,5.399) (5.498,5.359) (5.547,5.319) (5.596,5.278) (5.645,5.236) (5.694,5.194) (5.744,5.15) (5.794,5.106) (5.844,5.061) (5.894,5.014) (5.944,4.967) (5.995,4.919) (6.046,4.87) (6.097,4.82) (6.148,4.77) (6.2,4.718)};
\end{scope}
\node[above,font=\small\sffamily] at (3.1,7.9) {下側（壁〜極小）の電荷密度};
\draw[black!10,line width=0.4pt] (8.493,4.6)--(8.493,7.9);
\draw[black!10,line width=0.4pt] (9.979,4.6)--(9.979,7.9);
\draw[black!10,line width=0.4pt] (11.465,4.6)--(11.465,7.9);
\draw[black!10,line width=0.4pt] (12.951,4.6)--(12.951,7.9);
\draw[black!10,line width=0.4pt] (8,4.6)--(14.2,4.6);
\draw[black!10,line width=0.4pt] (8,5.7)--(14.2,5.7);
\draw[black!10,line width=0.4pt] (8,6.8)--(14.2,6.8);
\draw[black!10,line width=0.4pt] (8,7.9)--(14.2,7.9);
\draw[line width=0.6pt] (8,4.6) rectangle (14.2,7.9);
\draw[line width=0.6pt] (8.493,4.6)--++(0,0.1);
\node[below,font=\scriptsize\sffamily] at (8.493,4.6) {$0$};
\draw[line width=0.6pt] (9.979,4.6)--++(0,0.1);
\node[below,font=\scriptsize\sffamily] at (9.979,4.6) {$1$};
\draw[line width=0.6pt] (11.465,4.6)--++(0,0.1);
\node[below,font=\scriptsize\sffamily] at (11.465,4.6) {$2$};
\draw[line width=0.6pt] (12.951,4.6)--++(0,0.1);
\node[below,font=\scriptsize\sffamily] at (12.951,4.6) {$3$};
\draw[line width=0.6pt] (8,4.6)--++(0.1,0);
\node[left,font=\scriptsize\sffamily] at (8,4.6) {$0$};
\draw[line width=0.6pt] (8,5.7)--++(0.1,0);
\node[left,font=\scriptsize\sffamily] at (8,5.7) {$1$};
\draw[line width=0.6pt] (8,6.8)--++(0.1,0);
\node[left,font=\scriptsize\sffamily] at (8,6.8) {$2$};
\draw[line width=0.6pt] (8,7.9)--++(0.1,0);
\node[left,font=\scriptsize\sffamily] at (8,7.9) {$3$};
\node[font=\small\sffamily] at (11.1,4) {$\phi$ [V]};
\node[rotate=90,font=\small\sffamily] at (7.2,6.25) {$E^2$ [V$^2$/m$^2$]};
\begin{scope}
\clip (8,4.6) rectangle (14.2,7.9);
\draw[black!45,dashed,line width=0.5pt] (8,4.6)--(14.2,4.6);
\draw[BeachRed,line width=1.2pt] plot coordinates {(8,4.6) (8.015,4.6) (8.031,4.6) (8.047,4.601) (8.067,4.601) (8.084,4.601) (8.103,4.602) (8.124,4.602) (8.14,4.603) (8.155,4.603) (8.172,4.604) (8.19,4.604) (8.208,4.605) (8.228,4.606) (8.248,4.606) (8.269,4.607) (8.291,4.608) (8.314,4.609) (8.338,4.61) (8.362,4.611) (8.387,4.612) (8.414,4.614) (8.441,4.615) (8.469,4.617) (8.498,4.618) (8.527,4.62) (8.543,4.621) (8.558,4.622) (8.574,4.623) (8.589,4.624) (8.605,4.625) (8.622,4.626) (8.638,4.627) (8.655,4.629) (8.672,4.63) (8.689,4.631) (8.706,4.633) (8.724,4.634) (8.742,4.635) (8.76,4.637) (8.778,4.638) (8.796,4.64) (8.815,4.641) (8.834,4.643) (8.853,4.645) (8.872,4.647) (8.891,4.648) (8.911,4.65) (8.931,4.652) (8.951,4.654) (8.971,4.656) (8.992,4.658) (9.013,4.661) (9.034,4.663) (9.055,4.665) (9.076,4.668) (9.098,4.67) (9.12,4.672) (9.142,4.675) (9.164,4.678) (9.187,4.68) (9.209,4.683) (9.232,4.686) (9.256,4.689) (9.279,4.692) (9.302,4.695) (9.326,4.699) (9.35,4.702) (9.374,4.705) (9.399,4.709) (9.424,4.712) (9.448,4.716) (9.473,4.72) (9.499,4.724) (9.524,4.728) (9.55,4.732) (9.576,4.736) (9.602,4.741) (9.628,4.745) (9.655,4.75) (9.682,4.754) (9.709,4.759) (9.736,4.764) (9.764,4.769) (9.791,4.774) (9.819,4.779) (9.847,4.785) (9.875,4.79) (9.904,4.796) (9.933,4.802) (9.962,4.808) (9.991,4.814) (10.02,4.821) (10.05,4.827) (10.08,4.834) (10.11,4.841) (10.14,4.847) (10.17,4.855) (10.201,4.862) (10.232,4.869) (10.263,4.877) (10.294,4.885) (10.326,4.893) (10.358,4.901) (10.39,4.91) (10.422,4.918) (10.454,4.927) (10.487,4.936) (10.52,4.946) (10.553,4.955) (10.586,4.965) (10.619,4.975) (10.653,4.985) (10.687,4.995) (10.721,5.006) (10.756,5.017) (10.79,5.028) (10.825,5.04) (10.86,5.052) (10.895,5.064) (10.93,5.076) (10.966,5.088) (11.002,5.101) (11.038,5.114) (11.074,5.128) (11.111,5.142) (11.147,5.156) (11.184,5.17) (11.222,5.185) (11.259,5.2) (11.296,5.215) (11.334,5.231) (11.372,5.247) (11.41,5.264) (11.449,5.281) (11.488,5.298) (11.526,5.316) (11.565,5.334) (11.605,5.352) (11.644,5.371) (11.684,5.391) (11.724,5.41) (11.764,5.431) (11.804,5.451) (11.845,5.472) (11.886,5.494) (11.927,5.516) (11.968,5.539) (12.009,5.562) (12.051,5.586) (12.093,5.61) (12.135,5.635) (12.177,5.66) (12.22,5.686) (12.263,5.712) (12.306,5.739) (12.349,5.767) (12.392,5.795) (12.436,5.824) (12.48,5.854) (12.524,5.884) (12.568,5.915) (12.612,5.946) (12.657,5.979) (12.702,6.012) (12.747,6.045) (12.792,6.08) (12.838,6.115) (12.883,6.151) (12.929,6.188) (12.976,6.226) (13.022,6.264) (13.069,6.304) (13.115,6.344) (13.162,6.385) (13.21,6.427) (13.257,6.47) (13.305,6.515) (13.353,6.56) (13.401,6.606) (13.449,6.653) (13.498,6.701) (13.547,6.75) (13.596,6.8) (13.645,6.852) (13.694,6.904) (13.744,6.958) (13.794,7.013) (13.844,7.069) (13.894,7.127) (13.944,7.186) (13.995,7.246) (14.046,7.307) (14.097,7.37) (14.148,7.435) (14.2,7.5)};
\end{scope}
\node[above,font=\small\sffamily] at (11.1,7.9) {下側の $E^2(\phi)$};
\node[anchor=east,font=\scriptsize\sffamily] at (14.2,7.735) {$E_0^2$};
\fill[BeachOrange!35] (1.524,2.75) -- plot coordinates {(1.524,2.759) (1.602,2.795) (1.682,2.829) (1.764,2.861) (1.847,2.891) (1.933,2.919) (2.02,2.945) (2.11,2.969) (2.201,2.991) (2.294,3.012) (2.39,3.03) (2.487,3.047) (2.586,3.062) (2.687,3.075) (2.79,3.086) (2.895,3.095) (3.002,3.103) (3.111,3.108) (3.222,3.112) (3.334,3.115) (3.449,3.115) (3.565,3.114) (3.684,3.11) (3.804,3.106) (3.927,3.099) (4.051,3.091) (4.177,3.08) (4.306,3.069) (4.436,3.055) (4.568,3.04) (4.702,3.023) (4.838,3.004) (4.976,2.984) (5.115,2.962) (5.257,2.938) (5.401,2.913) (5.547,2.886) (5.694,2.857) (5.844,2.827) (5.995,2.795) (6.148,2.762)} -- (6.148,2.75) -- cycle;
\fill[BeachBlue!20] (0,2.75) -- plot coordinates {(0,-0.389) (0.001,-0.264) (0.004,-0.142) (0.009,-0.022) (0.016,0.095) (0.024,0.21) (0.035,0.322) (0.047,0.432) (0.062,0.539) (0.078,0.644) (0.097,0.746) (0.117,0.846) (0.139,0.944) (0.164,1.039) (0.19,1.132) (0.218,1.222) (0.248,1.31) (0.28,1.396) (0.314,1.48) (0.35,1.561) (0.387,1.64) (0.427,1.717) (0.469,1.792) (0.512,1.864) (0.558,1.934) (0.605,2.002) (0.655,2.068) (0.706,2.132) (0.76,2.193) (0.815,2.253) (0.872,2.31) (0.931,2.366) (0.992,2.419) (1.055,2.47) (1.12,2.519) (1.187,2.566) (1.255,2.611) (1.326,2.654) (1.399,2.695) (1.473,2.734)} -- (6.2,2.75) -- cycle;
\draw[black!10,line width=0.4pt] (0.592,0)--(0.592,3.3);
\draw[black!10,line width=0.4pt] (2.461,0)--(2.461,3.3);
\draw[black!10,line width=0.4pt] (4.331,0)--(4.331,3.3);
\draw[black!10,line width=0.4pt] (6.2,0)--(6.2,3.3);
\draw[black!10,line width=0.4pt] (0,0.55)--(6.2,0.55);
\draw[black!10,line width=0.4pt] (0,1.65)--(6.2,1.65);
\draw[black!10,line width=0.4pt] (0,2.75)--(6.2,2.75);
\draw[line width=0.6pt] (0,0) rectangle (6.2,3.3);
\draw[line width=0.6pt] (0.592,0)--++(0,0.1);
\node[below,font=\scriptsize\sffamily] at (0.592,0) {$-0.3$};
\draw[line width=0.6pt] (2.461,0)--++(0,0.1);
\node[below,font=\scriptsize\sffamily] at (2.461,0) {$-0.2$};
\draw[line width=0.6pt] (4.331,0)--++(0,0.1);
\node[below,font=\scriptsize\sffamily] at (4.331,0) {$-0.1$};
\draw[line width=0.6pt] (6.2,0)--++(0,0.1);
\node[below,font=\scriptsize\sffamily] at (6.2,0) {$0$};
\draw[line width=0.6pt] (0,0.55)--++(0.1,0);
\node[left,font=\scriptsize\sffamily] at (0,0.55) {$-0.2$};
\draw[line width=0.6pt] (0,1.65)--++(0.1,0);
\node[left,font=\scriptsize\sffamily] at (0,1.65) {$-0.1$};
\draw[line width=0.6pt] (0,2.75)--++(0.1,0);
\node[left,font=\scriptsize\sffamily] at (0,2.75) {$0$};
\node[font=\small\sffamily] at (3.1,-0.6) {$\phi$ [V]};
\node[rotate=90,font=\small\sffamily] at (-1,1.65) {$\rho/e$ [cm$^{-3}$]};
\begin{scope}
\clip (0,0) rectangle (6.2,3.3);
\draw[black!45,dashed,line width=0.5pt] (0,2.75)--(6.2,2.75);
\draw[BeachBlue,line width=1.2pt] plot coordinates {(0,-0.389) (0,-0.347) (0,-0.305) (0.001,-0.264) (0.002,-0.223) (0.003,-0.182) (0.004,-0.142) (0.005,-0.102) (0.007,-0.062) (0.009,-0.022) (0.011,0.017) (0.013,0.056) (0.016,0.095) (0.018,0.134) (0.021,0.172) (0.024,0.21) (0.028,0.248) (0.031,0.285) (0.035,0.322) (0.039,0.359) (0.043,0.396) (0.047,0.432) (0.052,0.468) (0.057,0.504) (0.062,0.539) (0.067,0.574) (0.073,0.609) (0.078,0.644) (0.084,0.678) (0.091,0.712) (0.097,0.746) (0.103,0.78) (0.11,0.813) (0.117,0.846) (0.124,0.879) (0.132,0.911) (0.139,0.944) (0.147,0.976) (0.155,1.007) (0.164,1.039) (0.172,1.07) (0.181,1.101) (0.19,1.132) (0.199,1.162) (0.208,1.192) (0.218,1.222) (0.228,1.252) (0.238,1.281) (0.248,1.31) (0.258,1.339) (0.269,1.368) (0.28,1.396) (0.291,1.424) (0.302,1.452) (0.314,1.48) (0.326,1.507) (0.338,1.534) (0.35,1.561) (0.362,1.588) (0.375,1.614) (0.387,1.64) (0.401,1.666) (0.414,1.692) (0.427,1.717) (0.441,1.742) (0.455,1.767) (0.469,1.792) (0.483,1.816) (0.498,1.84) (0.512,1.864) (0.527,1.888) (0.543,1.911) (0.558,1.934) (0.574,1.957) (0.589,1.98) (0.605,2.002) (0.622,2.024) (0.638,2.046) (0.655,2.068) (0.672,2.09) (0.689,2.111) (0.706,2.132) (0.724,2.153) (0.742,2.173) (0.76,2.193) (0.778,2.214) (0.796,2.233) (0.815,2.253) (0.834,2.272) (0.853,2.291) (0.872,2.31) (0.891,2.329) (0.911,2.347) (0.931,2.366) (0.951,2.384) (0.971,2.401) (0.992,2.419) (1.013,2.436) (1.034,2.453) (1.055,2.47) (1.076,2.487) (1.098,2.503) (1.12,2.519) (1.142,2.535) (1.164,2.551) (1.187,2.566) (1.209,2.581) (1.232,2.596) (1.255,2.611) (1.279,2.626) (1.302,2.64) (1.326,2.654) (1.35,2.668) (1.374,2.682) (1.399,2.695) (1.424,2.708) (1.448,2.721) (1.473,2.734) (1.499,2.747) (1.524,2.759) (1.55,2.771) (1.576,2.783) (1.602,2.795) (1.628,2.806) (1.655,2.818) (1.682,2.829) (1.709,2.84) (1.736,2.85) (1.764,2.861) (1.791,2.871) (1.819,2.881) (1.847,2.891) (1.875,2.9) (1.904,2.909) (1.933,2.919) (1.962,2.927) (1.991,2.936) (2.02,2.945) (2.05,2.953) (2.08,2.961) (2.11,2.969) (2.14,2.976) (2.17,2.984) (2.201,2.991) (2.232,2.998) (2.263,3.005) (2.294,3.012) (2.326,3.018) (2.358,3.024) (2.39,3.03) (2.422,3.036) (2.454,3.041) (2.487,3.047) (2.52,3.052) (2.553,3.057) (2.586,3.062) (2.62,3.066) (2.653,3.071) (2.687,3.075) (2.721,3.079) (2.756,3.082) (2.79,3.086) (2.825,3.089) (2.86,3.092) (2.895,3.095) (2.93,3.098) (2.966,3.1) (3.002,3.103) (3.038,3.105) (3.074,3.107) (3.111,3.108) (3.147,3.11) (3.184,3.111) (3.222,3.112) (3.259,3.113) (3.296,3.114) (3.334,3.115) (3.372,3.115) (3.41,3.115) (3.449,3.115) (3.488,3.115) (3.526,3.114) (3.565,3.114) (3.605,3.113) (3.644,3.112) (3.684,3.11) (3.724,3.109) (3.764,3.107) (3.804,3.106) (3.845,3.104) (3.886,3.101) (3.927,3.099) (3.968,3.096) (4.009,3.093) (4.051,3.091) (4.093,3.087) (4.135,3.084) (4.177,3.08) (4.22,3.077) (4.263,3.073) (4.306,3.069) (4.349,3.064) (4.392,3.06) (4.436,3.055) (4.48,3.05) (4.524,3.045) (4.568,3.04) (4.612,3.034) (4.657,3.029) (4.702,3.023) (4.747,3.017) (4.792,3.011) (4.838,3.004) (4.883,2.998) (4.929,2.991) (4.976,2.984) (5.022,2.977) (5.069,2.969) (5.115,2.962) (5.162,2.954) (5.21,2.946) (5.257,2.938) (5.305,2.93) (5.353,2.922) (5.401,2.913) (5.449,2.904) (5.498,2.895) (5.547,2.886) (5.596,2.877) (5.645,2.867) (5.694,2.857) (5.744,2.847) (5.794,2.837) (5.844,2.827) (5.894,2.817) (5.944,2.806) (5.995,2.795) (6.046,2.784) (6.097,2.773) (6.148,2.762) (6.2,2.75)};
\end{scope}
\node[above,font=\small\sffamily] at (3.1,3.3) {上側（極小〜上流）の電荷密度};
\node[font=\scriptsize\sffamily] at (1.153,2.255) {負};
\node[font=\scriptsize\sffamily] at (4.144,3.025) {正};
\draw[black!10,line width=0.4pt] (8.592,0)--(8.592,3.3);
\draw[black!10,line width=0.4pt] (10.461,0)--(10.461,3.3);
\draw[black!10,line width=0.4pt] (12.331,0)--(12.331,3.3);
\draw[black!10,line width=0.4pt] (14.2,0)--(14.2,3.3);
\draw[black!10,line width=0.4pt] (8,0)--(14.2,0);
\draw[black!10,line width=0.4pt] (8,1.32)--(14.2,1.32);
\draw[black!10,line width=0.4pt] (8,2.64)--(14.2,2.64);
\draw[line width=0.6pt] (8,0) rectangle (14.2,3.3);
\draw[line width=0.6pt] (8.592,0)--++(0,0.1);
\node[below,font=\scriptsize\sffamily] at (8.592,0) {$-0.3$};
\draw[line width=0.6pt] (10.461,0)--++(0,0.1);
\node[below,font=\scriptsize\sffamily] at (10.461,0) {$-0.2$};
\draw[line width=0.6pt] (12.331,0)--++(0,0.1);
\node[below,font=\scriptsize\sffamily] at (12.331,0) {$-0.1$};
\draw[line width=0.6pt] (14.2,0)--++(0,0.1);
\node[below,font=\scriptsize\sffamily] at (14.2,0) {$0$};
\draw[line width=0.6pt] (8,0)--++(0.1,0);
\node[left,font=\scriptsize\sffamily] at (8,0) {$0$};
\draw[line width=0.6pt] (8,1.32)--++(0.1,0);
\node[left,font=\scriptsize\sffamily] at (8,1.32) {$1$};
\draw[line width=0.6pt] (8,2.64)--++(0.1,0);
\node[left,font=\scriptsize\sffamily] at (8,2.64) {$2$};
\node[font=\small\sffamily] at (11.1,-0.6) {$\phi$ [V]};
\node[rotate=90,font=\small\sffamily] at (7.2,1.65) {$E^2$ [$10^{-4}$ V$^2$/m$^2$]};
\begin{scope}
\clip (8,0) rectangle (14.2,3.3);
\draw[black!45,dashed,line width=0.5pt] (8,0)--(14.2,0);
\draw[BeachBlue,line width=1.2pt] plot coordinates {(8,0) (8.003,0.019) (8.005,0.036) (8.009,0.059) (8.013,0.086) (8.016,0.101) (8.018,0.118) (8.021,0.135) (8.024,0.154) (8.028,0.173) (8.031,0.194) (8.035,0.215) (8.039,0.238) (8.043,0.261) (8.047,0.285) (8.052,0.309) (8.057,0.335) (8.062,0.361) (8.067,0.388) (8.073,0.415) (8.078,0.444) (8.084,0.472) (8.091,0.502) (8.097,0.531) (8.103,0.562) (8.11,0.593) (8.117,0.624) (8.124,0.655) (8.132,0.687) (8.139,0.72) (8.147,0.752) (8.155,0.785) (8.164,0.819) (8.172,0.852) (8.181,0.886) (8.19,0.92) (8.199,0.954) (8.208,0.988) (8.218,1.022) (8.228,1.057) (8.238,1.091) (8.248,1.126) (8.258,1.16) (8.269,1.195) (8.28,1.229) (8.291,1.264) (8.302,1.298) (8.314,1.333) (8.326,1.367) (8.338,1.401) (8.35,1.435) (8.362,1.469) (8.375,1.503) (8.387,1.536) (8.401,1.569) (8.414,1.602) (8.427,1.635) (8.441,1.667) (8.455,1.699) (8.469,1.731) (8.483,1.763) (8.498,1.794) (8.512,1.824) (8.527,1.855) (8.543,1.885) (8.558,1.914) (8.574,1.944) (8.589,1.972) (8.605,2.001) (8.622,2.028) (8.638,2.056) (8.655,2.083) (8.672,2.109) (8.689,2.135) (8.706,2.16) (8.724,2.185) (8.742,2.209) (8.76,2.233) (8.778,2.256) (8.796,2.278) (8.815,2.3) (8.834,2.322) (8.853,2.342) (8.872,2.362) (8.891,2.382) (8.911,2.401) (8.931,2.419) (8.951,2.436) (8.971,2.453) (8.992,2.47) (9.013,2.485) (9.034,2.5) (9.055,2.514) (9.076,2.528) (9.098,2.541) (9.12,2.553) (9.142,2.564) (9.164,2.575) (9.187,2.585) (9.209,2.594) (9.232,2.603) (9.255,2.611) (9.279,2.618) (9.302,2.624) (9.326,2.63) (9.35,2.635) (9.374,2.639) (9.399,2.642) (9.424,2.645) (9.448,2.647) (9.473,2.649) (9.499,2.649) (9.524,2.649) (9.55,2.648) (9.576,2.646) (9.602,2.644) (9.628,2.641) (9.655,2.637) (9.682,2.632) (9.709,2.627) (9.736,2.621) (9.764,2.614) (9.791,2.607) (9.819,2.599) (9.847,2.59) (9.875,2.58) (9.904,2.57) (9.933,2.559) (9.962,2.548) (9.991,2.535) (10.02,2.522) (10.05,2.509) (10.08,2.494) (10.11,2.479) (10.14,2.464) (10.17,2.447) (10.201,2.43) (10.232,2.413) (10.263,2.395) (10.294,2.376) (10.326,2.357) (10.358,2.337) (10.39,2.316) (10.422,2.295) (10.454,2.273) (10.487,2.251) (10.52,2.228) (10.553,2.205) (10.586,2.181) (10.62,2.156) (10.653,2.131) (10.687,2.106) (10.721,2.08) (10.756,2.054) (10.79,2.027) (10.825,2) (10.86,1.972) (10.895,1.944) (10.93,1.915) (10.966,1.886) (11.002,1.857) (11.038,1.827) (11.074,1.797) (11.111,1.767) (11.147,1.737) (11.184,1.706) (11.222,1.674) (11.259,1.643) (11.296,1.611) (11.334,1.579) (11.372,1.547) (11.41,1.515) (11.449,1.482) (11.488,1.449) (11.526,1.416) (11.565,1.383) (11.605,1.35) (11.644,1.317) (11.684,1.284) (11.724,1.25) (11.764,1.217) (11.804,1.183) (11.845,1.15) (11.886,1.117) (11.927,1.083) (11.968,1.05) (12.009,1.017) (12.051,0.984) (12.093,0.951) (12.135,0.918) (12.177,0.885) (12.22,0.853) (12.263,0.821) (12.306,0.789) (12.349,0.757) (12.392,0.725) (12.436,0.694) (12.48,0.663) (12.524,0.633) (12.568,0.603) (12.612,0.573) (12.657,0.544) (12.702,0.515) (12.747,0.487) (12.792,0.459) (12.838,0.432) (12.883,0.405) (12.929,0.379) (12.976,0.354) (13.022,0.329) (13.069,0.305) (13.115,0.281) (13.162,0.259) (13.21,0.237) (13.257,0.215) (13.305,0.195) (13.353,0.175) (13.401,0.157) (13.449,0.139) (13.498,0.122) (13.547,0.106) (13.596,0.091) (13.645,0.077) (13.694,0.064) (13.744,0.053) (13.794,0.042) (13.844,0.032) (13.894,0.024) (13.944,0.017) (13.995,0.011) (14.046,0.006) (14.097,0.003) (14.148,0.001) (14.2,0)};
\end{scope}
\node[above,font=\small\sffamily] at (11.1,3.3) {上側の $E^2(\phi)$};
\node[anchor=east,font=\scriptsize\sffamily] at (13.639,0.462) {$E^2(0)=0$};
\end{tikzpicture}
\caption[Type A の電荷密度と第一積分]{Type A（$\alpha=60^\circ$）の電荷密度 $\rho(\phi)$ と第一積分 $E^2(\phi)$。横軸は位置ではなく電位で、上段は壁から極小まで、下段は極小から上流までの区間である。二つの区間では粒子集団が異なるため、同じ電位でも $\rho$ が異なる。$E^2$ は極小 $\phi_m$ で 0 から始まり、上側では上流 $\phi=0$ で再び 0 に戻る。これは上側の電荷密度の負の面積と正の面積が打ち消すこと（式\eqref{eq:sm_connection}）に等しい。$E^2$ は数値求積で独立に計算した。}
\label{fig:sm_first_integral}
\end{figure}


\subsection{三つの型}
電位の形により解は三つに分かれる\citep{Zhao2020}。Type A は $\phi_m<\min(\phi_0,0)$ の内部極小を持つ非単調な構造、
Type B は正の壁から 0 へ単調に下がる構造、Type C は負の壁から 0 へ単調に上がる構造である
（図\ref{fig:sheath_types}）。Type A の壁電位は正とは限らず、負の壁電位を持つ Type A もある。
光電子がなければ壁を正にする機構がなく、Type C だけが成り立つ。
同じ条件で複数の型が候補になり得ることは Sana と Mishra が示しており\citep{SanaMishra2023}、
本モデルは型の間の安定性による選択をしない。

\begin{figure}[htbp]
\centering
\begin{tikzpicture}[beach diagram]
  \foreach \shift/\name in {0/A,5/B,10/C} {
    \begin{scope}[xshift=\shift cm]
      \node[beach heading] at (1.9,5.3) {Type \name};
      \fill[black!8] (-0.15,0.25) rectangle (0,4.8);
      \draw[beach axis] (0,0.25)--(4.2,0.25) node[right] {$z$};
      \draw[beach axis] (0,0.25)--(0,4.9) node[above] {$\phi$};
      \draw[black!45,dashed] (0,2.5)--(4,2.5);
      \node[left] at (0,2.5) {$0$};
      \node[below] at (0,0.25) {壁};
      \node[below] at (3.3,0.25) {上流へ};
    \end{scope}
  }
  \draw[BeachBlue,thick] (0,3.75) .. controls (0.45,1.9) and (0.8,0.85) .. (1.4,0.85)
    .. controls (2.05,0.85) and (2.25,2.5) .. (3.8,2.5);
  \draw[BeachBlue,dotted] (0,0.85)--(1.4,0.85);
  \node[left] at (0,0.85) {$\phi_m$};
  \node[left] at (0,3.75) {$\phi_0$};
  \node[beach dot] at (1.4,0.85) {};
  \draw[BeachTeal,thick,dashed] (5,4.1) .. controls (5.8,3.15) and (6.4,2.5) .. (8.8,2.5);
  \node[left] at (5,4.1) {$\phi_0$};
  \draw[BeachOrange,thick,dash dot] (10,0.85) .. controls (10.8,1.8) and (11.5,2.5) .. (13.8,2.5);
  \node[left] at (10,0.85) {$\phi_0$};
  \node[beach note] at (1.9,-0.65) {非単調、負の極小を持つ};
  \node[beach note] at (6.9,-0.65) {正の壁から単調に低下};
  \node[beach note] at (11.9,-0.65) {負の壁から単調に上昇};
\end{tikzpicture}
\caption[外部シースの Type A、B、C]{外部シースの電位形状による Type A、B、C の分類。曲線は形の違いだけを示す模式図で、
壁電位、距離、勾配は数値解を表していない。Type A は正の壁電位を持つ例を描いた。
電子 drift による上流の帯はこの分類図に含めず、図\ref{fig:sm_upstream_band}で示す。}
\label{fig:sheath_types}
\end{figure}


\subsection{電位分布の再構成}
空間分布は第一積分から
\begin{equation}
 z(\phi)=\int_\phi^{\phi_0}\frac{\dd\psi}{|E(\psi)|}
 \label{eq:sm_z_of_phi}
\end{equation}
で得る。Type A では下側の積分で極小の位置 $z_m$ を求め、上側を $z_m$ から積分する。
極小の近くでは $E^2\approx-2\rho(\phi_m)(\phi-\phi_m)/\eps$ なので $\phi-\phi_m\propto(z-z_m)^2$ となり、被積分関数の特異点は積分できる。
上流の近くでは $\rho\approx\rho'(0)\phi$ から $E^2\approx-\rho'(0)\phi^2/\eps$ となり、$\phi$ は長さ $L=\sqrt{\eps/(-\rho'(0))}$ で指数関数的に 0 に近づく。
上流までの距離は無限なので、分布は $|\phi|$ が小さな閾値に達した点で打ち切る。

\section{閉じる条件}
\label{sec:sm_closure}
未知数は壁電位 $\phi_0$、Type A の極小 $\phi_m$、太陽風電子の規格化 $N_e$ である。零電流の closure は次の条件で根を決める。
\begin{enumerate}
\item \textbf{上流の準中性：}
 \begin{equation}
  n_{i\infty}=n_e(0)+n_{\mathrm{pe}}(0),\qquad
  n_{\mathrm{pe}}(0)=\frac{N_{\mathrm{pe}}}2e^{-\phi_0/T_{\mathrm{pe}}}\operatorname{erfc}\sqrt{-\phi_{\min}/T_{\mathrm{pe}}}
  \label{eq:sm_neutrality}
 \end{equation}
 （bin の光電子源では式\eqref{eq:sm_pe_bins}を $\phi=0$ で使う）。
\item \textbf{壁への零電流：}電子の吸収 $J_e=-e\Gamma_e$、イオンの吸収 $J_i=e\Gamma_i$、光電子の正味の放出 $J_{\mathrm{escape}}=e\Gamma_{\mathrm{esc}}$ に対して
 \begin{equation}
  J_e+J_i+J_{\mathrm{escape}}=0
  \label{eq:zero_current}
 \end{equation}
 である。壁へ戻る光電子は放出と相殺するので、正味には脱出分だけが残る。
\item \textbf{Type A の接続：}式\eqref{eq:sm_connection}。
\end{enumerate}
Type B、C の未知数は 2 つ、Type A は 3 つで、方程式の数と一致する。$N_e$ は準中性と電流の両方に線形に入る。
これらは電位空間の代数方程式と積分条件であり、$z$ 方向の積分を含まない。
太陽高度 $\alpha$ で照射とプラズマ流を与える場合は、Zhao らに従って
\begin{equation}
 N_{\mathrm{pe}}=n_{\mathrm{pe,ref}}\sin\alpha,\qquad u_i=v_{\mathrm{sw}}\sin\alpha
 \label{eq:sm_illumination}
\end{equation}
とする。$n_{\mathrm{pe,ref}}$ は垂直照射での光電子の規格化、$v_{\mathrm{sw}}$ は太陽風の速さである。
電子の drift は、太陽風の法線速度 $v_d=v_{\mathrm{sw}}\sin\alpha$ と $v_d=0$ を物理条件として明示的に選ぶ（\ref{sec:sm_band}節）。

零電流のほかに二つの閉じ方がある（表\ref{tab:sm_closures}）。境界電場を指定する closure は、式\eqref{eq:zero_current}の代わりに
$E^2(\phi_0)=E_H^2$ と電場の符号を課し、電流を結果として返す。$N_e$ は準中性で決まるので、$E_H$ を変えると $N_e$ も変わる。
$E_H=0$ は平坦な解だけでなく、非平坦な Type B、C の根も含む。静的な評価は、与えた電位で準中性から $N_e$ を求め、
符号付きの $E^2(\phi_0)$、Type A の接続の残差、各電流を返す。外側の時間発展や結合方程式は呼び出し側が解き、
その結果を物理解の判定にかける。

\begin{table}[htbp]
\centering\small
\caption{外部シースの閉じ方と未知数。いずれも同じ密度、流束、第一積分、物理解の判定を使う。}
\label{tab:sm_closures}
\begin{tabularx}{\textwidth}{p{27mm}YY}
\toprule
閉じ方 & 固定する量 & 決める量 \\
\midrule
零電流 & 上流のイオン、電子温度と drift、光電子源 & $\phi_0$、$\phi_m$、$N_e$、各電流 \\
境界電場の指定 & 上記と $E_H$ & $\phi_0$、$\phi_m$、$N_e$、正味の電流 \\
静的な評価 & 上記と $\phi_0$（と $\phi_m$） & $N_e$、符号付きの $E^2(\phi_0)$、接続の残差、各電流 \\
\bottomrule
\end{tabularx}
\end{table}

\section{物理解の条件}
\label{sec:sm_admissibility}
代数方程式の根であることと、上流から壁まで実数の電位分布がつながることは別の条件である。根は次をすべて満たすときだけ受理する。
\begin{enumerate}
\item \textbf{イオンの到達：}経路上の最大電位 $\max(\phi_0,0)$ が式\eqref{eq:sm_ion_crit}の $\phi_{\mathrm{crit}}$ 以下である。
\item \textbf{実数の電場：}全区間で $E^2(\phi)\ge0$ であり、Type A では極小で $\rho<0$、接続条件\eqref{eq:sm_connection}が満たされる。
\item \textbf{上流での漸近：}上流の近くでも $E^2\ge0$ が保たれる。
\end{enumerate}

\subsection{上流の漸近条件}
上流の近くで $\rho$ が $\phi$ について滑らかなら、$E^2\approx-\rho'(0)\phi^2/\eps\ge0$ から $\dd\rho/\dd\phi\le0$ が必要である。
冷たいイオンと Boltzmann 電子では、これは Bohm の条件 $u_i\ge\sqrt{eT_e/m_i}$ に一致する\citep{Riemann1991}。
本モデルでは入射速度を Bohm 速度へ補正せず、指定した入射状態でこの条件が成り立つかを判定する。

上流の分布が速度 0 で切れている場合、$\rho$ は滑らかでない。Type B（$\phi>0$）では通過電子の下限速度 $\sqrt\psi$ と
捕捉光電子のために $\rho/e=C\sqrt\phi+O(\phi)$ となり、Maxwell の光電子源では
\begin{equation}
 C=\frac{N_ee^{-u^2}}{\sqrt{\pi T_e}}-\frac{N_{\mathrm{pe}}e^{-\phi_0/T_{\mathrm{pe}}}}{\sqrt{\pi T_{\mathrm{pe}}}}
 \label{eq:sm_b_coefficient}
\end{equation}
である。式\eqref{eq:sagdeev}から $E^2\approx-4eC\phi^{3/2}/(3\eps)$ なので、$C\le0$ が必要である。
bin の光電子源では第 2 項を $2(2g_-{-}g_+)/v_*$ に置き換える。$g_\mp$ は $K=\phi_0$ の左右の bin の値で、
戻る軌道と脱出する軌道が別の側の値を使う。この条件は有限個の点での検査では見逃し得るので、係数で判定する。

\subsection{電子 drift と上流の帯}
\label{sec:sm_band}
上流で壁へ向かう電子の分布を、規格化した速さ $a$ の関数 $g(a)=N_ee^{-(a-u)^2}/\sqrt\pi$ と書く。
Type A の上側と Type C で上流から深さ $h=-\phi/T_e$ の点では、遅い電子は両方向に存在するので、密度の変化は
\begin{equation}
 n_e(h)-n_e(0)=2\int_0^{\infty}\left[g\!\left(\sqrt{w^2+h}\right)-g(w)\right]\dd w+O(h)
 \approx g'(0)\,\frac h2\log\frac1h+O(h)
 \label{eq:hlogh}
\end{equation}
となる。drift のない Maxwell 分布は $g'(0)=0$ でこの項を持たないが、drift のある分布は $g'(0)=2u\,g(0)>0$ なので、
上流に近づくほど電子の過剰が $h\log(1/h)$ で現れる。イオンと光電子の密度は $\phi$ について滑らかなのでこれを打ち消せず、
上流電位に接する帯 $\phi_b<\phi<0$ で $E^2<0$ になる。$E^2<0$ の区間には実数の $z(\phi)$ がないので、
厳密な半無限の上流条件のもとでは、drift のある Type A、C の根は物理解にならない。
この上流分布は $v_z=0$ に尖った極小を持ち、Penrose の条件\citep{Penrose1960}では drift の大きさによらず運動論的に不安定である。
実際には緩和や散乱で滑らかになると考えられ、遅い入射電子の drift だけを除いて $g'(0)=0$ にすると帯は消える。

帯の内側で $E^2$ が負になる量から、場の不足
\begin{equation}
 E_{\mathrm{def}}=\sqrt{-\min_{\phi_b<\phi<0}E^2(\phi)}
 \label{eq:sm_deficit}
\end{equation}
を定める。これは、静電シースが上流へつながるために外から供給されなければならない電場である。
drift のある根は、帯の電位幅が $|\phi_m|$（A）または $|\phi_0|$（C）の指定した割合以下であるか、
$E_{\mathrm{def}}$ が指定した外部の場のレベル $E_{\mathrm{floor}}$ 以下であるときに受理できる。
後者は、上流境界を $E=0$ の無限遠ではなく、場が外部のレベルまで下がった点に置くことに当たる。
どちらの判定も根の方程式を変えないので、受理した根の壁側の量は判定の仕方によらない。
帯より外側だけを物理的な電位分布として使い、分布の再構成も帯の手前で打ち切る。
図\ref{fig:sm_upstream_band}に、drift の有無による上流付近の $E^2$ の違いを示す。BEACH での判定の選び方は\ref{sec:upstream_band}節で述べる。

% Generated by tools/sheath_model_figures.py from upstream_band.csv; edit the generator, not this file.
\begin{figure}[htbp]
\centering
\begin{tikzpicture}[beach diagram]
\draw[black!10,line width=0.4pt] (0,0)--(0,5.6);
\draw[black!10,line width=0.4pt] (3.02,0)--(3.02,5.6);
\draw[black!10,line width=0.4pt] (6.039,0)--(6.039,5.6);
\draw[black!10,line width=0.4pt] (9.059,0)--(9.059,5.6);
\draw[black!10,line width=0.4pt] (0,0.329)--(10.5,0.329);
\draw[black!10,line width=0.4pt] (0,1.647)--(10.5,1.647);
\draw[black!10,line width=0.4pt] (0,2.965)--(10.5,2.965);
\draw[black!10,line width=0.4pt] (0,4.282)--(10.5,4.282);
\draw[black!10,line width=0.4pt] (0,5.6)--(10.5,5.6);
\draw[line width=0.6pt] (0,0) rectangle (10.5,5.6);
\draw[line width=0.6pt] (0,0)--++(0,0.1);
\node[below,font=\scriptsize\sffamily] at (0,0) {$10^{-4}$};
\draw[line width=0.6pt] (3.02,0)--++(0,0.1);
\node[below,font=\scriptsize\sffamily] at (3.02,0) {$10^{-3}$};
\draw[line width=0.6pt] (6.039,0)--++(0,0.1);
\node[below,font=\scriptsize\sffamily] at (6.039,0) {$10^{-2}$};
\draw[line width=0.6pt] (9.059,0)--++(0,0.1);
\node[below,font=\scriptsize\sffamily] at (9.059,0) {$10^{-1}$};
\draw[line width=0.6pt] (0,0.329)--++(0.1,0);
\node[left,font=\scriptsize\sffamily] at (0,0.329) {$-2$};
\draw[line width=0.6pt] (0,1.647)--++(0.1,0);
\node[left,font=\scriptsize\sffamily] at (0,1.647) {$0$};
\draw[line width=0.6pt] (0,2.965)--++(0.1,0);
\node[left,font=\scriptsize\sffamily] at (0,2.965) {$2$};
\draw[line width=0.6pt] (0,4.282)--++(0.1,0);
\node[left,font=\scriptsize\sffamily] at (0,4.282) {$4$};
\draw[line width=0.6pt] (0,5.6)--++(0.1,0);
\node[left,font=\scriptsize\sffamily] at (0,5.6) {$6$};
\node[font=\small\sffamily] at (5.25,-0.75) {上流からの電位の深さ $-\phi$ [V]};
\node[rotate=90,font=\small\sffamily] at (-1.05,2.8) {$\mathrm{sgn}(E^2)\sqrt{|E^2|}$ [mV/m]};
\fill[BeachOrange!12] (0,0) rectangle (7.922,5.6);
\begin{scope}
\clip (0,0) rectangle (10.5,5.6);
\draw[black!45,dashed,line width=0.5pt] (0,1.647)--(10.5,1.647);
\draw[BeachRed,dashed,line width=0.7pt] (0,0.329)--(10.5,0.329);
\draw[BeachBlue,line width=1.2pt] plot coordinates {(-3.02,1.648) (-2.934,1.648) (-2.849,1.648) (-2.764,1.648) (-2.678,1.648) (-2.593,1.648) (-2.508,1.648) (-2.423,1.648) (-2.337,1.648) (-2.252,1.648) (-2.167,1.648) (-2.081,1.648) (-1.996,1.648) (-1.911,1.648) (-1.825,1.648) (-1.74,1.649) (-1.655,1.649) (-1.569,1.649) (-1.484,1.649) (-1.399,1.649) (-1.313,1.649) (-1.228,1.649) (-1.143,1.649) (-1.058,1.65) (-0.972,1.65) (-0.887,1.65) (-0.802,1.65) (-0.716,1.65) (-0.631,1.65) (-0.546,1.651) (-0.46,1.651) (-0.375,1.651) (-0.29,1.651) (-0.204,1.652) (-0.119,1.652) (-0.034,1.652) (0.052,1.653) (0.137,1.653) (0.222,1.654) (0.307,1.654) (0.393,1.654) (0.478,1.655) (0.563,1.655) (0.649,1.656) (0.734,1.657) (0.819,1.657) (0.905,1.658) (0.99,1.659) (1.075,1.66) (1.161,1.66) (1.246,1.661) (1.331,1.662) (1.417,1.663) (1.502,1.664) (1.587,1.665) (1.672,1.667) (1.758,1.668) (1.843,1.669) (1.928,1.671) (2.014,1.673) (2.099,1.674) (2.184,1.676) (2.27,1.678) (2.355,1.68) (2.44,1.682) (2.526,1.685) (2.611,1.687) (2.696,1.69) (2.782,1.693) (2.867,1.696) (2.952,1.699) (3.037,1.703) (3.123,1.706) (3.208,1.71) (3.293,1.715) (3.379,1.719) (3.464,1.724) (3.549,1.729) (3.635,1.735) (3.72,1.74) (3.805,1.747) (3.891,1.753) (3.976,1.761) (4.061,1.768) (4.147,1.776) (4.232,1.785) (4.317,1.794) (4.402,1.804) (4.488,1.815) (4.573,1.826) (4.658,1.838) (4.744,1.851) (4.829,1.864) (4.914,1.879) (5,1.894) (5.085,1.911) (5.17,1.929) (5.256,1.948) (5.341,1.968) (5.426,1.989) (5.512,2.012) (5.597,2.036) (5.682,2.062) (5.767,2.09) (5.853,2.12) (5.938,2.151) (6.023,2.185) (6.109,2.221) (6.194,2.259) (6.279,2.3) (6.365,2.343) (6.45,2.39) (6.535,2.439) (6.621,2.492) (6.706,2.548) (6.791,2.608) (6.877,2.671) (6.962,2.739) (7.047,2.811) (7.132,2.888) (7.218,2.97) (7.303,3.057) (7.388,3.15) (7.474,3.249) (7.559,3.354) (7.644,3.465) (7.73,3.584) (7.815,3.71) (7.9,3.844) (7.986,3.987) (8.071,4.138) (8.156,4.298) (8.242,4.468) (8.327,4.648) (8.412,4.839) (8.497,5.041) (8.583,5.255) (8.668,5.481) (8.753,5.719) (8.839,5.97) (8.924,6.233) (9.009,6.511) (9.095,6.801) (9.18,7.105) (9.265,7.422) (9.351,7.751) (9.436,8.091) (9.521,8.441) (9.607,8.798) (9.692,9.159) (9.777,9.519) (9.862,9.872) (9.948,10.207) (10.033,10.511) (10.118,10.764) (10.204,10.935) (10.289,10.974) (10.374,10.793) (10.46,10.221) (10.545,8.806) (10.63,2.825)};
\draw[BeachOrange,line width=1.2pt] plot coordinates {(-3.02,1.646) (-2.924,1.646) (-2.828,1.646) (-2.732,1.646) (-2.637,1.646) (-2.541,1.646) (-2.445,1.646) (-2.349,1.646) (-2.253,1.646) (-2.158,1.646) (-2.062,1.646) (-1.966,1.646) (-1.87,1.646) (-1.775,1.645) (-1.679,1.645) (-1.583,1.645) (-1.487,1.645) (-1.391,1.645) (-1.296,1.645) (-1.2,1.645) (-1.104,1.644) (-1.008,1.644) (-0.913,1.644) (-0.817,1.644) (-0.721,1.644) (-0.625,1.643) (-0.529,1.643) (-0.434,1.643) (-0.338,1.643) (-0.242,1.642) (-0.146,1.642) (-0.05,1.641) (0.045,1.641) (0.141,1.641) (0.237,1.64) (0.333,1.64) (0.428,1.639) (0.524,1.639) (0.62,1.638) (0.716,1.638) (0.812,1.637) (0.907,1.636) (1.003,1.635) (1.099,1.635) (1.195,1.634) (1.29,1.633) (1.386,1.632) (1.482,1.631) (1.578,1.63) (1.674,1.629) (1.769,1.627) (1.865,1.626) (1.961,1.625) (2.057,1.623) (2.153,1.622) (2.248,1.62) (2.344,1.618) (2.44,1.616) (2.536,1.614) (2.631,1.612) (2.727,1.61) (2.823,1.607) (2.919,1.605) (3.015,1.602) (3.11,1.599) (3.206,1.596) (3.302,1.593) (3.398,1.589) (3.493,1.585) (3.589,1.581) (3.685,1.577) (3.781,1.573) (3.877,1.568) (3.972,1.563) (4.068,1.558) (4.164,1.552) (4.26,1.546) (4.356,1.54) (4.451,1.534) (4.547,1.527) (4.643,1.52) (4.739,1.512) (4.834,1.504) (4.93,1.496) (5.026,1.487) (5.122,1.478) (5.218,1.468) (5.313,1.458) (5.409,1.447) (5.505,1.436) (5.601,1.425) (5.696,1.413) (5.792,1.401) (5.888,1.389) (5.984,1.376) (6.08,1.363) (6.175,1.349) (6.271,1.336) (6.367,1.322) (6.463,1.308) (6.559,1.295) (6.654,1.282) (6.75,1.27) (6.846,1.258) (6.942,1.248) (7.037,1.24) (7.133,1.233) (7.229,1.23) (7.325,1.231) (7.421,1.237) (7.516,1.251) (7.612,1.275) (7.708,1.314) (7.804,1.381) (7.899,1.521) (7.995,1.888) (8.091,2.041) (8.187,2.177) (8.283,2.311) (8.378,2.45) (8.474,2.596) (8.57,2.752) (8.666,2.918) (8.762,3.098) (8.857,3.292) (8.953,3.502) (9.049,3.729) (9.145,3.975) (9.24,4.242) (9.336,4.532) (9.432,4.846) (9.528,5.187) (9.624,5.556) (9.719,5.956) (9.815,6.389) (9.911,6.857) (10.007,7.364) (10.102,7.911) (10.198,8.501) (10.294,9.137) (10.39,9.822) (10.486,10.558) (10.581,11.349) (10.677,12.195) (10.773,13.1) (10.869,14.065) (10.965,15.09) (11.06,16.175) (11.156,17.317) (11.252,18.513) (11.348,19.753) (11.443,21.025) (11.539,22.307) (11.635,23.568) (11.731,24.755) (11.827,25.788) (11.922,26.528) (12.018,26.726) (12.114,25.871) (12.21,22.593) (12.305,5.126)};
\end{scope}
\node[anchor=west,font=\scriptsize\sffamily,text=BeachRed] at (0.532,0.527) {$-E_{\mathrm{floor}}=-2$ mV/m};
\node[anchor=west,font=\scriptsize\sffamily] at (7.922,5.205) {帯の縁 42 mV};
\node[font=\scriptsize\sffamily,text=BeachOrange] at (4.461,0.938) {場の不足 0.63 mV/m};
\draw[BeachBlue,line width=1.2pt] (0.3,5.2)--(0.9,5.2);
\node[right,font=\scriptsize\sffamily] at (0.9,5.2) {drift 0（$\phi_0=3.84$ V）};
\draw[BeachOrange,line width=1.2pt] (0.3,4.7)--(0.9,4.7);
\node[right,font=\scriptsize\sffamily] at (0.9,4.7) {drift あり（$\phi_0=2.65$ V）};
\end{tikzpicture}
\caption[上流の帯と場の不足の数値例]{Type A（$\alpha=60^\circ$）の上側の区間で、上流 $\phi=0$ の近くの $E^2(\phi)$。縦軸は符号を保った $\sqrt{|E^2|}$、横軸は上流からの電位の深さ（対数）である。電子 drift 0 の根は $E^2\propto\phi^2$ で 0 に近づく。太陽風の法線速度で drift した電子では、$h\log(1/h)$ 項（式\eqref{eq:hlogh}）によって上流に接する帯（橙の網掛け）で $E^2<0$ になる。帯の中の最小値の大きさが場の不足であり、この例では下限 $E_{\mathrm{floor}}=2$ mV/m を下回るので根を受理する。drift の有無で根そのもの（$\phi_0$、$\phi_m$）も変わる。}
\label{fig:sm_upstream_band}
\end{figure}


\section{数値アルゴリズム}
\label{sec:sm_algorithm}
図\ref{fig:sm_algorithm}に、一つの型について零電流根を求める手順を示す。

\subsection{規格化と未知数の変換}
電位は光電子温度 $T_{\mathrm{pe}}$（bin の光電子源では $T_e$ 以下の最大の 2 の冪）で、密度は上流のイオン密度で規格化する。
2 の冪を使うのは、規格化の往復で bin の端の値が丸め誤差で動かないようにするためである。
未知数は符号と順序を保つ対数変数に変換する。Type A は $(\log[(\phi_0-\phi_m)/T],\ \log[-\phi_m/T],\ \log[N_e/n_{i\infty}])$、
Type B、C は $(\log[|\phi_0|/T],\ \log[N_e/n_{i\infty}])$ を使い、$\phi_m<\min(\phi_0,0)$ や $N_e>0$ を反復の途中でも保つ。
Type A の接続の残差は $n_{i\infty}(-\phi_m/T)^{3/2}$ で割る。浅い極小では積分区間とその中の電荷密度がともに小さくなり、
規格化しなければ極小が浅いだけで残差が収束したように見えるからである。

\subsection{根の探索}
根は保護付き Newton 法で求める。Jacobian は中心差分（定義域の端では片側差分）で作り、Newton 方向に沿って
残差の最大ノルムが減るまで歩幅を半分にする。減らせなければ減衰係数を増やしながら Levenberg--Marquardt 方向を試す。
定義域外の試行点は評価せずに退け、元の方程式の残差がすべて $10^{-10}$ 以下になったときだけ収束とみなす。
歩幅が小さくなったことや、最小二乗の停留点に達したことは収束と扱わない。

初期値は、呼び出し側が与えた前回の根、近い条件の根を記録した解の表からの予測、既定の多点初期値の順に試す。Type B、C では既定の初期値の前に、次の一変数の探索も行う。
既定の初期値は、温度と密度で無次元化した電位の候補の組であり、$N_e$ は準中性が $N_e$ について線形であることから決める。
$N_e$ を電流の式で消去すると $\phi_0$ だけの方程式になるので、対数と線形の格子で符号変化を探し、二分法で根を求める。
近い条件の根から目標の条件まで物理条件を少しずつ動かす継続法も使え、擬弧長の継続では単純な折り返しを越えられる。
複数の根を列挙するときは、見つけた根から反発する deflation で新しい根を探し、元の方程式と物理解の判定に合格したものだけを残す。
型を指定しない場合は、太陽高度が $20^\circ$ 未満なら C、A、B、それ以外は A、B、C の順に試し、最初に受理した根を返す。
いずれも有限個の初期値と反復による探索であり、根が見つからないことは物理解の不存在を意味しない。

\begin{figure}[htbp]
\centering
\begin{tikzpicture}[beach diagram,node distance=4mm and 6mm,
    flowstep/.style={beach box,text width=46mm,minimum height=10mm,font=\scriptsize\sffamily},
    gate/.style={beach box,text width=46mm,minimum height=10mm,font=\scriptsize\sffamily,fill=BeachOrange!8},
    side/.style={beach box,text width=40mm,minimum height=10mm,font=\scriptsize\sffamily,fill=BeachBlue!6}]
  \node[flowstep] (in) {物理条件を固定\\$n_i,u_i,T_i,T_e,u,$ 光電子源、型};
  \node[flowstep,below=of in] (enc) {未知数を対数変数へ変換\\符号と順序 $\phi_m<\min(\phi_0,0)$ を保証};
  \node[flowstep,below=of enc] (seed) {初期値の候補を並べる\\前回の根 → 解の表 → 既定の多点初期値};
  \node[flowstep,below=of seed] (newton) {保護付き Newton 法\\後退で残差を減らせなければ LM 法};
  \node[gate,below=of newton] (res) {元の方程式の残差 $\le 10^{-10}$ か};
  \node[gate,below=of res] (adm) {物理解の判定\\イオン到達、接続、$E^2\ge0$、上流漸近};
  \node[flowstep,below=of adm,fill=BeachTeal!8] (out) {根を受理し、$z(\phi)$ を再構成\\電流、密度、場の不足を返す};
  \node[side,left=8mm of newton] (scalar) {B、C の一変数探索\\$N_e$ を電流式で消去し、\\$\phi_0$ の格子で符号変化を二分};
  \node[side,right=8mm of res] (next) {次の初期値へ\\（候補列挙では収束済みの根を\\deflation で避ける）};
  \node[side,right=8mm of adm] (rej) {棄却の理由を記録\\非収束と物理的棄却を区別};
  \foreach \a/\b in {in/enc,enc/seed,seed/newton,newton/res,adm/out}
    \draw[beach arrow] (\a)--(\b);
  \draw[beach arrow] (res)--node[right,font=\scriptsize] {是}(adm);
  \draw[beach arrow] (seed.west) -| (scalar.north);
  \draw[beach arrow] (scalar.south) |- (adm.west);
  \draw[beach arrow] (res)--node[above,font=\scriptsize] {否}(next);
  \draw[beach arrow] (adm)--node[above,font=\scriptsize] {否}(rej);
  \draw[beach arrow] (rej)--(next);
  \draw[beach arrow] (next.north) |- (newton.east);
\end{tikzpicture}
\caption[零電流根の求解手順]{一つの型について零電流根を求める手順。橙の箱が受理の判定である。
代数方程式が収束したことと、半無限の物理的な電位分布が存在することを別々に判定する。
型を指定しない場合はこの手順を型の順に繰り返し、最初に受理した根を返す。
有限個の初期値と反復で探すため、見つからないことは解の不存在を意味しない。}
\label{fig:sm_algorithm}
\end{figure}


\subsection{物理解の判定の離散化}
収束した根について、まずイオンの到達と Type A の接続を確かめる。次に各区間で等間隔の 129 点の電位で $E^2$ を評価し、
規格化した $E^2$ の尺度の $10^{-8}$ 倍を超える負の値があれば棄却する。Type B の漸近条件は式\eqref{eq:sm_b_coefficient}の係数で判定し、
丸め誤差に $128$ 倍の機械精度の余裕を与える。drift のある Type A、C では、上流へ向けて
$\phi=-\phi_{\mathrm{ref}}2^{-j}$（$j=1,\dots,48$）の点を加えて帯を見つけ、二分法で帯の縁 $\phi_b$ を決めて、
その幅と場の不足を受理の条件と比べる。下側の区間には帯を許さない。
有限個の点での判定なので、任意の条件に対する存在定理ではない。

\subsection{速度積分と電位積分}
drift が 0 の電子密度は式\eqref{eq:sm_e_density_zero}の閉じた形で計算し、$e^{\psi}\operatorname{erfc}$ の積は
スケールした相補誤差関数 $e^{x^2}\operatorname{erfc}x$ で評価して桁あふれを避ける。drift があるときは、速度積分の区間を
$w=w_{\mathrm{lo}}+(w_{\mathrm{hi}}-w_{\mathrm{lo}})t^2$ で変換し、端点の平方根特異性を除いて複合 Gauss--Legendre 求積で積分する。
速度の上限は drift または切断速度から熱速度の 10 倍先で打ち切る。

第一積分の電位積分は $\phi=\phi_{\mathrm{lo}}+(\phi_{\mathrm{hi}}-\phi_{\mathrm{lo}})\sin^2(\pi t/2)$ で両端に点を集めた
複合 Gauss--Legendre 求積で行う。bin の光電子源の密度の電位積分は $3/2$ 乗の差として解析的に求め、平方根の差は有理化し、
$3/2$ 乗の差は bin 幅と電位の区間で因数分解して評価する。狭い bin や隣り合う浮動小数点数の電位で、
近い数の引き算による桁落ちを避けるためである。

\subsection{電位分布の再構成}
Type B、C では、上流から壁へ $\phi=\phi_0(1-t)^2$ で点を取り、電荷密度の台形積分で $E^2$ を、$1/|E|$ の台形積分で $z$ を求める。
Type A では極小のすぐ外側から両方向へ点を取り、最初の区間だけ極小近くの 2 次の漸近形で距離を与えてから積分する。
上流側は $|\phi|$ が $10^{-3}T_{\mathrm{pe}}$（帯がある場合は帯の縁の少し外側）に達した点で止め、その先に人工的な 0 V の区間を付け足さない。
密度、流束、速度分布は、再構成した電位から式\eqref{eq:sm_e_density}--\eqref{eq:sm_pe_bins}で直接評価する。
このため、表示用の格子の粗さは流束の値に入らない。

\section{数値例}
\label{sec:sm_examples}
本節の図は、表\ref{tab:sm_example_inputs}の条件で外部シースモデルを解いた数値解である。電子の drift は図\ref{fig:sm_upstream_band}を除いて 0 とし、
零電流で閉じた。BEACH の表面帯電計算は含まない。

\begin{table}[htbp]
\centering
\caption{数値例の条件。太陽高度 $\alpha$ は図ごとに示す。}
\label{tab:sm_example_inputs}
\begin{tabular}{lr}
\toprule
量 & 値 \\
\midrule
上流のイオン密度 $n_{i\infty}$（陽子） & 8.7 cm$^{-3}$ \\
太陽風電子の温度 $T_e$ & 12 eV \\
イオン温度 $T_i$ & 0（冷たいビーム） \\
光電子の温度 $T_{\mathrm{pe}}$ & 2.2 eV \\
垂直照射での光電子の規格化 $n_{\mathrm{pe,ref}}$ & 64 cm$^{-3}$ \\
太陽風の速さ $v_{\mathrm{sw}}$ & 468 km/s \\
法線速度と光電子源 & 式\eqref{eq:sm_illumination} \\
電子の drift & 0（図\ref{fig:sm_upstream_band}では $v_{\mathrm{sw}}\sin\alpha$） \\
\bottomrule
\end{tabular}
\end{table}

\subsection{三つの型の電位構造}
図\ref{fig:sm_profiles}は、同じプラズマ条件で太陽高度を変え、Type A、B、C をそれぞれ指定して解いた電位と電場である。
$\alpha=60^\circ$ の Type A は壁電位 3.84 V、壁から 14.9 m に $-0.33$ V の極小を持つ。壁の電場は 1.6 V/m で、極小より上流では向きが反転する。
$\alpha=20^\circ$ の Type B は 1.62 V の壁から光電子の Debye 長程度で減衰し、$\alpha=10^\circ$ の Type C は $-4.24$ V の壁から
太陽風電子の Debye 長程度でゆっくり 0 V に近づく。Type A の上側の電場は 0.01 V/m 程度で、壁近くの電場より 2 桁小さい。

% Generated by tools/sheath_model_figures.py from profile_[ABC].csv; edit the generator, not this file.
\begin{figure}[htbp]
\centering
\begin{tikzpicture}[beach diagram]
\draw[black!10,line width=0.4pt] (0,4.3)--(0,7.9);
\draw[black!10,line width=0.4pt] (1.75,4.3)--(1.75,7.9);
\draw[black!10,line width=0.4pt] (3.5,4.3)--(3.5,7.9);
\draw[black!10,line width=0.4pt] (5.25,4.3)--(5.25,7.9);
\draw[black!10,line width=0.4pt] (7,4.3)--(7,7.9);
\draw[black!10,line width=0.4pt] (8.75,4.3)--(8.75,7.9);
\draw[black!10,line width=0.4pt] (10.5,4.3)--(10.5,7.9);
\draw[black!10,line width=0.4pt] (0,4.66)--(10.5,4.66);
\draw[black!10,line width=0.4pt] (0,5.38)--(10.5,5.38);
\draw[black!10,line width=0.4pt] (0,6.1)--(10.5,6.1);
\draw[black!10,line width=0.4pt] (0,6.82)--(10.5,6.82);
\draw[black!10,line width=0.4pt] (0,7.54)--(10.5,7.54);
\draw[line width=0.6pt] (0,4.3) rectangle (10.5,7.9);
\draw[line width=0.6pt] (0,4.3)--++(0,0.1);
\draw[line width=0.6pt] (1.75,4.3)--++(0,0.1);
\draw[line width=0.6pt] (3.5,4.3)--++(0,0.1);
\draw[line width=0.6pt] (5.25,4.3)--++(0,0.1);
\draw[line width=0.6pt] (7,4.3)--++(0,0.1);
\draw[line width=0.6pt] (8.75,4.3)--++(0,0.1);
\draw[line width=0.6pt] (10.5,4.3)--++(0,0.1);
\draw[line width=0.6pt] (0,4.66)--++(0.1,0);
\node[left,font=\scriptsize\sffamily] at (0,4.66) {$-4$};
\draw[line width=0.6pt] (0,5.38)--++(0.1,0);
\node[left,font=\scriptsize\sffamily] at (0,5.38) {$-2$};
\draw[line width=0.6pt] (0,6.1)--++(0.1,0);
\node[left,font=\scriptsize\sffamily] at (0,6.1) {$0$};
\draw[line width=0.6pt] (0,6.82)--++(0.1,0);
\node[left,font=\scriptsize\sffamily] at (0,6.82) {$2$};
\draw[line width=0.6pt] (0,7.54)--++(0.1,0);
\node[left,font=\scriptsize\sffamily] at (0,7.54) {$4$};
\node[rotate=90,font=\small\sffamily] at (-1.05,6.1) {電位 $\phi$ [V]};
\begin{scope}
\clip (0,4.3) rectangle (10.5,7.9);
\draw[black!45,dashed,line width=0.5pt] (0,6.1)--(10.5,6.1);
\draw[BeachBlue,solid,line width=1.1pt] plot coordinates {(0,7.482) (0.005,7.467) (0.027,7.397) (0.044,7.345) (0.062,7.294) (0.079,7.246) (0.097,7.2) (0.114,7.157) (0.132,7.115) (0.15,7.075) (0.167,7.036) (0.185,7) (0.203,6.965) (0.22,6.931) (0.238,6.898) (0.256,6.867) (0.273,6.837) (0.291,6.808) (0.309,6.781) (0.326,6.754) (0.344,6.729) (0.361,6.704) (0.379,6.68) (0.397,6.657) (0.415,6.634) (0.432,6.613) (0.45,6.592) (0.468,6.572) (0.486,6.552) (0.503,6.533) (0.521,6.515) (0.538,6.498) (0.556,6.481) (0.574,6.465) (0.592,6.448) (0.61,6.433) (0.627,6.418) (0.645,6.404) (0.663,6.39) (0.68,6.376) (0.698,6.363) (0.716,6.351) (0.734,6.338) (0.751,6.326) (0.769,6.315) (0.787,6.304) (0.805,6.293) (0.822,6.283) (0.84,6.272) (0.858,6.262) (0.876,6.253) (0.894,6.243) (0.912,6.234) (0.93,6.226) (0.948,6.217) (0.965,6.209) (0.983,6.201) (1.001,6.193) (1.019,6.186) (1.036,6.179) (1.054,6.172) (1.072,6.165) (1.09,6.158) (1.108,6.152) (1.126,6.145) (1.143,6.14) (1.161,6.134) (1.179,6.128) (1.197,6.122) (1.215,6.117) (1.233,6.112) (1.251,6.107) (1.268,6.102) (1.286,6.098) (1.304,6.093) (1.322,6.089) (1.34,6.084) (1.357,6.08) (1.375,6.076) (1.393,6.072) (1.412,6.068) (1.429,6.065) (1.447,6.061) (1.465,6.058) (1.484,6.054) (1.501,6.051) (1.519,6.048) (1.537,6.045) (1.555,6.042) (1.574,6.04) (1.591,6.037) (1.609,6.034) (1.627,6.032) (1.645,6.029) (1.664,6.027) (1.681,6.025) (1.699,6.023) (1.717,6.02) (1.735,6.018) (1.754,6.016) (1.772,6.014) (1.79,6.013) (1.808,6.011) (1.826,6.009) (1.844,6.008) (1.863,6.006) (1.881,6.004) (1.899,6.003) (1.917,6.002) (1.935,6) (1.953,5.999) (1.972,5.998) (1.99,5.997) (2.008,5.996) (2.026,5.995) (2.044,5.994) (2.062,5.993) (2.08,5.992) (2.099,5.991) (2.117,5.99) (2.135,5.989) (2.153,5.989) (2.171,5.988) (2.189,5.987) (2.207,5.987) (2.225,5.986) (2.244,5.985) (2.263,5.985) (2.281,5.984) (2.298,5.984) (2.316,5.984) (2.334,5.983) (2.353,5.983) (2.371,5.982) (2.39,5.982) (2.409,5.982) (2.428,5.982) (2.445,5.981) (2.463,5.981) (2.481,5.981) (2.499,5.981) (2.517,5.981) (2.535,5.981) (2.553,5.981) (2.572,5.981) (2.619,5.981) (2.637,5.981) (2.659,5.981) (2.681,5.981) (2.699,5.981) (2.719,5.981) (2.74,5.981) (2.758,5.981) (2.778,5.981) (2.797,5.982) (2.816,5.982) (2.834,5.982) (2.854,5.982) (2.872,5.983) (2.891,5.983) (2.91,5.983) (2.928,5.983) (2.946,5.984) (2.964,5.984) (2.984,5.984) (3.002,5.985) (3.02,5.985) (3.039,5.985) (3.057,5.986) (3.075,5.986) (3.094,5.986) (3.112,5.987) (3.13,5.987) (3.148,5.987) (3.166,5.988) (3.184,5.988) (3.202,5.989) (3.221,5.989) (3.239,5.989) (3.257,5.99) (3.275,5.99) (3.293,5.991) (3.311,5.991) (3.329,5.992) (3.347,5.992) (3.365,5.993) (3.383,5.993) (3.401,5.993) (3.419,5.994) (3.437,5.994) (3.455,5.995) (3.473,5.995) (3.49,5.996) (3.509,5.996) (3.527,5.997) (3.545,5.997) (3.563,5.998) (3.581,5.998) (3.598,5.999) (3.616,5.999) (3.634,6) (3.652,6) (3.671,6.001) (3.689,6.001) (3.706,6.002) (3.724,6.002) (3.742,6.003) (3.76,6.003) (3.777,6.004) (3.796,6.004) (3.814,6.005) (3.832,6.006) (3.85,6.006) (3.868,6.007) (3.886,6.007) (3.904,6.008) (3.922,6.008) (3.94,6.009) (3.958,6.009) (3.975,6.01) (3.993,6.01) (4.011,6.011) (4.028,6.011) (4.046,6.012) (4.064,6.012) (4.082,6.013) (4.101,6.013) (4.119,6.014) (4.137,6.014) (4.155,6.015) (4.173,6.015) (4.192,6.016) (4.21,6.016) (4.228,6.017) (4.246,6.018) (4.264,6.018) (4.282,6.019) (4.3,6.019) (4.318,6.02) (4.336,6.02) (4.354,6.021) (4.371,6.021) (4.389,6.022) (4.407,6.022) (4.425,6.023) (4.443,6.023) (4.461,6.024) (4.479,6.024) (4.497,6.025) (4.515,6.025) (4.532,6.026) (4.55,6.026) (4.568,6.027) (4.586,6.027) (4.604,6.028) (4.622,6.028) (4.64,6.028) (4.658,6.029) (4.676,6.029) (4.693,6.03) (4.711,6.03) (4.729,6.031) (4.747,6.031) (4.765,6.032) (4.783,6.032) (4.801,6.033) (4.819,6.033) (4.837,6.034) (4.855,6.034) (4.873,6.035) (4.891,6.035) (4.91,6.036) (4.928,6.036) (4.946,6.036) (4.964,6.037) (4.982,6.037) (5,6.038) (5.019,6.038) (5.037,6.039) (5.055,6.039) (5.074,6.04) (5.092,6.04) (5.109,6.04) (5.127,6.041) (5.145,6.041) (5.162,6.042) (5.18,6.042) (5.197,6.042) (5.215,6.043) (5.233,6.043) (5.251,6.044) (5.268,6.044) (5.286,6.045) (5.304,6.045) (5.322,6.045) (5.34,6.046) (5.358,6.046) (5.376,6.047) (5.394,6.047) (5.412,6.047) (5.43,6.048) (5.448,6.048) (5.466,6.048) (5.484,6.049) (5.503,6.049) (5.521,6.05) (5.539,6.05) (5.558,6.05) (5.576,6.051) (5.594,6.051) (5.611,6.051) (5.629,6.052) (5.647,6.052) (5.664,6.053) (5.682,6.053) (5.7,6.053) (5.718,6.054) (5.735,6.054) (5.753,6.054) (5.771,6.055) (5.789,6.055) (5.807,6.055) (5.825,6.056) (5.844,6.056) (5.862,6.056) (5.88,6.057) (5.898,6.057) (5.917,6.057) (5.935,6.058) (5.954,6.058) (5.971,6.058) (5.989,6.059) (6.006,6.059) (6.024,6.059) (6.042,6.06) (6.059,6.06) (6.077,6.06) (6.095,6.061) (6.113,6.061) (6.131,6.061) (6.149,6.061) (6.167,6.062) (6.185,6.062) (6.203,6.062) (6.221,6.063) (6.24,6.063) (6.258,6.063) (6.276,6.064) (6.295,6.064) (6.313,6.064) (6.331,6.064) (6.348,6.065) (6.366,6.065) (6.384,6.065) (6.402,6.065) (6.419,6.066) (6.437,6.066) (6.455,6.066) (6.473,6.067) (6.491,6.067) (6.509,6.067) (6.527,6.067) (6.545,6.068) (6.564,6.068) (6.582,6.068) (6.6,6.068) (6.619,6.069) (6.637,6.069) (6.656,6.069) (6.674,6.069) (6.692,6.07) (6.709,6.07) (6.727,6.07) (6.745,6.07) (6.762,6.071) (6.78,6.071) (6.798,6.071) (6.816,6.071) (6.834,6.071) (6.852,6.072) (6.87,6.072) (6.889,6.072) (6.907,6.072) (6.925,6.073) (6.943,6.073) (6.962,6.073) (6.98,6.073) (6.999,6.073) (7.018,6.074) (7.035,6.074) (7.053,6.074) (7.07,6.074) (7.088,6.075) (7.106,6.075) (7.124,6.075) (7.141,6.075) (7.159,6.075) (7.177,6.076) (7.195,6.076) (7.213,6.076) (7.232,6.076) (7.25,6.076) (7.268,6.077) (7.287,6.077) (7.305,6.077) (7.324,6.077) (7.342,6.077) (7.361,6.077) (7.379,6.078) (7.397,6.078) (7.414,6.078) (7.432,6.078) (7.45,6.078) (7.468,6.079) (7.485,6.079) (7.503,6.079) (7.521,6.079) (7.539,6.079) (7.557,6.079) (7.575,6.08) (7.594,6.08) (7.612,6.08) (7.63,6.08) (7.649,6.08) (7.667,6.08) (7.686,6.081) (7.704,6.081) (7.723,6.081) (7.742,6.081) (7.76,6.081) (7.778,6.081) (7.795,6.082) (7.813,6.082) (7.831,6.082) (7.849,6.082) (7.866,6.082) (7.884,6.082) (7.902,6.082) (7.92,6.083) (7.938,6.083) (7.957,6.083) (7.975,6.083) (7.993,6.083) (8.011,6.083) (8.03,6.083) (8.048,6.084) (8.067,6.084) (8.085,6.084) (8.104,6.084) (8.123,6.084) (8.142,6.084) (8.161,6.084) (8.178,6.084) (8.196,6.085) (8.213,6.085) (8.231,6.085) (8.249,6.085) (8.267,6.085) (8.285,6.085) (8.303,6.085) (8.321,6.085) (8.339,6.086) (8.357,6.086) (8.375,6.086) (8.393,6.086) (8.412,6.086) (8.43,6.086) (8.449,6.086) (8.467,6.086) (8.486,6.087) (8.505,6.087) (8.523,6.087) (8.542,6.087) (8.561,6.087) (8.58,6.087) (8.599,6.087) (8.617,6.087) (8.634,6.087) (8.652,6.088) (8.67,6.088) (8.688,6.088) (8.706,6.088) (8.724,6.088) (8.742,6.088) (8.76,6.088) (8.778,6.088) (8.796,6.088) (8.814,6.088) (8.833,6.089) (8.851,6.089) (8.869,6.089) (8.888,6.089) (8.907,6.089) (8.925,6.089) (8.944,6.089) (8.963,6.089) (8.982,6.089) (9.001,6.089) (9.02,6.09) (9.039,6.09) (9.058,6.09) (9.076,6.09) (9.093,6.09) (9.111,6.09) (9.129,6.09) (9.146,6.09) (9.164,6.09) (9.182,6.09) (9.2,6.09) (9.218,6.09) (9.236,6.091) (9.254,6.091) (9.272,6.091) (9.291,6.091) (9.309,6.091) (9.328,6.091) (9.346,6.091) (9.365,6.091) (9.383,6.091) (9.402,6.091) (9.421,6.091) (9.44,6.091) (9.459,6.091) (9.478,6.092) (9.497,6.092) (9.516,6.092) (9.535,6.092) (9.554,6.092) (9.574,6.092) (9.593,6.092) (9.611,6.092) (9.628,6.092) (9.646,6.092) (9.664,6.092) (9.682,6.092) (9.699,6.092) (9.717,6.092) (9.735,6.092) (9.753,6.093) (9.772,6.093) (9.79,6.093) (9.808,6.093) (9.826,6.093) (9.845,6.093) (9.863,6.093) (9.882,6.093) (9.9,6.093) (9.919,6.093) (9.938,6.093) (9.957,6.093) (9.975,6.093) (9.994,6.093) (10.013,6.093) (10.032,6.093) (10.052,6.094) (10.071,6.094) (10.09,6.094) (10.11,6.094) (10.129,6.094) (10.148,6.094) (10.168,6.094) (10.188,6.094) (10.205,6.094) (10.223,6.094) (10.241,6.094) (10.258,6.094) (10.276,6.094) (10.294,6.094) (10.312,6.094) (10.33,6.094) (10.348,6.094) (10.366,6.094) (10.384,6.094) (10.402,6.095) (10.421,6.095) (10.439,6.095) (10.457,6.095) (10.476,6.095) (10.494,6.095) (10.513,6.095) (10.532,6.095) (10.55,6.095) (10.569,6.095) (10.588,6.095) (10.607,6.095) (10.626,6.095) (10.645,6.095) (10.664,6.095) (10.683,6.095) (10.703,6.095) (10.722,6.095) (10.741,6.095) (10.761,6.095) (10.78,6.095) (10.8,6.095) (10.82,6.095) (10.84,6.096) (10.859,6.096) (10.879,6.096) (10.899,6.096) (10.917,6.096) (10.934,6.096) (10.952,6.096) (10.97,6.096) (10.987,6.096) (11.005,6.096) (11.023,6.096) (11.041,6.096) (11.059,6.096) (11.077,6.096) (11.095,6.096) (11.113,6.096) (11.131,6.096) (11.15,6.096) (11.168,6.096) (11.186,6.096) (11.205,6.096) (11.223,6.096) (11.242,6.096) (11.26,6.096) (11.279,6.096) (11.298,6.096) (11.316,6.096) (11.335,6.096) (11.354,6.096) (11.373,6.097) (11.392,6.097) (11.411,6.097) (11.43,6.097) (11.449,6.097) (11.469,6.097) (11.488,6.097) (11.507,6.097) (11.527,6.097) (11.546,6.097) (11.566,6.097) (11.585,6.097) (11.605,6.097) (11.625,6.097) (11.644,6.097) (11.664,6.097) (11.684,6.097) (11.704,6.097) (11.724,6.097) (11.744,6.097) (11.764,6.097) (11.784,6.097) (11.804,6.097) (11.825,6.097) (11.845,6.097) (11.865,6.097) (11.883,6.097) (11.9,6.097) (11.918,6.097) (11.936,6.097) (11.953,6.097) (11.971,6.097) (11.989,6.097) (12.006,6.097) (12.024,6.097) (12.042,6.097) (12.06,6.097) (12.078,6.098) (12.096,6.098) (12.114,6.098) (12.132,6.098) (12.15,6.098) (12.168,6.098) (12.186,6.098) (12.205,6.098) (12.223,6.098) (12.241,6.098) (12.259,6.098) (12.278,6.098) (12.296,6.098) (12.315,6.098) (12.333,6.098) (12.351,6.098) (12.37,6.098) (12.389,6.098) (12.407,6.098) (12.426,6.098) (12.444,6.098) (12.463,6.098) (12.482,6.098) (12.5,6.098) (12.519,6.098) (12.538,6.098) (12.557,6.098) (12.576,6.098) (12.594,6.098) (12.613,6.098) (12.632,6.098) (12.651,6.098) (12.67,6.098) (12.689,6.098) (12.708,6.098) (12.727,6.098) (12.746,6.098) (12.765,6.098) (12.784,6.098) (12.803,6.098) (12.822,6.098) (12.841,6.098) (12.86,6.098) (12.879,6.098) (12.898,6.098) (12.917,6.098) (12.936,6.098) (12.955,6.098) (12.974,6.098) (12.993,6.098) (13.012,6.098) (13.032,6.098) (13.051,6.098) (13.07,6.098) (13.089,6.098) (13.108,6.098) (13.127,6.098) (13.146,6.099) (13.165,6.099) (13.184,6.099) (13.203,6.099) (13.221,6.099) (13.24,6.099) (13.259,6.099) (13.278,6.099) (13.297,6.099) (13.316,6.099) (13.334,6.099) (13.353,6.099) (13.372,6.099) (13.39,6.099) (13.409,6.099) (13.427,6.099) (13.446,6.099) (13.464,6.099) (13.482,6.099) (13.501,6.099) (13.519,6.099) (13.537,6.099) (13.555,6.099) (13.573,6.099) (13.591,6.099) (13.609,6.099) (13.626,6.099) (13.644,6.099) (13.661,6.099) (13.682,6.099) (13.702,6.099) (13.722,6.099) (13.742,6.099) (13.762,6.099) (13.781,6.099) (13.801,6.099) (13.82,6.099) (13.839,6.099) (13.858,6.099) (13.877,6.099) (13.895,6.099) (13.914,6.099) (13.932,6.099) (13.95,6.099) (13.967,6.099) (13.987,6.099) (14.007,6.099) (14.026,6.099) (14.045,6.099) (14.064,6.099) (14.082,6.099) (14.1,6.099) (14.118,6.099) (14.137,6.099) (14.156,6.099) (14.174,6.099) (14.192,6.099) (14.211,6.099) (14.229,6.099) (14.247,6.099) (14.266,6.099) (14.284,6.099) (14.302,6.099) (14.321,6.099) (14.339,6.099) (14.357,6.099) (14.375,6.099) (14.394,6.099) (14.411,6.099) (14.429,6.099) (14.431,6.099)};
\draw[BeachOrange,dashed,line width=1.1pt] plot coordinates {(0,6.682) (0.01,6.669) (0.02,6.657) (0.031,6.645) (0.049,6.624) (0.067,6.604) (0.084,6.585) (0.102,6.567) (0.121,6.549) (0.139,6.532) (0.157,6.515) (0.174,6.5) (0.192,6.485) (0.21,6.47) (0.228,6.456) (0.246,6.443) (0.264,6.429) (0.283,6.416) (0.3,6.404) (0.318,6.392) (0.337,6.381) (0.354,6.37) (0.372,6.36) (0.39,6.349) (0.409,6.339) (0.427,6.33) (0.445,6.321) (0.463,6.312) (0.481,6.303) (0.499,6.295) (0.517,6.288) (0.535,6.28) (0.553,6.273) (0.572,6.265) (0.589,6.259) (0.607,6.252) (0.626,6.246) (0.644,6.239) (0.663,6.233) (0.681,6.228) (0.698,6.222) (0.717,6.217) (0.735,6.212) (0.754,6.207) (0.773,6.202) (0.791,6.198) (0.809,6.193) (0.827,6.189) (0.846,6.185) (0.865,6.181) (0.884,6.177) (0.902,6.173) (0.92,6.17) (0.938,6.167) (0.957,6.163) (0.976,6.16) (0.995,6.157) (1.013,6.155) (1.031,6.152) (1.049,6.149) (1.067,6.147) (1.086,6.144) (1.105,6.142) (1.125,6.14) (1.142,6.138) (1.16,6.136) (1.178,6.134) (1.197,6.132) (1.216,6.13) (1.235,6.128) (1.254,6.127) (1.274,6.125) (1.292,6.124) (1.31,6.122) (1.328,6.121) (1.347,6.12) (1.366,6.119) (1.385,6.117) (1.405,6.116) (1.425,6.115) (1.442,6.114) (1.46,6.113) (1.478,6.112) (1.497,6.112) (1.516,6.111) (1.535,6.11) (1.555,6.109) (1.575,6.109) (1.596,6.108) (1.614,6.107) (1.632,6.107) (1.65,6.106) (1.669,6.106) (1.688,6.105) (1.708,6.105) (1.728,6.104) (1.749,6.104) (1.77,6.104) (1.792,6.103) (1.81,6.103) (1.828,6.103) (1.847,6.102) (1.867,6.102) (1.887,6.102) (1.907,6.102) (1.929,6.102) (1.951,6.101) (1.974,6.101) (1.991,6.101) (2.01,6.101) (2.022,6.101)};
\draw[BeachTeal,dash dot,line width=1.1pt] plot coordinates {(0,4.574) (0.015,4.585) (0.03,4.595) (0.045,4.605) (0.06,4.615) (0.075,4.625) (0.09,4.635) (0.105,4.645) (0.12,4.655) (0.135,4.665) (0.165,4.684) (0.183,4.697) (0.202,4.709) (0.221,4.721) (0.24,4.733) (0.258,4.745) (0.277,4.757) (0.296,4.769) (0.315,4.781) (0.334,4.792) (0.353,4.804) (0.372,4.816) (0.391,4.828) (0.41,4.839) (0.429,4.851) (0.448,4.862) (0.467,4.874) (0.486,4.885) (0.506,4.896) (0.525,4.908) (0.544,4.919) (0.564,4.93) (0.583,4.941) (0.603,4.952) (0.622,4.963) (0.642,4.974) (0.662,4.985) (0.682,4.996) (0.702,5.007) (0.721,5.017) (0.741,5.028) (0.762,5.039) (0.782,5.049) (0.802,5.06) (0.822,5.07) (0.843,5.081) (0.863,5.091) (0.884,5.102) (0.904,5.112) (0.925,5.122) (0.946,5.132) (0.966,5.142) (0.987,5.152) (1.008,5.162) (1.029,5.172) (1.051,5.182) (1.072,5.192) (1.093,5.202) (1.115,5.212) (1.136,5.221) (1.158,5.231) (1.18,5.24) (1.201,5.25) (1.219,5.258) (1.237,5.265) (1.254,5.273) (1.272,5.28) (1.29,5.288) (1.308,5.295) (1.326,5.302) (1.344,5.31) (1.362,5.317) (1.38,5.324) (1.398,5.331) (1.416,5.339) (1.435,5.346) (1.453,5.353) (1.472,5.36) (1.49,5.367) (1.509,5.374) (1.528,5.381) (1.546,5.388) (1.565,5.395) (1.584,5.402) (1.603,5.409) (1.622,5.416) (1.642,5.423) (1.661,5.429) (1.68,5.436) (1.7,5.443) (1.719,5.449) (1.739,5.456) (1.759,5.463) (1.778,5.469) (1.798,5.476) (1.818,5.482) (1.838,5.489) (1.859,5.495) (1.879,5.501) (1.899,5.508) (1.92,5.514) (1.94,5.52) (1.961,5.527) (1.982,5.533) (2.003,5.539) (2.024,5.545) (2.045,5.551) (2.066,5.557) (2.087,5.564) (2.109,5.57) (2.13,5.576) (2.152,5.581) (2.174,5.587) (2.195,5.593) (2.217,5.599) (2.24,5.605) (2.262,5.611) (2.284,5.616) (2.306,5.622) (2.329,5.628) (2.352,5.634) (2.375,5.639) (2.398,5.645) (2.421,5.65) (2.444,5.656) (2.467,5.661) (2.485,5.665) (2.502,5.669) (2.52,5.673) (2.538,5.677) (2.556,5.681) (2.574,5.685) (2.592,5.689) (2.61,5.693) (2.628,5.697) (2.647,5.701) (2.665,5.705) (2.684,5.709) (2.702,5.713) (2.721,5.717) (2.74,5.721) (2.758,5.724) (2.777,5.728) (2.796,5.732) (2.815,5.736) (2.835,5.739) (2.854,5.743) (2.873,5.747) (2.893,5.75) (2.912,5.754) (2.932,5.758) (2.952,5.761) (2.972,5.765) (2.992,5.768) (3.012,5.772) (3.032,5.775) (3.052,5.779) (3.072,5.782) (3.093,5.786) (3.114,5.789) (3.134,5.793) (3.155,5.796) (3.176,5.8) (3.197,5.803) (3.218,5.806) (3.239,5.81) (3.261,5.813) (3.282,5.816) (3.304,5.82) (3.326,5.823) (3.348,5.826) (3.37,5.829) (3.392,5.833) (3.414,5.836) (3.436,5.839) (3.459,5.842) (3.481,5.845) (3.504,5.848) (3.527,5.851) (3.55,5.854) (3.573,5.858) (3.597,5.861) (3.62,5.864) (3.644,5.867) (3.667,5.87) (3.691,5.873) (3.715,5.875) (3.739,5.878) (3.764,5.881) (3.788,5.884) (3.813,5.887) (3.838,5.89) (3.863,5.893) (3.888,5.895) (3.913,5.898) (3.939,5.901) (3.964,5.904) (3.99,5.907) (4.016,5.909) (4.042,5.912) (4.06,5.914) (4.078,5.915) (4.095,5.917) (4.113,5.919) (4.131,5.921) (4.149,5.922) (4.167,5.924) (4.185,5.926) (4.204,5.928) (4.222,5.929) (4.24,5.931) (4.259,5.933) (4.277,5.934) (4.296,5.936) (4.315,5.938) (4.334,5.939) (4.353,5.941) (4.372,5.943) (4.391,5.944) (4.41,5.946) (4.43,5.948) (4.449,5.949) (4.469,5.951) (4.488,5.952) (4.508,5.954) (4.528,5.956) (4.548,5.957) (4.568,5.959) (4.588,5.96) (4.608,5.962) (4.629,5.963) (4.649,5.965) (4.67,5.966) (4.691,5.968) (4.712,5.969) (4.733,5.971) (4.754,5.972) (4.775,5.974) (4.796,5.975) (4.818,5.977) (4.839,5.978) (4.861,5.979) (4.883,5.981) (4.905,5.982) (4.927,5.984) (4.949,5.985) (4.971,5.987) (4.994,5.988) (5.017,5.989) (5.039,5.991) (5.062,5.992) (5.085,5.993) (5.109,5.995) (5.132,5.996) (5.155,5.997) (5.179,5.999) (5.203,6) (5.227,6.001) (5.251,6.003) (5.275,6.004) (5.3,6.005) (5.324,6.006) (5.349,6.008) (5.374,6.009) (5.399,6.01) (5.424,6.011) (5.449,6.013) (5.475,6.014) (5.501,6.015) (5.527,6.016) (5.553,6.017) (5.579,6.019) (5.606,6.02) (5.632,6.021) (5.659,6.022) (5.686,6.023) (5.714,6.024) (5.741,6.025) (5.769,6.027) (5.797,6.028) (5.825,6.029) (5.853,6.03) (5.882,6.031) (5.911,6.032) (5.94,6.033) (5.969,6.034) (5.998,6.035) (6.028,6.036) (6.058,6.037) (6.088,6.038) (6.119,6.039) (6.149,6.04) (6.18,6.041) (6.212,6.042) (6.243,6.043) (6.275,6.044) (6.307,6.045) (6.339,6.046) (6.372,6.047) (6.405,6.048) (6.438,6.049) (6.471,6.05) (6.505,6.051) (6.539,6.052) (6.574,6.053) (6.608,6.054) (6.626,6.054) (6.643,6.055) (6.661,6.055) (6.679,6.055) (6.697,6.056) (6.715,6.056) (6.733,6.057) (6.751,6.057) (6.769,6.058) (6.787,6.058) (6.806,6.058) (6.824,6.059) (6.843,6.059) (6.861,6.06) (6.88,6.06) (6.899,6.06) (6.918,6.061) (6.937,6.061) (6.956,6.062) (6.975,6.062) (6.995,6.062) (7.014,6.063) (7.034,6.063) (7.053,6.064) (7.073,6.064) (7.093,6.064) (7.113,6.065) (7.133,6.065) (7.153,6.066) (7.174,6.066) (7.194,6.066) (7.214,6.067) (7.235,6.067) (7.256,6.068) (7.277,6.068) (7.298,6.068) (7.319,6.069) (7.34,6.069) (7.362,6.069) (7.383,6.07) (7.405,6.07) (7.426,6.07) (7.448,6.071) (7.47,6.071) (7.493,6.071) (7.515,6.072) (7.537,6.072) (7.56,6.072) (7.583,6.073) (7.605,6.073) (7.628,6.074) (7.652,6.074) (7.675,6.074) (7.698,6.075) (7.722,6.075) (7.746,6.075) (7.769,6.075) (7.793,6.076) (7.818,6.076) (7.842,6.076) (7.867,6.077) (7.891,6.077) (7.916,6.077) (7.941,6.078) (7.966,6.078) (7.992,6.078) (8.017,6.079) (8.043,6.079) (8.069,6.079) (8.095,6.079) (8.121,6.08) (8.148,6.08) (8.175,6.08) (8.201,6.081) (8.229,6.081) (8.256,6.081) (8.283,6.082) (8.311,6.082) (8.339,6.082) (8.367,6.082) (8.395,6.083) (8.424,6.083) (8.453,6.083) (8.482,6.083) (8.511,6.084) (8.54,6.084) (8.57,6.084) (8.6,6.084) (8.63,6.085) (8.661,6.085) (8.692,6.085) (8.723,6.085) (8.754,6.086) (8.785,6.086) (8.817,6.086) (8.849,6.086) (8.882,6.087) (8.914,6.087) (8.947,6.087) (8.98,6.087) (9.014,6.088) (9.048,6.088) (9.082,6.088) (9.117,6.088) (9.151,6.089) (9.187,6.089) (9.222,6.089) (9.258,6.089) (9.294,6.089) (9.331,6.09) (9.368,6.09) (9.405,6.09) (9.443,6.09) (9.481,6.09) (9.519,6.091) (9.558,6.091) (9.598,6.091) (9.637,6.091) (9.678,6.091) (9.718,6.092) (9.759,6.092) (9.801,6.092) (9.843,6.092) (9.886,6.092) (9.929,6.093) (9.972,6.093) (10.016,6.093) (10.061,6.093) (10.106,6.093) (10.152,6.093) (10.198,6.094) (10.245,6.094) (10.293,6.094) (10.341,6.094) (10.39,6.094) (10.439,6.094) (10.49,6.094) (10.541,6.095) (10.592,6.095) (10.645,6.095) (10.698,6.095) (10.752,6.095) (10.806,6.095) (10.862,6.096) (10.918,6.096) (10.976,6.096) (11.034,6.096) (11.093,6.096) (11.153,6.096) (11.214,6.096) (11.276,6.096) (11.34,6.097) (11.404,6.097) (11.47,6.097) (11.536,6.097) (11.604,6.097) (11.674,6.097) (11.744,6.097) (11.816,6.097) (11.89,6.097) (11.965,6.098) (12.042,6.098) (12.12,6.098) (12.2,6.098) (12.282,6.098) (12.365,6.098) (12.451,6.098) (12.539,6.098) (12.628,6.098) (12.721,6.098) (12.815,6.098) (12.912,6.099) (13.012,6.099) (13.114,6.099) (13.22,6.099) (13.328,6.099) (13.44,6.099) (13.556,6.099) (13.675,6.099) (13.799,6.099) (13.926,6.099)};
\end{scope}
\draw[BeachBlue,dotted,line width=0.6pt] (2.603,4.3)--(2.603,5.981);
\node[beach dot,fill=BeachBlue] at (2.603,5.981) {};
\node[anchor=west,font=\scriptsize\sffamily,text=BeachBlue] at (2.743,6.305) {$z_m$, $\phi_m$};
\draw[black!10,line width=0.4pt] (0,0)--(0,3.6);
\draw[black!10,line width=0.4pt] (1.75,0)--(1.75,3.6);
\draw[black!10,line width=0.4pt] (3.5,0)--(3.5,3.6);
\draw[black!10,line width=0.4pt] (5.25,0)--(5.25,3.6);
\draw[black!10,line width=0.4pt] (7,0)--(7,3.6);
\draw[black!10,line width=0.4pt] (8.75,0)--(8.75,3.6);
\draw[black!10,line width=0.4pt] (10.5,0)--(10.5,3.6);
\draw[black!10,line width=0.4pt] (0,0)--(10.5,0);
\draw[black!10,line width=0.4pt] (0,0.8)--(10.5,0.8);
\draw[black!10,line width=0.4pt] (0,1.6)--(10.5,1.6);
\draw[black!10,line width=0.4pt] (0,2.4)--(10.5,2.4);
\draw[black!10,line width=0.4pt] (0,3.2)--(10.5,3.2);
\draw[line width=0.6pt] (0,0) rectangle (10.5,3.6);
\draw[line width=0.6pt] (0,0)--++(0,0.1);
\node[below,font=\scriptsize\sffamily] at (0,0) {$0$};
\draw[line width=0.6pt] (1.75,0)--++(0,0.1);
\node[below,font=\scriptsize\sffamily] at (1.75,0) {$10$};
\draw[line width=0.6pt] (3.5,0)--++(0,0.1);
\node[below,font=\scriptsize\sffamily] at (3.5,0) {$20$};
\draw[line width=0.6pt] (5.25,0)--++(0,0.1);
\node[below,font=\scriptsize\sffamily] at (5.25,0) {$30$};
\draw[line width=0.6pt] (7,0)--++(0,0.1);
\node[below,font=\scriptsize\sffamily] at (7,0) {$40$};
\draw[line width=0.6pt] (8.75,0)--++(0,0.1);
\node[below,font=\scriptsize\sffamily] at (8.75,0) {$50$};
\draw[line width=0.6pt] (10.5,0)--++(0,0.1);
\node[below,font=\scriptsize\sffamily] at (10.5,0) {$60$};
\draw[line width=0.6pt] (0,0)--++(0.1,0);
\node[left,font=\scriptsize\sffamily] at (0,0) {$-0.5$};
\draw[line width=0.6pt] (0,0.8)--++(0.1,0);
\node[left,font=\scriptsize\sffamily] at (0,0.8) {$0$};
\draw[line width=0.6pt] (0,1.6)--++(0.1,0);
\node[left,font=\scriptsize\sffamily] at (0,1.6) {$0.5$};
\draw[line width=0.6pt] (0,2.4)--++(0.1,0);
\node[left,font=\scriptsize\sffamily] at (0,2.4) {$1$};
\draw[line width=0.6pt] (0,3.2)--++(0.1,0);
\node[left,font=\scriptsize\sffamily] at (0,3.2) {$1.5$};
\node[font=\small\sffamily] at (5.25,-0.75) {壁からの距離 $z$ [m]};
\node[rotate=90,font=\small\sffamily] at (-1.05,1.8) {電場 $E_z$ [V/m]};
\begin{scope}
\clip (0,0) rectangle (10.5,3.6);
\draw[black!45,dashed,line width=0.5pt] (0,0.8)--(10.5,0.8);
\draw[BeachBlue,solid,line width=1.1pt] plot coordinates {(0,3.398) (0.002,3.382) (0.004,3.366) (0.006,3.35) (0.008,3.334) (0.027,3.2) (0.044,3.082) (0.062,2.972) (0.079,2.87) (0.097,2.775) (0.114,2.687) (0.132,2.603) (0.15,2.525) (0.167,2.452) (0.185,2.384) (0.203,2.32) (0.22,2.259) (0.238,2.201) (0.256,2.146) (0.273,2.095) (0.291,2.046) (0.309,2) (0.326,1.955) (0.344,1.914) (0.361,1.874) (0.379,1.836) (0.397,1.799) (0.415,1.763) (0.432,1.73) (0.45,1.698) (0.468,1.668) (0.486,1.638) (0.503,1.61) (0.521,1.583) (0.538,1.557) (0.556,1.532) (0.574,1.509) (0.592,1.485) (0.61,1.463) (0.627,1.442) (0.645,1.422) (0.663,1.402) (0.68,1.383) (0.698,1.364) (0.716,1.347) (0.734,1.33) (0.751,1.313) (0.769,1.297) (0.787,1.282) (0.805,1.267) (0.822,1.253) (0.84,1.238) (0.858,1.225) (0.876,1.212) (0.894,1.199) (0.912,1.187) (0.93,1.175) (0.948,1.164) (0.965,1.152) (0.983,1.142) (1.001,1.131) (1.019,1.121) (1.036,1.112) (1.054,1.102) (1.072,1.093) (1.09,1.084) (1.108,1.075) (1.126,1.067) (1.143,1.059) (1.161,1.051) (1.179,1.043) (1.197,1.035) (1.215,1.028) (1.233,1.021) (1.251,1.014) (1.268,1.008) (1.286,1.001) (1.304,0.995) (1.322,0.989) (1.34,0.983) (1.357,0.977) (1.375,0.972) (1.393,0.966) (1.412,0.961) (1.429,0.956) (1.447,0.951) (1.465,0.946) (1.484,0.941) (1.501,0.936) (1.519,0.932) (1.537,0.927) (1.555,0.923) (1.574,0.919) (1.591,0.915) (1.609,0.911) (1.627,0.907) (1.645,0.903) (1.664,0.9) (1.681,0.896) (1.699,0.893) (1.717,0.889) (1.735,0.886) (1.754,0.883) (1.772,0.88) (1.79,0.877) (1.808,0.874) (1.826,0.871) (1.844,0.868) (1.863,0.865) (1.881,0.863) (1.899,0.86) (1.917,0.858) (1.935,0.855) (1.953,0.853) (1.972,0.851) (1.99,0.848) (2.008,0.846) (2.026,0.844) (2.044,0.842) (2.062,0.84) (2.08,0.838) (2.099,0.836) (2.117,0.834) (2.135,0.833) (2.153,0.831) (2.171,0.829) (2.189,0.828) (2.207,0.826) (2.225,0.824) (2.244,0.823) (2.263,0.821) (2.281,0.82) (2.298,0.818) (2.316,0.817) (2.334,0.816) (2.353,0.814) (2.371,0.813) (2.39,0.812) (2.409,0.811) (2.428,0.81) (2.445,0.808) (2.463,0.807) (2.481,0.806) (2.499,0.805) (2.517,0.804) (2.535,0.803) (2.553,0.802) (2.572,0.801) (2.619,0.799) (2.637,0.798) (2.659,0.797) (2.681,0.796) (2.699,0.796) (2.719,0.795) (2.74,0.794) (2.758,0.793) (2.778,0.793) (2.797,0.792) (2.816,0.791) (2.834,0.791) (2.854,0.79) (2.872,0.79) (2.891,0.789) (2.91,0.789) (2.928,0.788) (2.946,0.788) (2.964,0.787) (2.984,0.787) (3.002,0.786) (3.02,0.786) (3.039,0.785) (3.057,0.785) (3.075,0.785) (3.094,0.784) (3.112,0.784) (3.13,0.784) (3.148,0.783) (3.166,0.783) (3.184,0.783) (3.202,0.782) (3.221,0.782) (3.239,0.782) (3.257,0.782) (3.275,0.781) (3.293,0.781) (3.311,0.781) (3.329,0.781) (3.347,0.78) (3.365,0.78) (3.383,0.78) (3.401,0.78) (3.419,0.78) (3.437,0.779) (3.455,0.779) (3.473,0.779) (3.49,0.779) (3.509,0.779) (3.527,0.779) (3.545,0.779) (3.563,0.778) (3.581,0.778) (3.598,0.778) (3.616,0.778) (3.634,0.778) (3.652,0.778) (3.671,0.778) (3.689,0.778) (3.706,0.778) (3.724,0.778) (3.742,0.778) (3.76,0.778) (3.777,0.778) (3.796,0.778) (3.814,0.777) (3.832,0.777) (3.85,0.777) (3.868,0.777) (3.886,0.777) (3.904,0.777) (3.922,0.777) (3.94,0.777) (3.958,0.777) (3.975,0.777) (3.993,0.777) (4.011,0.777) (4.028,0.777) (4.046,0.777) (4.064,0.777) (4.082,0.777) (4.101,0.777) (4.119,0.777) (4.137,0.777) (4.155,0.778) (4.173,0.778) (4.192,0.778) (4.21,0.778) (4.228,0.778) (4.246,0.778) (4.264,0.778) (4.282,0.778) (4.3,0.778) (4.318,0.778) (4.336,0.778) (4.354,0.778) (4.371,0.778) (4.389,0.778) (4.407,0.778) (4.425,0.778) (4.443,0.778) (4.461,0.778) (4.479,0.778) (4.497,0.779) (4.515,0.779) (4.532,0.779) (4.55,0.779) (4.568,0.779) (4.586,0.779) (4.604,0.779) (4.622,0.779) (4.64,0.779) (4.658,0.779) (4.676,0.779) (4.693,0.779) (4.711,0.779) (4.729,0.78) (4.747,0.78) (4.765,0.78) (4.783,0.78) (4.801,0.78) (4.819,0.78) (4.837,0.78) (4.855,0.78) (4.873,0.78) (4.891,0.78) (4.91,0.78) (4.928,0.781) (4.946,0.781) (4.964,0.781) (4.982,0.781) (5,0.781) (5.019,0.781) (5.037,0.781) (5.055,0.781) (5.074,0.781) (5.092,0.781) (5.109,0.781) (5.127,0.782) (5.145,0.782) (5.162,0.782) (5.18,0.782) (5.197,0.782) (5.215,0.782) (5.233,0.782) (5.251,0.782) (5.268,0.782) (5.286,0.782) (5.304,0.783) (5.322,0.783) (5.34,0.783) (5.358,0.783) (5.376,0.783) (5.394,0.783) (5.412,0.783) (5.43,0.783) (5.448,0.783) (5.466,0.783) (5.484,0.784) (5.503,0.784) (5.521,0.784) (5.539,0.784) (5.558,0.784) (5.576,0.784) (5.594,0.784) (5.611,0.784) (5.629,0.784) (5.647,0.784) (5.664,0.784) (5.682,0.785) (5.7,0.785) (5.718,0.785) (5.735,0.785) (5.753,0.785) (5.771,0.785) (5.789,0.785) (5.807,0.785) (5.825,0.785) (5.844,0.785) (5.862,0.786) (5.88,0.786) (5.898,0.786) (5.917,0.786) (5.935,0.786) (5.954,0.786) (5.971,0.786) (5.989,0.786) (6.006,0.786) (6.024,0.786) (6.042,0.786) (6.059,0.787) (6.077,0.787) (6.095,0.787) (6.113,0.787) (6.131,0.787) (6.149,0.787) (6.167,0.787) (6.185,0.787) (6.203,0.787) (6.221,0.787) (6.24,0.787) (6.258,0.788) (6.276,0.788) (6.295,0.788) (6.313,0.788) (6.331,0.788) (6.348,0.788) (6.366,0.788) (6.384,0.788) (6.402,0.788) (6.419,0.788) (6.437,0.788) (6.455,0.788) (6.473,0.789) (6.491,0.789) (6.509,0.789) (6.527,0.789) (6.545,0.789) (6.564,0.789) (6.582,0.789) (6.6,0.789) (6.619,0.789) (6.637,0.789) (6.656,0.789) (6.674,0.789) (6.692,0.79) (6.709,0.79) (6.727,0.79) (6.745,0.79) (6.762,0.79) (6.78,0.79) (6.798,0.79) (6.816,0.79) (6.834,0.79) (6.852,0.79) (6.87,0.79) (6.889,0.79) (6.907,0.79) (6.925,0.79) (6.943,0.791) (6.962,0.791) (6.98,0.791) (6.999,0.791) (7.018,0.791) (7.035,0.791) (7.053,0.791) (7.07,0.791) (7.088,0.791) (7.106,0.791) (7.124,0.791) (7.141,0.791) (7.159,0.791) (7.177,0.791) (7.195,0.791) (7.213,0.792) (7.232,0.792) (7.25,0.792) (7.268,0.792) (7.287,0.792) (7.305,0.792) (7.324,0.792) (7.342,0.792) (7.361,0.792) (7.379,0.792) (7.397,0.792) (7.414,0.792) (7.432,0.792) (7.45,0.792) (7.468,0.792) (7.485,0.792) (7.503,0.793) (7.521,0.793) (7.539,0.793) (7.557,0.793) (7.575,0.793) (7.594,0.793) (7.612,0.793) (7.63,0.793) (7.649,0.793) (7.667,0.793) (7.686,0.793) (7.704,0.793) (7.723,0.793) (7.742,0.793) (7.76,0.793) (7.778,0.793) (7.795,0.793) (7.813,0.793) (7.831,0.794) (7.849,0.794) (7.866,0.794) (7.884,0.794) (7.902,0.794) (7.92,0.794) (7.938,0.794) (7.957,0.794) (7.975,0.794) (7.993,0.794) (8.011,0.794) (8.03,0.794) (8.048,0.794) (8.067,0.794) (8.085,0.794) (8.104,0.794) (8.123,0.794) (8.142,0.794) (8.161,0.794) (8.178,0.794) (8.196,0.794) (8.213,0.795) (8.231,0.795) (8.249,0.795) (8.267,0.795) (8.285,0.795) (8.303,0.795) (8.321,0.795) (8.339,0.795) (8.357,0.795) (8.375,0.795) (8.393,0.795) (8.412,0.795) (8.43,0.795) (8.449,0.795) (8.467,0.795) (8.486,0.795) (8.505,0.795) (8.523,0.795) (8.542,0.795) (8.561,0.795) (8.58,0.795) (8.599,0.795) (8.617,0.795) (8.634,0.795) (8.652,0.796) (8.67,0.796) (8.688,0.796) (8.706,0.796) (8.724,0.796) (8.742,0.796) (8.76,0.796) (8.778,0.796) (8.796,0.796) (8.814,0.796) (8.833,0.796) (8.851,0.796) (8.869,0.796) (8.888,0.796) (8.907,0.796) (8.925,0.796) (8.944,0.796) (8.963,0.796) (8.982,0.796) (9.001,0.796) (9.02,0.796) (9.039,0.796) (9.058,0.796) (9.076,0.796) (9.093,0.796) (9.111,0.796) (9.129,0.796) (9.146,0.796) (9.164,0.796) (9.182,0.796) (9.2,0.796) (9.218,0.797) (9.236,0.797) (9.254,0.797) (9.272,0.797) (9.291,0.797) (9.309,0.797) (9.328,0.797) (9.346,0.797) (9.365,0.797) (9.383,0.797) (9.402,0.797) (9.421,0.797) (9.44,0.797) (9.459,0.797) (9.478,0.797) (9.497,0.797) (9.516,0.797) (9.535,0.797) (9.554,0.797) (9.574,0.797) (9.593,0.797) (9.611,0.797) (9.628,0.797) (9.646,0.797) (9.664,0.797) (9.682,0.797) (9.699,0.797) (9.717,0.797) (9.735,0.797) (9.753,0.797) (9.772,0.797) (9.79,0.797) (9.808,0.797) (9.826,0.797) (9.845,0.797) (9.863,0.797) (9.882,0.797) (9.9,0.797) (9.919,0.797) (9.938,0.798) (9.957,0.798) (9.975,0.798) (9.994,0.798) (10.013,0.798) (10.032,0.798) (10.052,0.798) (10.071,0.798) (10.09,0.798) (10.11,0.798) (10.129,0.798) (10.148,0.798) (10.168,0.798) (10.188,0.798) (10.205,0.798) (10.223,0.798) (10.241,0.798) (10.258,0.798) (10.276,0.798) (10.294,0.798) (10.312,0.798) (10.33,0.798) (10.348,0.798) (10.366,0.798) (10.384,0.798) (10.402,0.798) (10.421,0.798) (10.439,0.798) (10.457,0.798) (10.476,0.798) (10.494,0.798) (10.513,0.798) (10.532,0.798) (10.55,0.798) (10.569,0.798) (10.588,0.798) (10.607,0.798) (10.626,0.798) (10.645,0.798) (10.664,0.798) (10.683,0.798) (10.703,0.798) (10.722,0.798) (10.741,0.798) (10.761,0.798) (10.78,0.798) (10.8,0.798) (10.82,0.798) (10.84,0.798) (10.859,0.798) (10.879,0.798) (10.899,0.798) (10.917,0.798) (10.934,0.798) (10.952,0.798) (10.97,0.798) (10.987,0.798) (11.005,0.798) (11.023,0.798) (11.041,0.799) (11.059,0.799) (11.077,0.799) (11.095,0.799) (11.113,0.799) (11.131,0.799) (11.15,0.799) (11.168,0.799) (11.186,0.799) (11.205,0.799) (11.223,0.799) (11.242,0.799) (11.26,0.799) (11.279,0.799) (11.298,0.799) (11.316,0.799) (11.335,0.799) (11.354,0.799) (11.373,0.799) (11.392,0.799) (11.411,0.799) (11.43,0.799) (11.449,0.799) (11.469,0.799) (11.488,0.799) (11.507,0.799) (11.527,0.799) (11.546,0.799) (11.566,0.799) (11.585,0.799) (11.605,0.799) (11.625,0.799) (11.644,0.799) (11.664,0.799) (11.684,0.799) (11.704,0.799) (11.724,0.799) (11.744,0.799) (11.764,0.799) (11.784,0.799) (11.804,0.799) (11.825,0.799) (11.845,0.799) (11.865,0.799) (11.883,0.799) (11.9,0.799) (11.918,0.799) (11.936,0.799) (11.953,0.799) (11.971,0.799) (11.989,0.799) (12.006,0.799) (12.024,0.799) (12.042,0.799) (12.06,0.799) (12.078,0.799) (12.096,0.799) (12.114,0.799) (12.132,0.799) (12.15,0.799) (12.168,0.799) (12.186,0.799) (12.205,0.799) (12.223,0.799) (12.241,0.799) (12.259,0.799) (12.278,0.799) (12.296,0.799) (12.315,0.799) (12.333,0.799) (12.351,0.799) (12.37,0.799) (12.389,0.799) (12.407,0.799) (12.426,0.799) (12.444,0.799) (12.463,0.799) (12.482,0.799) (12.5,0.799) (12.519,0.799) (12.538,0.799) (12.557,0.799) (12.576,0.799) (12.594,0.799) (12.613,0.799) (12.632,0.799) (12.651,0.799) (12.67,0.799) (12.689,0.799) (12.708,0.799) (12.727,0.799) (12.746,0.799) (12.765,0.799) (12.784,0.799) (12.803,0.799) (12.822,0.799) (12.841,0.799) (12.86,0.799) (12.879,0.799) (12.898,0.799) (12.917,0.799) (12.936,0.799) (12.955,0.799) (12.974,0.799) (12.993,0.799) (13.012,0.799) (13.032,0.799) (13.051,0.799) (13.07,0.799) (13.089,0.799) (13.108,0.799) (13.127,0.799) (13.146,0.799) (13.165,0.799) (13.184,0.799) (13.203,0.799) (13.221,0.799) (13.24,0.799) (13.259,0.799) (13.278,0.799) (13.297,0.799) (13.316,0.799) (13.334,0.799) (13.353,0.799) (13.372,0.799) (13.39,0.8) (13.409,0.8) (13.427,0.8) (13.446,0.8) (13.464,0.8) (13.482,0.8) (13.501,0.8) (13.519,0.8) (13.537,0.8) (13.555,0.8) (13.573,0.8) (13.591,0.8) (13.609,0.8) (13.626,0.8) (13.644,0.8) (13.661,0.8) (13.682,0.8) (13.702,0.8) (13.722,0.8) (13.742,0.8) (13.762,0.8) (13.781,0.8) (13.801,0.8) (13.82,0.8) (13.839,0.8) (13.858,0.8) (13.877,0.8) (13.895,0.8) (13.914,0.8) (13.932,0.8) (13.95,0.8) (13.967,0.8) (13.987,0.8) (14.007,0.8) (14.026,0.8) (14.045,0.8) (14.064,0.8) (14.082,0.8) (14.1,0.8) (14.118,0.8) (14.137,0.8) (14.156,0.8) (14.174,0.8) (14.192,0.8) (14.211,0.8) (14.229,0.8) (14.247,0.8) (14.266,0.8) (14.284,0.8) (14.302,0.8) (14.321,0.8) (14.339,0.8) (14.357,0.8) (14.375,0.8) (14.394,0.8) (14.411,0.8) (14.429,0.8) (14.431,0.8)};
\draw[BeachOrange,dashed,line width=1.1pt] plot coordinates {(0,1.774) (0.007,1.759) (0.014,1.744) (0.021,1.73) (0.028,1.715) (0.049,1.676) (0.067,1.642) (0.084,1.611) (0.102,1.58) (0.121,1.55) (0.139,1.523) (0.157,1.496) (0.174,1.471) (0.192,1.446) (0.21,1.422) (0.228,1.4) (0.246,1.379) (0.264,1.357) (0.283,1.337) (0.3,1.318) (0.318,1.299) (0.337,1.281) (0.354,1.264) (0.372,1.247) (0.39,1.231) (0.409,1.215) (0.427,1.201) (0.445,1.186) (0.463,1.172) (0.481,1.159) (0.499,1.146) (0.517,1.134) (0.535,1.122) (0.553,1.11) (0.572,1.099) (0.589,1.088) (0.607,1.078) (0.626,1.067) (0.644,1.058) (0.663,1.048) (0.681,1.039) (0.698,1.03) (0.717,1.022) (0.735,1.013) (0.754,1.005) (0.773,0.997) (0.791,0.99) (0.809,0.983) (0.827,0.976) (0.846,0.969) (0.865,0.962) (0.884,0.956) (0.902,0.95) (0.92,0.944) (0.938,0.938) (0.957,0.933) (0.976,0.927) (0.995,0.922) (1.013,0.917) (1.031,0.912) (1.049,0.908) (1.067,0.903) (1.086,0.899) (1.105,0.895) (1.125,0.89) (1.142,0.887) (1.16,0.883) (1.178,0.879) (1.197,0.876) (1.216,0.872) (1.235,0.869) (1.254,0.866) (1.274,0.863) (1.292,0.86) (1.31,0.857) (1.328,0.854) (1.347,0.852) (1.366,0.849) (1.385,0.847) (1.405,0.844) (1.425,0.842) (1.442,0.84) (1.46,0.838) (1.478,0.836) (1.497,0.834) (1.516,0.832) (1.535,0.83) (1.555,0.829) (1.575,0.827) (1.596,0.825) (1.614,0.824) (1.632,0.822) (1.65,0.821) (1.669,0.82) (1.688,0.818) (1.708,0.817) (1.728,0.816) (1.749,0.815) (1.77,0.814) (1.792,0.813) (1.81,0.812) (1.828,0.811) (1.847,0.81) (1.867,0.809) (1.887,0.809) (1.907,0.808) (1.929,0.807) (1.951,0.806) (1.974,0.806) (1.991,0.805) (2.01,0.805) (2.022,0.804)};
\draw[BeachTeal,dash dot,line width=1.1pt] plot coordinates {(0,0.276) (0.015,0.275) (0.03,0.276) (0.045,0.277) (0.06,0.278) (0.075,0.279) (0.09,0.281) (0.105,0.282) (0.12,0.284) (0.135,0.286) (0.165,0.29) (0.183,0.292) (0.202,0.295) (0.221,0.298) (0.24,0.301) (0.258,0.303) (0.277,0.306) (0.296,0.309) (0.315,0.312) (0.334,0.315) (0.353,0.319) (0.372,0.322) (0.391,0.325) (0.41,0.328) (0.429,0.332) (0.448,0.335) (0.467,0.338) (0.486,0.341) (0.506,0.345) (0.525,0.348) (0.544,0.352) (0.564,0.355) (0.583,0.358) (0.603,0.362) (0.622,0.365) (0.642,0.369) (0.662,0.372) (0.682,0.376) (0.702,0.379) (0.721,0.383) (0.741,0.386) (0.762,0.39) (0.782,0.393) (0.802,0.397) (0.822,0.4) (0.843,0.404) (0.863,0.407) (0.884,0.411) (0.904,0.414) (0.925,0.417) (0.946,0.421) (0.966,0.424) (0.987,0.428) (1.008,0.431) (1.029,0.435) (1.051,0.438) (1.072,0.442) (1.093,0.445) (1.115,0.449) (1.136,0.452) (1.158,0.456) (1.18,0.459) (1.201,0.462) (1.219,0.465) (1.237,0.468) (1.254,0.471) (1.272,0.473) (1.29,0.476) (1.308,0.479) (1.326,0.481) (1.344,0.484) (1.362,0.487) (1.38,0.489) (1.398,0.492) (1.416,0.495) (1.435,0.497) (1.453,0.5) (1.472,0.503) (1.49,0.505) (1.509,0.508) (1.528,0.511) (1.546,0.513) (1.565,0.516) (1.584,0.518) (1.603,0.521) (1.622,0.523) (1.642,0.526) (1.661,0.529) (1.68,0.531) (1.7,0.534) (1.719,0.536) (1.739,0.539) (1.759,0.541) (1.778,0.544) (1.798,0.546) (1.818,0.549) (1.838,0.551) (1.859,0.554) (1.879,0.556) (1.899,0.559) (1.92,0.561) (1.94,0.563) (1.961,0.566) (1.982,0.568) (2.003,0.571) (2.024,0.573) (2.045,0.575) (2.066,0.578) (2.087,0.58) (2.109,0.582) (2.13,0.585) (2.152,0.587) (2.174,0.589) (2.195,0.592) (2.217,0.594) (2.24,0.596) (2.262,0.599) (2.284,0.601) (2.306,0.603) (2.329,0.605) (2.352,0.608) (2.375,0.61) (2.398,0.612) (2.421,0.614) (2.444,0.616) (2.467,0.619) (2.485,0.62) (2.502,0.622) (2.52,0.624) (2.538,0.625) (2.556,0.627) (2.574,0.628) (2.592,0.63) (2.61,0.631) (2.628,0.633) (2.647,0.635) (2.665,0.636) (2.684,0.638) (2.702,0.639) (2.721,0.641) (2.74,0.642) (2.758,0.644) (2.777,0.645) (2.796,0.647) (2.815,0.648) (2.835,0.65) (2.854,0.651) (2.873,0.653) (2.893,0.654) (2.912,0.656) (2.932,0.657) (2.952,0.659) (2.972,0.66) (2.992,0.662) (3.012,0.663) (3.032,0.665) (3.052,0.666) (3.072,0.667) (3.093,0.669) (3.114,0.67) (3.134,0.672) (3.155,0.673) (3.176,0.674) (3.197,0.676) (3.218,0.677) (3.239,0.679) (3.261,0.68) (3.282,0.681) (3.304,0.683) (3.326,0.684) (3.348,0.685) (3.37,0.687) (3.392,0.688) (3.414,0.689) (3.436,0.691) (3.459,0.692) (3.481,0.693) (3.504,0.694) (3.527,0.696) (3.55,0.697) (3.573,0.698) (3.597,0.699) (3.62,0.701) (3.644,0.702) (3.667,0.703) (3.691,0.704) (3.715,0.706) (3.739,0.707) (3.764,0.708) (3.788,0.709) (3.813,0.71) (3.838,0.712) (3.863,0.713) (3.888,0.714) (3.913,0.715) (3.939,0.716) (3.964,0.717) (3.99,0.718) (4.016,0.72) (4.042,0.721) (4.06,0.721) (4.078,0.722) (4.095,0.723) (4.113,0.724) (4.131,0.724) (4.149,0.725) (4.167,0.726) (4.185,0.727) (4.204,0.727) (4.222,0.728) (4.24,0.729) (4.259,0.729) (4.277,0.73) (4.296,0.731) (4.315,0.731) (4.334,0.732) (4.353,0.733) (4.372,0.733) (4.391,0.734) (4.41,0.735) (4.43,0.736) (4.449,0.736) (4.469,0.737) (4.488,0.738) (4.508,0.738) (4.528,0.739) (4.548,0.74) (4.568,0.74) (4.588,0.741) (4.608,0.741) (4.629,0.742) (4.649,0.743) (4.67,0.743) (4.691,0.744) (4.712,0.745) (4.733,0.745) (4.754,0.746) (4.775,0.746) (4.796,0.747) (4.818,0.748) (4.839,0.748) (4.861,0.749) (4.883,0.75) (4.905,0.75) (4.927,0.751) (4.949,0.751) (4.971,0.752) (4.994,0.752) (5.017,0.753) (5.039,0.754) (5.062,0.754) (5.085,0.755) (5.109,0.755) (5.132,0.756) (5.155,0.756) (5.179,0.757) (5.203,0.758) (5.227,0.758) (5.251,0.759) (5.275,0.759) (5.3,0.76) (5.324,0.76) (5.349,0.761) (5.374,0.761) (5.399,0.762) (5.424,0.762) (5.449,0.763) (5.475,0.763) (5.501,0.764) (5.527,0.764) (5.553,0.765) (5.579,0.765) (5.606,0.766) (5.632,0.766) (5.659,0.767) (5.686,0.767) (5.714,0.768) (5.741,0.768) (5.769,0.769) (5.797,0.769) (5.825,0.77) (5.853,0.77) (5.882,0.771) (5.911,0.771) (5.94,0.772) (5.969,0.772) (5.998,0.772) (6.028,0.773) (6.058,0.773) (6.088,0.774) (6.119,0.774) (6.149,0.775) (6.18,0.775) (6.212,0.775) (6.243,0.776) (6.275,0.776) (6.307,0.777) (6.339,0.777) (6.372,0.778) (6.405,0.778) (6.438,0.778) (6.471,0.779) (6.505,0.779) (6.539,0.779) (6.574,0.78) (6.608,0.78) (6.626,0.78) (6.643,0.781) (6.661,0.781) (6.679,0.781) (6.697,0.781) (6.715,0.781) (6.733,0.782) (6.751,0.782) (6.769,0.782) (6.787,0.782) (6.806,0.782) (6.824,0.782) (6.843,0.783) (6.861,0.783) (6.88,0.783) (6.899,0.783) (6.918,0.783) (6.937,0.783) (6.956,0.784) (6.975,0.784) (6.995,0.784) (7.014,0.784) (7.034,0.784) (7.053,0.785) (7.073,0.785) (7.093,0.785) (7.113,0.785) (7.133,0.785) (7.153,0.785) (7.174,0.785) (7.194,0.786) (7.214,0.786) (7.235,0.786) (7.256,0.786) (7.277,0.786) (7.298,0.786) (7.319,0.787) (7.34,0.787) (7.362,0.787) (7.383,0.787) (7.405,0.787) (7.426,0.787) (7.448,0.788) (7.47,0.788) (7.493,0.788) (7.515,0.788) (7.537,0.788) (7.56,0.788) (7.583,0.788) (7.605,0.789) (7.628,0.789) (7.652,0.789) (7.675,0.789) (7.698,0.789) (7.722,0.789) (7.746,0.789) (7.769,0.79) (7.793,0.79) (7.818,0.79) (7.842,0.79) (7.867,0.79) (7.891,0.79) (7.916,0.79) (7.941,0.79) (7.966,0.791) (7.992,0.791) (8.017,0.791) (8.043,0.791) (8.069,0.791) (8.095,0.791) (8.121,0.791) (8.148,0.791) (8.175,0.792) (8.201,0.792) (8.229,0.792) (8.256,0.792) (8.283,0.792) (8.311,0.792) (8.339,0.792) (8.367,0.792) (8.395,0.793) (8.424,0.793) (8.453,0.793) (8.482,0.793) (8.511,0.793) (8.54,0.793) (8.57,0.793) (8.6,0.793) (8.63,0.793) (8.661,0.794) (8.692,0.794) (8.723,0.794) (8.754,0.794) (8.785,0.794) (8.817,0.794) (8.849,0.794) (8.882,0.794) (8.914,0.794) (8.947,0.795) (8.98,0.795) (9.014,0.795) (9.048,0.795) (9.082,0.795) (9.117,0.795) (9.151,0.795) (9.187,0.795) (9.222,0.795) (9.258,0.795) (9.294,0.795) (9.331,0.796) (9.368,0.796) (9.405,0.796) (9.443,0.796) (9.481,0.796) (9.519,0.796) (9.558,0.796) (9.598,0.796) (9.637,0.796) (9.678,0.796) (9.718,0.796) (9.759,0.796) (9.801,0.797) (9.843,0.797) (9.886,0.797) (9.929,0.797) (9.972,0.797) (10.016,0.797) (10.061,0.797) (10.106,0.797) (10.152,0.797) (10.198,0.797) (10.245,0.797) (10.293,0.797) (10.341,0.797) (10.39,0.798) (10.439,0.798) (10.49,0.798) (10.541,0.798) (10.592,0.798) (10.645,0.798) (10.698,0.798) (10.752,0.798) (10.806,0.798) (10.862,0.798) (10.918,0.798) (10.976,0.798) (11.034,0.798) (11.093,0.798) (11.153,0.798) (11.214,0.798) (11.276,0.798) (11.34,0.799) (11.404,0.799) (11.47,0.799) (11.536,0.799) (11.604,0.799) (11.674,0.799) (11.744,0.799) (11.816,0.799) (11.89,0.799) (11.965,0.799) (12.042,0.799) (12.12,0.799) (12.2,0.799) (12.282,0.799) (12.365,0.799) (12.451,0.799) (12.539,0.799) (12.628,0.799) (12.721,0.799) (12.815,0.799) (12.912,0.799) (13.012,0.799) (13.114,0.799) (13.22,0.799) (13.328,0.8) (13.44,0.8) (13.556,0.8) (13.675,0.8) (13.799,0.8) (13.926,0.8)};
\end{scope}
\draw[BeachBlue,solid,line width=1.1pt] (4.6,7.55)--(5.2,7.55);
\node[right,font=\scriptsize\sffamily] at (5.2,7.55) {Type A：$\alpha=60^\circ$、$\phi_0=3.84$ V};
\draw[BeachOrange,dashed,line width=1.1pt] (4.6,7.1)--(5.2,7.1);
\node[right,font=\scriptsize\sffamily] at (5.2,7.1) {Type B：$\alpha=20^\circ$、$\phi_0=1.62$ V};
\draw[BeachTeal,dash dot,line width=1.1pt] (4.6,6.65)--(5.2,6.65);
\node[right,font=\scriptsize\sffamily] at (5.2,6.65) {Type C：$\alpha=10^\circ$、$\phi_0=-4.24$ V};
\end{tikzpicture}
\caption[Type A、B、C の電位と電場の数値解]{零電流条件で解いた Type A、B、C の電位と電場。条件は表\ref{tab:sm_example_inputs}で、型ごとに太陽高度 $\alpha$ を変えた。Type A は $z_m\approx14.9$ m に負の極小 $\phi_m=-0.332$ V を持ち、極小の手前で電場の向きが反転する。Type B は Debye 長 $\lambda_{D,\mathrm{pe}}\approx1.4$ m 程度で減衰し、Type C は太陽風電子の Debye 長 $\lambda_{D,e}\approx8.7$ m 程度でゆっくり 0 V に近づく。電子 drift は 0 とした。曲線は $|\phi|$ が 2.2 mV（$10^{-3}T_{\mathrm{pe}}$）に達した点で止め、その先に 0 V の線を補っていない。}
\label{fig:sm_profiles}
\end{figure}


\subsection{粒子集団と速度分布}
図\ref{fig:sm_densities}は、Type A の粒子集団ごとの密度である。壁の近くでは捕捉された光電子が約 50 cm$^{-3}$ あり、
イオン密度の約 6 倍の負電荷が壁の正電位を遮蔽する。捕捉光電子は極小で 0 になり、その上流側は反射した太陽風電子に置き換わる。
上流では通過電子 4.0 cm$^{-3}$、反射電子 1.8 cm$^{-3}$、脱出した光電子 2.8 cm$^{-3}$ の和がイオン密度 8.7 cm$^{-3}$ と釣り合う。
このとき規格化 $N_e=9.9$ cm$^{-3}$ は実際の上流の太陽風電子密度 5.9 cm$^{-3}$ と異なる。
壁への数流束は太陽風電子 $5.6\times10^{12}$、イオン $3.5\times10^{12}$、脱出光電子 $2.1\times10^{12}\,\mathrm{m^{-2}\,s^{-1}}$ で、
$\Gamma_e=\Gamma_i+\Gamma_{\mathrm{esc}}$ が成り立つ。放出した光電子の 85\% は障壁で戻され、脱出は 15\% である。
Type B（$\alpha=20^\circ$）では脱出が 48\%、Type C では障壁がないので 100\% になる。

% Generated by tools/sheath_model_figures.py from profile_A.csv; edit the generator, not this file.
\begin{figure}[htbp]
\centering
\begin{tikzpicture}[beach diagram]
\draw[black!10,line width=0.4pt] (0,0)--(0,6);
\draw[black!10,line width=0.4pt] (1.75,0)--(1.75,6);
\draw[black!10,line width=0.4pt] (3.5,0)--(3.5,6);
\draw[black!10,line width=0.4pt] (5.25,0)--(5.25,6);
\draw[black!10,line width=0.4pt] (7,0)--(7,6);
\draw[black!10,line width=0.4pt] (8.75,0)--(8.75,6);
\draw[black!10,line width=0.4pt] (10.5,0)--(10.5,6);
\draw[black!10,line width=0.4pt] (0,0.891)--(10.5,0.891);
\draw[black!10,line width=0.4pt] (0,2.594)--(10.5,2.594);
\draw[black!10,line width=0.4pt] (0,4.297)--(10.5,4.297);
\draw[black!10,line width=0.4pt] (0,6)--(10.5,6);
\draw[line width=0.6pt] (0,0) rectangle (10.5,6);
\draw[line width=0.6pt] (0,0)--++(0,0.1);
\node[below,font=\scriptsize\sffamily] at (0,0) {$0$};
\draw[line width=0.6pt] (1.75,0)--++(0,0.1);
\node[below,font=\scriptsize\sffamily] at (1.75,0) {$10$};
\draw[line width=0.6pt] (3.5,0)--++(0,0.1);
\node[below,font=\scriptsize\sffamily] at (3.5,0) {$20$};
\draw[line width=0.6pt] (5.25,0)--++(0,0.1);
\node[below,font=\scriptsize\sffamily] at (5.25,0) {$30$};
\draw[line width=0.6pt] (7,0)--++(0,0.1);
\node[below,font=\scriptsize\sffamily] at (7,0) {$40$};
\draw[line width=0.6pt] (8.75,0)--++(0,0.1);
\node[below,font=\scriptsize\sffamily] at (8.75,0) {$50$};
\draw[line width=0.6pt] (10.5,0)--++(0,0.1);
\node[below,font=\scriptsize\sffamily] at (10.5,0) {$60$};
\draw[line width=0.6pt] (0,0.891)--++(0.1,0);
\node[left,font=\scriptsize\sffamily] at (0,0.891) {$0.1$};
\draw[line width=0.6pt] (0,2.594)--++(0.1,0);
\node[left,font=\scriptsize\sffamily] at (0,2.594) {$1$};
\draw[line width=0.6pt] (0,4.297)--++(0.1,0);
\node[left,font=\scriptsize\sffamily] at (0,4.297) {$10$};
\draw[line width=0.6pt] (0,6)--++(0.1,0);
\node[left,font=\scriptsize\sffamily] at (0,6) {$100$};
\node[font=\small\sffamily] at (5.25,-0.75) {壁からの距離 $z$ [m]};
\node[rotate=90,font=\small\sffamily] at (-1.05,3) {数密度 [cm$^{-3}$]};
\begin{scope}
\clip (0,0) rectangle (10.5,6);
\draw[BeachInk,dotted,line width=0.6pt] (2.603,0)--(2.603,6);
\draw[BeachInk,line width=1.2pt] plot coordinates {(0,4.195) (0.027,4.195) (0.044,4.195) (0.062,4.195) (0.079,4.195) (0.097,4.195) (0.114,4.195) (0.132,4.195) (0.15,4.195) (0.167,4.195) (0.185,4.195) (0.203,4.195) (0.22,4.195) (0.238,4.195) (0.256,4.195) (0.273,4.195) (0.291,4.195) (0.309,4.195) (0.326,4.195) (0.344,4.195) (0.361,4.195) (0.379,4.195) (0.397,4.195) (0.415,4.194) (0.432,4.194) (0.45,4.194) (0.468,4.194) (0.486,4.194) (0.503,4.194) (0.521,4.194) (0.538,4.194) (0.556,4.194) (0.574,4.194) (0.592,4.194) (0.61,4.194) (0.627,4.194) (0.645,4.194) (0.663,4.194) (0.68,4.194) (0.698,4.194) (0.716,4.194) (0.734,4.194) (0.751,4.194) (0.769,4.194) (0.787,4.194) (0.805,4.194) (0.822,4.194) (0.84,4.194) (0.858,4.194) (0.876,4.194) (0.894,4.194) (0.912,4.194) (0.93,4.194) (0.948,4.194) (0.965,4.194) (0.983,4.194) (1.001,4.194) (1.019,4.194) (1.036,4.194) (1.054,4.194) (1.072,4.194) (1.09,4.194) (1.108,4.194) (1.126,4.194) (1.143,4.194) (1.161,4.194) (1.179,4.194) (1.197,4.194) (1.215,4.194) (1.233,4.194) (1.251,4.194) (1.268,4.194) (1.286,4.194) (1.304,4.194) (1.322,4.194) (1.34,4.194) (1.357,4.194) (1.375,4.194) (1.393,4.194) (1.412,4.194) (1.429,4.194) (1.447,4.194) (1.465,4.194) (1.484,4.194) (1.501,4.194) (1.519,4.194) (1.537,4.194) (1.555,4.194) (1.574,4.194) (1.591,4.194) (1.609,4.194) (1.627,4.194) (1.645,4.194) (1.664,4.194) (1.681,4.194) (1.699,4.194) (1.717,4.194) (1.735,4.194) (1.754,4.194) (1.772,4.194) (1.79,4.194) (1.808,4.194) (1.826,4.194) (1.844,4.194) (1.863,4.194) (1.881,4.194) (1.899,4.194) (1.917,4.194) (1.935,4.194) (1.953,4.194) (1.972,4.194) (1.99,4.194) (2.008,4.194) (2.026,4.194) (2.044,4.194) (2.062,4.194) (2.08,4.194) (2.099,4.194) (2.117,4.194) (2.135,4.194) (2.153,4.194) (2.171,4.194) (2.189,4.194) (2.207,4.194) (2.225,4.194) (2.244,4.194) (2.263,4.194) (2.281,4.194) (2.298,4.194) (2.316,4.194) (2.334,4.194) (2.353,4.194) (2.371,4.194) (2.39,4.194) (2.409,4.194) (2.428,4.194) (2.445,4.194) (2.463,4.194) (2.481,4.194) (2.499,4.194) (2.517,4.194) (2.535,4.194) (2.553,4.194) (2.572,4.194) (2.619,4.194) (2.637,4.194) (2.659,4.194) (2.681,4.194) (2.699,4.194) (2.719,4.194) (2.74,4.194) (2.758,4.194) (2.778,4.194) (2.797,4.194) (2.816,4.194) (2.834,4.194) (2.854,4.194) (2.872,4.194) (2.891,4.194) (2.91,4.194) (2.928,4.194) (2.946,4.194) (2.964,4.194) (2.984,4.194) (3.002,4.194) (3.02,4.194) (3.039,4.194) (3.057,4.194) (3.075,4.194) (3.094,4.194) (3.112,4.194) (3.13,4.194) (3.148,4.194) (3.166,4.194) (3.184,4.194) (3.202,4.194) (3.221,4.194) (3.239,4.194) (3.257,4.194) (3.275,4.194) (3.293,4.194) (3.311,4.194) (3.329,4.194) (3.347,4.194) (3.365,4.194) (3.383,4.194) (3.401,4.194) (3.419,4.194) (3.437,4.194) (3.455,4.194) (3.473,4.194) (3.49,4.194) (3.509,4.194) (3.527,4.194) (3.545,4.194) (3.563,4.194) (3.581,4.194) (3.598,4.194) (3.616,4.194) (3.634,4.194) (3.652,4.194) (3.671,4.194) (3.689,4.194) (3.706,4.194) (3.724,4.194) (3.742,4.194) (3.76,4.194) (3.777,4.194) (3.796,4.194) (3.814,4.194) (3.832,4.194) (3.85,4.194) (3.868,4.194) (3.886,4.194) (3.904,4.194) (3.922,4.194) (3.94,4.194) (3.958,4.194) (3.975,4.194) (3.993,4.194) (4.011,4.194) (4.028,4.194) (4.046,4.194) (4.064,4.194) (4.082,4.194) (4.101,4.194) (4.119,4.194) (4.137,4.194) (4.155,4.194) (4.173,4.194) (4.192,4.194) (4.21,4.194) (4.228,4.194) (4.246,4.194) (4.264,4.194) (4.282,4.194) (4.3,4.194) (4.318,4.194) (4.336,4.194) (4.354,4.194) (4.371,4.194) (4.389,4.194) (4.407,4.194) (4.425,4.194) (4.443,4.194) (4.461,4.194) (4.479,4.194) (4.497,4.194) (4.515,4.194) (4.532,4.194) (4.55,4.194) (4.568,4.194) (4.586,4.194) (4.604,4.194) (4.622,4.194) (4.64,4.194) (4.658,4.194) (4.676,4.194) (4.693,4.194) (4.711,4.194) (4.729,4.194) (4.747,4.194) (4.765,4.194) (4.783,4.194) (4.801,4.194) (4.819,4.194) (4.837,4.194) (4.855,4.194) (4.873,4.194) (4.891,4.194) (4.91,4.194) (4.928,4.194) (4.946,4.194) (4.964,4.194) (4.982,4.194) (5,4.194) (5.019,4.194) (5.037,4.194) (5.055,4.194) (5.074,4.194) (5.092,4.194) (5.109,4.194) (5.127,4.194) (5.145,4.194) (5.162,4.194) (5.18,4.194) (5.197,4.194) (5.215,4.194) (5.233,4.194) (5.251,4.194) (5.268,4.194) (5.286,4.194) (5.304,4.194) (5.322,4.194) (5.34,4.194) (5.358,4.194) (5.376,4.194) (5.394,4.194) (5.412,4.194) (5.43,4.194) (5.448,4.194) (5.466,4.194) (5.484,4.194) (5.503,4.194) (5.521,4.194) (5.539,4.194) (5.558,4.194) (5.576,4.194) (5.594,4.194) (5.611,4.194) (5.629,4.194) (5.647,4.194) (5.664,4.194) (5.682,4.194) (5.7,4.194) (5.718,4.194) (5.735,4.194) (5.753,4.194) (5.771,4.194) (5.789,4.194) (5.807,4.194) (5.825,4.194) (5.844,4.194) (5.862,4.194) (5.88,4.194) (5.898,4.194) (5.917,4.194) (5.935,4.194) (5.954,4.194) (5.971,4.194) (5.989,4.194) (6.006,4.194) (6.024,4.194) (6.042,4.194) (6.059,4.194) (6.077,4.194) (6.095,4.194) (6.113,4.194) (6.131,4.194) (6.149,4.194) (6.167,4.194) (6.185,4.194) (6.203,4.194) (6.221,4.194) (6.24,4.194) (6.258,4.194) (6.276,4.194) (6.295,4.194) (6.313,4.194) (6.331,4.194) (6.348,4.194) (6.366,4.194) (6.384,4.194) (6.402,4.194) (6.419,4.194) (6.437,4.194) (6.455,4.194) (6.473,4.194) (6.491,4.194) (6.509,4.194) (6.527,4.194) (6.545,4.194) (6.564,4.194) (6.582,4.194) (6.6,4.194) (6.619,4.194) (6.637,4.194) (6.656,4.194) (6.674,4.194) (6.692,4.194) (6.709,4.194) (6.727,4.194) (6.745,4.194) (6.762,4.194) (6.78,4.194) (6.798,4.194) (6.816,4.194) (6.834,4.194) (6.852,4.194) (6.87,4.194) (6.889,4.194) (6.907,4.194) (6.925,4.194) (6.943,4.194) (6.962,4.194) (6.98,4.194) (6.999,4.194) (7.018,4.194) (7.035,4.194) (7.053,4.194) (7.07,4.194) (7.088,4.194) (7.106,4.194) (7.124,4.194) (7.141,4.194) (7.159,4.194) (7.177,4.194) (7.195,4.194) (7.213,4.194) (7.232,4.194) (7.25,4.194) (7.268,4.194) (7.287,4.194) (7.305,4.194) (7.324,4.194) (7.342,4.194) (7.361,4.194) (7.379,4.194) (7.397,4.194) (7.414,4.194) (7.432,4.194) (7.45,4.194) (7.468,4.194) (7.485,4.194) (7.503,4.194) (7.521,4.194) (7.539,4.194) (7.557,4.194) (7.575,4.194) (7.594,4.194) (7.612,4.194) (7.63,4.194) (7.649,4.194) (7.667,4.194) (7.686,4.194) (7.704,4.194) (7.723,4.194) (7.742,4.194) (7.76,4.194) (7.778,4.194) (7.795,4.194) (7.813,4.194) (7.831,4.194) (7.849,4.194) (7.866,4.194) (7.884,4.194) (7.902,4.194) (7.92,4.194) (7.938,4.194) (7.957,4.194) (7.975,4.194) (7.993,4.194) (8.011,4.194) (8.03,4.194) (8.048,4.194) (8.067,4.194) (8.085,4.194) (8.104,4.194) (8.123,4.194) (8.142,4.194) (8.161,4.194) (8.178,4.194) (8.196,4.194) (8.213,4.194) (8.231,4.194) (8.249,4.194) (8.267,4.194) (8.285,4.194) (8.303,4.194) (8.321,4.194) (8.339,4.194) (8.357,4.194) (8.375,4.194) (8.393,4.194) (8.412,4.194) (8.43,4.194) (8.449,4.194) (8.467,4.194) (8.486,4.194) (8.505,4.194) (8.523,4.194) (8.542,4.194) (8.561,4.194) (8.58,4.194) (8.599,4.194) (8.617,4.194) (8.634,4.194) (8.652,4.194) (8.67,4.194) (8.688,4.194) (8.706,4.194) (8.724,4.194) (8.742,4.194) (8.76,4.194) (8.778,4.194) (8.796,4.194) (8.814,4.194) (8.833,4.194) (8.851,4.194) (8.869,4.194) (8.888,4.194) (8.907,4.194) (8.925,4.194) (8.944,4.194) (8.963,4.194) (8.982,4.194) (9.001,4.194) (9.02,4.194) (9.039,4.194) (9.058,4.194) (9.076,4.194) (9.093,4.194) (9.111,4.194) (9.129,4.194) (9.146,4.194) (9.164,4.194) (9.182,4.194) (9.2,4.194) (9.218,4.194) (9.236,4.194) (9.254,4.194) (9.272,4.194) (9.291,4.194) (9.309,4.194) (9.328,4.194) (9.346,4.194) (9.365,4.194) (9.383,4.194) (9.402,4.194) (9.421,4.194) (9.44,4.194) (9.459,4.194) (9.478,4.194) (9.497,4.194) (9.516,4.194) (9.535,4.194) (9.554,4.194) (9.574,4.194) (9.593,4.194) (9.611,4.194) (9.628,4.194) (9.646,4.194) (9.664,4.194) (9.682,4.194) (9.699,4.194) (9.717,4.194) (9.735,4.194) (9.753,4.194) (9.772,4.194) (9.79,4.194) (9.808,4.194) (9.826,4.194) (9.845,4.194) (9.863,4.194) (9.882,4.194) (9.9,4.194) (9.919,4.194) (9.938,4.194) (9.957,4.194) (9.975,4.194) (9.994,4.194) (10.013,4.194) (10.032,4.194) (10.052,4.194) (10.071,4.194) (10.09,4.194) (10.11,4.194) (10.129,4.194) (10.148,4.194) (10.168,4.194) (10.188,4.194) (10.205,4.194) (10.223,4.194) (10.241,4.194) (10.258,4.194) (10.276,4.194) (10.294,4.194) (10.312,4.194) (10.33,4.194) (10.348,4.194) (10.366,4.194) (10.384,4.194) (10.402,4.194) (10.421,4.194) (10.439,4.194) (10.457,4.194) (10.476,4.194) (10.494,4.194) (10.513,4.194) (10.532,4.194) (10.55,4.194) (10.569,4.194) (10.588,4.194) (10.607,4.194) (10.626,4.194) (10.645,4.194) (10.664,4.194) (10.683,4.194) (10.703,4.194) (10.722,4.194) (10.741,4.194) (10.761,4.194) (10.78,4.194) (10.8,4.194) (10.82,4.194) (10.84,4.194) (10.859,4.194) (10.879,4.194) (10.899,4.194) (10.917,4.194) (10.934,4.194) (10.952,4.194) (10.97,4.194) (10.987,4.194) (11.005,4.194) (11.023,4.194) (11.041,4.194) (11.059,4.194) (11.077,4.194) (11.095,4.194) (11.113,4.194) (11.131,4.194) (11.15,4.194) (11.168,4.194) (11.186,4.194) (11.205,4.194) (11.223,4.194) (11.242,4.194) (11.26,4.194) (11.279,4.194) (11.298,4.194) (11.316,4.194) (11.335,4.194) (11.354,4.194) (11.373,4.194) (11.392,4.194) (11.411,4.194) (11.43,4.194) (11.449,4.194) (11.469,4.194) (11.488,4.194) (11.507,4.194) (11.527,4.194) (11.546,4.194) (11.566,4.194) (11.585,4.194) (11.605,4.194) (11.625,4.194) (11.644,4.194) (11.664,4.194) (11.684,4.194) (11.704,4.194) (11.724,4.194) (11.744,4.194) (11.764,4.194) (11.784,4.194) (11.804,4.194) (11.825,4.194) (11.845,4.194) (11.865,4.194) (11.883,4.194) (11.9,4.194) (11.918,4.194) (11.936,4.194) (11.953,4.194) (11.971,4.194) (11.989,4.194) (12.006,4.194) (12.024,4.194) (12.042,4.194) (12.06,4.194) (12.078,4.194) (12.096,4.194) (12.114,4.194) (12.132,4.194) (12.15,4.194) (12.168,4.194) (12.186,4.194) (12.205,4.194) (12.223,4.194) (12.241,4.194) (12.259,4.194) (12.278,4.194) (12.296,4.194) (12.315,4.194) (12.333,4.194) (12.351,4.194) (12.37,4.194) (12.389,4.194) (12.407,4.194) (12.426,4.194) (12.444,4.194) (12.463,4.194) (12.482,4.194) (12.5,4.194) (12.519,4.194) (12.538,4.194) (12.557,4.194) (12.576,4.194) (12.594,4.194) (12.613,4.194) (12.632,4.194) (12.651,4.194) (12.67,4.194) (12.689,4.194) (12.708,4.194) (12.727,4.194) (12.746,4.194) (12.765,4.194) (12.784,4.194) (12.803,4.194) (12.822,4.194) (12.841,4.194) (12.86,4.194) (12.879,4.194) (12.898,4.194) (12.917,4.194) (12.936,4.194) (12.955,4.194) (12.974,4.194) (12.993,4.194) (13.012,4.194) (13.032,4.194) (13.051,4.194) (13.07,4.194) (13.089,4.194) (13.108,4.194) (13.127,4.194) (13.146,4.194) (13.165,4.194) (13.184,4.194) (13.203,4.194) (13.221,4.194) (13.24,4.194) (13.259,4.194) (13.278,4.194) (13.297,4.194) (13.316,4.194) (13.334,4.194) (13.353,4.194) (13.372,4.194) (13.39,4.194) (13.409,4.194) (13.427,4.194) (13.446,4.194) (13.464,4.194) (13.482,4.194) (13.501,4.194) (13.519,4.194) (13.537,4.194) (13.555,4.194) (13.573,4.194) (13.591,4.194) (13.609,4.194) (13.626,4.194) (13.644,4.194) (13.661,4.194) (13.682,4.194) (13.702,4.194) (13.722,4.194) (13.742,4.194) (13.762,4.194) (13.781,4.194) (13.801,4.194) (13.82,4.194) (13.839,4.194) (13.858,4.194) (13.877,4.194) (13.895,4.194) (13.914,4.194) (13.932,4.194) (13.95,4.194) (13.967,4.194) (13.987,4.194) (14.007,4.194) (14.026,4.194) (14.045,4.194) (14.064,4.194) (14.082,4.194) (14.1,4.194) (14.118,4.194) (14.137,4.194) (14.156,4.194) (14.174,4.194) (14.192,4.194) (14.211,4.194) (14.229,4.194) (14.247,4.194) (14.266,4.194) (14.284,4.194) (14.302,4.194) (14.321,4.194) (14.339,4.194) (14.357,4.194) (14.375,4.194) (14.394,4.194) (14.411,4.194) (14.429,4.194) (14.431,4.194)};
\draw[BeachBlue,line width=1.1pt] plot coordinates {(0,3.345) (0.027,3.355) (0.044,3.361) (0.062,3.368) (0.079,3.374) (0.097,3.38) (0.114,3.385) (0.132,3.391) (0.15,3.397) (0.167,3.402) (0.185,3.408) (0.203,3.413) (0.22,3.418) (0.238,3.423) (0.256,3.428) (0.273,3.433) (0.291,3.438) (0.309,3.443) (0.326,3.447) (0.344,3.452) (0.361,3.456) (0.379,3.461) (0.397,3.465) (0.415,3.469) (0.432,3.474) (0.45,3.478) (0.468,3.482) (0.486,3.486) (0.503,3.49) (0.521,3.494) (0.538,3.498) (0.556,3.502) (0.574,3.506) (0.592,3.509) (0.61,3.513) (0.627,3.517) (0.645,3.52) (0.663,3.524) (0.68,3.528) (0.698,3.531) (0.716,3.534) (0.734,3.538) (0.751,3.541) (0.769,3.545) (0.787,3.548) (0.805,3.551) (0.822,3.554) (0.84,3.558) (0.858,3.561) (0.876,3.564) (0.894,3.567) (0.912,3.57) (0.93,3.573) (0.948,3.576) (0.965,3.579) (0.983,3.582) (1.001,3.585) (1.019,3.588) (1.036,3.59) (1.054,3.593) (1.072,3.596) (1.09,3.599) (1.108,3.601) (1.126,3.604) (1.143,3.607) (1.161,3.609) (1.179,3.612) (1.197,3.615) (1.215,3.617) (1.233,3.62) (1.251,3.622) (1.268,3.625) (1.286,3.627) (1.304,3.63) (1.322,3.632) (1.34,3.634) (1.357,3.637) (1.375,3.639) (1.393,3.641) (1.412,3.644) (1.429,3.646) (1.447,3.648) (1.465,3.651) (1.484,3.653) (1.501,3.655) (1.519,3.657) (1.537,3.659) (1.555,3.661) (1.574,3.664) (1.591,3.666) (1.609,3.668) (1.627,3.67) (1.645,3.672) (1.664,3.674) (1.681,3.676) (1.699,3.678) (1.717,3.68) (1.735,3.682) (1.754,3.684) (1.772,3.686) (1.79,3.688) (1.808,3.689) (1.826,3.691) (1.844,3.693) (1.863,3.695) (1.881,3.697) (1.899,3.699) (1.917,3.7) (1.935,3.702) (1.953,3.704) (1.972,3.706) (1.99,3.708) (2.008,3.709) (2.026,3.711) (2.044,3.713) (2.062,3.714) (2.08,3.716) (2.099,3.718) (2.117,3.719) (2.135,3.721) (2.153,3.722) (2.171,3.724) (2.189,3.725) (2.207,3.727) (2.225,3.729) (2.244,3.73) (2.263,3.732) (2.281,3.733) (2.298,3.735) (2.316,3.736) (2.334,3.737) (2.353,3.739) (2.371,3.74) (2.39,3.742) (2.409,3.743) (2.428,3.745) (2.445,3.746) (2.463,3.747) (2.481,3.749) (2.499,3.75) (2.517,3.751) (2.535,3.753) (2.553,3.754) (2.572,3.755) (2.619,3.756) (2.637,3.755) (2.659,3.754) (2.681,3.752) (2.699,3.751) (2.719,3.75) (2.74,3.748) (2.758,3.747) (2.778,3.746) (2.797,3.745) (2.816,3.743) (2.834,3.742) (2.854,3.741) (2.872,3.74) (2.891,3.739) (2.91,3.738) (2.928,3.737) (2.946,3.736) (2.964,3.735) (2.984,3.734) (3.002,3.733) (3.02,3.732) (3.039,3.731) (3.057,3.73) (3.075,3.729) (3.094,3.728) (3.112,3.727) (3.13,3.726) (3.148,3.725) (3.166,3.724) (3.184,3.723) (3.202,3.722) (3.221,3.721) (3.239,3.72) (3.257,3.719) (3.275,3.719) (3.293,3.718) (3.311,3.717) (3.329,3.716) (3.347,3.715) (3.365,3.714) (3.383,3.714) (3.401,3.713) (3.419,3.712) (3.437,3.711) (3.455,3.71) (3.473,3.71) (3.49,3.709) (3.509,3.708) (3.527,3.707) (3.545,3.707) (3.563,3.706) (3.581,3.705) (3.598,3.704) (3.616,3.704) (3.634,3.703) (3.652,3.702) (3.671,3.702) (3.689,3.701) (3.706,3.7) (3.724,3.7) (3.742,3.699) (3.76,3.698) (3.777,3.698) (3.796,3.697) (3.814,3.696) (3.832,3.696) (3.85,3.695) (3.868,3.694) (3.886,3.694) (3.904,3.693) (3.922,3.693) (3.94,3.692) (3.958,3.691) (3.975,3.691) (3.993,3.69) (4.011,3.69) (4.028,3.689) (4.046,3.689) (4.064,3.688) (4.082,3.687) (4.101,3.687) (4.119,3.686) (4.137,3.686) (4.155,3.685) (4.173,3.685) (4.192,3.684) (4.21,3.684) (4.228,3.683) (4.246,3.683) (4.264,3.682) (4.282,3.682) (4.3,3.681) (4.318,3.681) (4.336,3.68) (4.354,3.68) (4.371,3.679) (4.389,3.679) (4.407,3.678) (4.425,3.678) (4.443,3.677) (4.461,3.677) (4.479,3.676) (4.497,3.676) (4.515,3.676) (4.532,3.675) (4.55,3.675) (4.568,3.674) (4.586,3.674) (4.604,3.673) (4.622,3.673) (4.64,3.673) (4.658,3.672) (4.676,3.672) (4.693,3.671) (4.711,3.671) (4.729,3.671) (4.747,3.67) (4.765,3.67) (4.783,3.669) (4.801,3.669) (4.819,3.669) (4.837,3.668) (4.855,3.668) (4.873,3.667) (4.891,3.667) (4.91,3.667) (4.928,3.666) (4.946,3.666) (4.964,3.666) (4.982,3.665) (5,3.665) (5.019,3.665) (5.037,3.664) (5.055,3.664) (5.074,3.664) (5.092,3.663) (5.109,3.663) (5.127,3.663) (5.145,3.662) (5.162,3.662) (5.18,3.662) (5.197,3.661) (5.215,3.661) (5.233,3.661) (5.251,3.66) (5.268,3.66) (5.286,3.66) (5.304,3.66) (5.322,3.659) (5.34,3.659) (5.358,3.659) (5.376,3.658) (5.394,3.658) (5.412,3.658) (5.43,3.658) (5.448,3.657) (5.466,3.657) (5.484,3.657) (5.503,3.656) (5.521,3.656) (5.539,3.656) (5.558,3.656) (5.576,3.655) (5.594,3.655) (5.611,3.655) (5.629,3.655) (5.647,3.654) (5.664,3.654) (5.682,3.654) (5.7,3.654) (5.718,3.653) (5.735,3.653) (5.753,3.653) (5.771,3.653) (5.789,3.652) (5.807,3.652) (5.825,3.652) (5.844,3.652) (5.862,3.652) (5.88,3.651) (5.898,3.651) (5.917,3.651) (5.935,3.651) (5.954,3.65) (5.971,3.65) (5.989,3.65) (6.006,3.65) (6.024,3.65) (6.042,3.649) (6.059,3.649) (6.077,3.649) (6.095,3.649) (6.113,3.649) (6.131,3.648) (6.149,3.648) (6.167,3.648) (6.185,3.648) (6.203,3.648) (6.221,3.647) (6.24,3.647) (6.258,3.647) (6.276,3.647) (6.295,3.647) (6.313,3.647) (6.331,3.646) (6.348,3.646) (6.366,3.646) (6.384,3.646) (6.402,3.646) (6.419,3.645) (6.437,3.645) (6.455,3.645) (6.473,3.645) (6.491,3.645) (6.509,3.645) (6.527,3.644) (6.545,3.644) (6.564,3.644) (6.582,3.644) (6.6,3.644) (6.619,3.644) (6.637,3.644) (6.656,3.643) (6.674,3.643) (6.692,3.643) (6.709,3.643) (6.727,3.643) (6.745,3.643) (6.762,3.642) (6.78,3.642) (6.798,3.642) (6.816,3.642) (6.834,3.642) (6.852,3.642) (6.87,3.642) (6.889,3.642) (6.907,3.641) (6.925,3.641) (6.943,3.641) (6.962,3.641) (6.98,3.641) (6.999,3.641) (7.018,3.641) (7.035,3.64) (7.053,3.64) (7.07,3.64) (7.088,3.64) (7.106,3.64) (7.124,3.64) (7.141,3.64) (7.159,3.64) (7.177,3.639) (7.195,3.639) (7.213,3.639) (7.232,3.639) (7.25,3.639) (7.268,3.639) (7.287,3.639) (7.305,3.639) (7.324,3.639) (7.342,3.638) (7.361,3.638) (7.379,3.638) (7.397,3.638) (7.414,3.638) (7.432,3.638) (7.45,3.638) (7.468,3.638) (7.485,3.638) (7.503,3.638) (7.521,3.637) (7.539,3.637) (7.557,3.637) (7.575,3.637) (7.594,3.637) (7.612,3.637) (7.63,3.637) (7.649,3.637) (7.667,3.637) (7.686,3.637) (7.704,3.636) (7.723,3.636) (7.742,3.636) (7.76,3.636) (7.778,3.636) (7.795,3.636) (7.813,3.636) (7.831,3.636) (7.849,3.636) (7.866,3.636) (7.884,3.636) (7.902,3.635) (7.92,3.635) (7.938,3.635) (7.957,3.635) (7.975,3.635) (7.993,3.635) (8.011,3.635) (8.03,3.635) (8.048,3.635) (8.067,3.635) (8.085,3.635) (8.104,3.635) (8.123,3.635) (8.142,3.634) (8.161,3.634) (8.178,3.634) (8.196,3.634) (8.213,3.634) (8.231,3.634) (8.249,3.634) (8.267,3.634) (8.285,3.634) (8.303,3.634) (8.321,3.634) (8.339,3.634) (8.357,3.634) (8.375,3.634) (8.393,3.633) (8.412,3.633) (8.43,3.633) (8.449,3.633) (8.467,3.633) (8.486,3.633) (8.505,3.633) (8.523,3.633) (8.542,3.633) (8.561,3.633) (8.58,3.633) (8.599,3.633) (8.617,3.633) (8.634,3.633) (8.652,3.633) (8.67,3.633) (8.688,3.632) (8.706,3.632) (8.724,3.632) (8.742,3.632) (8.76,3.632) (8.778,3.632) (8.796,3.632) (8.814,3.632) (8.833,3.632) (8.851,3.632) (8.869,3.632) (8.888,3.632) (8.907,3.632) (8.925,3.632) (8.944,3.632) (8.963,3.632) (8.982,3.632) (9.001,3.632) (9.02,3.632) (9.039,3.631) (9.058,3.631) (9.076,3.631) (9.093,3.631) (9.111,3.631) (9.129,3.631) (9.146,3.631) (9.164,3.631) (9.182,3.631) (9.2,3.631) (9.218,3.631) (9.236,3.631) (9.254,3.631) (9.272,3.631) (9.291,3.631) (9.309,3.631) (9.328,3.631) (9.346,3.631) (9.365,3.631) (9.383,3.631) (9.402,3.631) (9.421,3.631) (9.44,3.63) (9.459,3.63) (9.478,3.63) (9.497,3.63) (9.516,3.63) (9.535,3.63) (9.554,3.63) (9.574,3.63) (9.593,3.63) (9.611,3.63) (9.628,3.63) (9.646,3.63) (9.664,3.63) (9.682,3.63) (9.699,3.63) (9.717,3.63) (9.735,3.63) (9.753,3.63) (9.772,3.63) (9.79,3.63) (9.808,3.63) (9.826,3.63) (9.845,3.63) (9.863,3.63) (9.882,3.63) (9.9,3.63) (9.919,3.63) (9.938,3.63) (9.957,3.629) (9.975,3.629) (9.994,3.629) (10.013,3.629) (10.032,3.629) (10.052,3.629) (10.071,3.629) (10.09,3.629) (10.11,3.629) (10.129,3.629) (10.148,3.629) (10.168,3.629) (10.188,3.629) (10.205,3.629) (10.223,3.629) (10.241,3.629) (10.258,3.629) (10.276,3.629) (10.294,3.629) (10.312,3.629) (10.33,3.629) (10.348,3.629) (10.366,3.629) (10.384,3.629) (10.402,3.629) (10.421,3.629) (10.439,3.629) (10.457,3.629) (10.476,3.629) (10.494,3.629) (10.513,3.629) (10.532,3.629) (10.55,3.629) (10.569,3.629) (10.588,3.629) (10.607,3.629) (10.626,3.628) (10.645,3.628) (10.664,3.628) (10.683,3.628) (10.703,3.628) (10.722,3.628) (10.741,3.628) (10.761,3.628) (10.78,3.628) (10.8,3.628) (10.82,3.628) (10.84,3.628) (10.859,3.628) (10.879,3.628) (10.899,3.628) (10.917,3.628) (10.934,3.628) (10.952,3.628) (10.97,3.628) (10.987,3.628) (11.005,3.628) (11.023,3.628) (11.041,3.628) (11.059,3.628) (11.077,3.628) (11.095,3.628) (11.113,3.628) (11.131,3.628) (11.15,3.628) (11.168,3.628) (11.186,3.628) (11.205,3.628) (11.223,3.628) (11.242,3.628) (11.26,3.628) (11.279,3.628) (11.298,3.628) (11.316,3.628) (11.335,3.628) (11.354,3.628) (11.373,3.628) (11.392,3.628) (11.411,3.628) (11.43,3.628) (11.449,3.628) (11.469,3.628) (11.488,3.628) (11.507,3.628) (11.527,3.628) (11.546,3.628) (11.566,3.628) (11.585,3.628) (11.605,3.628) (11.625,3.627) (11.644,3.627) (11.664,3.627) (11.684,3.627) (11.704,3.627) (11.724,3.627) (11.744,3.627) (11.764,3.627) (11.784,3.627) (11.804,3.627) (11.825,3.627) (11.845,3.627) (11.865,3.627) (11.883,3.627) (11.9,3.627) (11.918,3.627) (11.936,3.627) (11.953,3.627) (11.971,3.627) (11.989,3.627) (12.006,3.627) (12.024,3.627) (12.042,3.627) (12.06,3.627) (12.078,3.627) (12.096,3.627) (12.114,3.627) (12.132,3.627) (12.15,3.627) (12.168,3.627) (12.186,3.627) (12.205,3.627) (12.223,3.627) (12.241,3.627) (12.259,3.627) (12.278,3.627) (12.296,3.627) (12.315,3.627) (12.333,3.627) (12.351,3.627) (12.37,3.627) (12.389,3.627) (12.407,3.627) (12.426,3.627) (12.444,3.627) (12.463,3.627) (12.482,3.627) (12.5,3.627) (12.519,3.627) (12.538,3.627) (12.557,3.627) (12.576,3.627) (12.594,3.627) (12.613,3.627) (12.632,3.627) (12.651,3.627) (12.67,3.627) (12.689,3.627) (12.708,3.627) (12.727,3.627) (12.746,3.627) (12.765,3.627) (12.784,3.627) (12.803,3.627) (12.822,3.627) (12.841,3.627) (12.86,3.627) (12.879,3.627) (12.898,3.627) (12.917,3.627) (12.936,3.627) (12.955,3.627) (12.974,3.627) (12.993,3.627) (13.012,3.627) (13.032,3.627) (13.051,3.627) (13.07,3.627) (13.089,3.627) (13.108,3.627) (13.127,3.627) (13.146,3.627) (13.165,3.627) (13.184,3.627) (13.203,3.627) (13.221,3.627) (13.24,3.627) (13.259,3.627) (13.278,3.627) (13.297,3.627) (13.316,3.627) (13.334,3.627) (13.353,3.627) (13.372,3.627) (13.39,3.627) (13.409,3.627) (13.427,3.627) (13.446,3.627) (13.464,3.627) (13.482,3.627) (13.501,3.627) (13.519,3.627) (13.537,3.627) (13.555,3.627) (13.573,3.627) (13.591,3.627) (13.609,3.626) (13.626,3.626) (13.644,3.626) (13.661,3.626) (13.682,3.626) (13.702,3.626) (13.722,3.626) (13.742,3.626) (13.762,3.626) (13.781,3.626) (13.801,3.626) (13.82,3.626) (13.839,3.626) (13.858,3.626) (13.877,3.626) (13.895,3.626) (13.914,3.626) (13.932,3.626) (13.95,3.626) (13.967,3.626) (13.987,3.626) (14.007,3.626) (14.026,3.626) (14.045,3.626) (14.064,3.626) (14.082,3.626) (14.1,3.626) (14.118,3.626) (14.137,3.626) (14.156,3.626) (14.174,3.626) (14.192,3.626) (14.211,3.626) (14.229,3.626) (14.247,3.626) (14.266,3.626) (14.284,3.626) (14.302,3.626) (14.321,3.626) (14.339,3.626) (14.357,3.626) (14.375,3.626) (14.394,3.626) (14.411,3.626) (14.429,3.626) (14.431,3.626)};
\draw[BeachCyan,dashed,line width=1.1pt] plot coordinates {(2.619,-0.526) (2.637,0.05) (2.659,0.4) (2.681,0.641) (2.699,0.786) (2.719,0.918) (2.74,1.037) (2.758,1.126) (2.778,1.211) (2.797,1.28) (2.816,1.346) (2.834,1.403) (2.854,1.458) (2.872,1.507) (2.891,1.554) (2.91,1.596) (2.928,1.634) (2.946,1.672) (2.964,1.706) (2.984,1.741) (3.002,1.772) (3.02,1.801) (3.039,1.83) (3.057,1.857) (3.075,1.882) (3.094,1.907) (3.112,1.931) (3.13,1.953) (3.148,1.974) (3.166,1.996) (3.184,2.016) (3.202,2.035) (3.221,2.054) (3.239,2.072) (3.257,2.09) (3.275,2.106) (3.293,2.122) (3.311,2.138) (3.329,2.154) (3.347,2.169) (3.365,2.183) (3.383,2.197) (3.401,2.211) (3.419,2.224) (3.437,2.237) (3.455,2.25) (3.473,2.262) (3.49,2.273) (3.509,2.285) (3.527,2.297) (3.545,2.308) (3.563,2.319) (3.581,2.329) (3.598,2.339) (3.616,2.349) (3.634,2.359) (3.652,2.369) (3.671,2.379) (3.689,2.388) (3.706,2.397) (3.724,2.406) (3.742,2.414) (3.76,2.423) (3.777,2.431) (3.796,2.439) (3.814,2.448) (3.832,2.456) (3.85,2.463) (3.868,2.471) (3.886,2.478) (3.904,2.486) (3.922,2.493) (3.94,2.5) (3.958,2.507) (3.975,2.513) (3.993,2.52) (4.011,2.526) (4.028,2.533) (4.046,2.539) (4.064,2.545) (4.082,2.551) (4.101,2.557) (4.119,2.563) (4.137,2.569) (4.155,2.575) (4.173,2.581) (4.192,2.586) (4.21,2.592) (4.228,2.597) (4.246,2.602) (4.264,2.608) (4.282,2.613) (4.3,2.618) (4.318,2.623) (4.336,2.628) (4.354,2.632) (4.371,2.637) (4.389,2.642) (4.407,2.646) (4.425,2.651) (4.443,2.655) (4.461,2.66) (4.479,2.664) (4.497,2.668) (4.515,2.672) (4.532,2.676) (4.55,2.68) (4.568,2.684) (4.586,2.688) (4.604,2.692) (4.622,2.696) (4.64,2.7) (4.658,2.704) (4.676,2.707) (4.693,2.711) (4.711,2.715) (4.729,2.718) (4.747,2.722) (4.765,2.725) (4.783,2.729) (4.801,2.732) (4.819,2.735) (4.837,2.739) (4.855,2.742) (4.873,2.745) (4.891,2.748) (4.91,2.751) (4.928,2.755) (4.946,2.758) (4.964,2.761) (4.982,2.764) (5,2.767) (5.019,2.77) (5.037,2.772) (5.055,2.775) (5.074,2.778) (5.092,2.781) (5.109,2.784) (5.127,2.786) (5.145,2.789) (5.162,2.791) (5.18,2.794) (5.197,2.796) (5.215,2.799) (5.233,2.801) (5.251,2.804) (5.268,2.806) (5.286,2.809) (5.304,2.811) (5.322,2.813) (5.34,2.816) (5.358,2.818) (5.376,2.82) (5.394,2.822) (5.412,2.825) (5.43,2.827) (5.448,2.829) (5.466,2.831) (5.484,2.833) (5.503,2.836) (5.521,2.838) (5.539,2.84) (5.558,2.842) (5.576,2.844) (5.594,2.846) (5.611,2.848) (5.629,2.85) (5.647,2.851) (5.664,2.853) (5.682,2.855) (5.7,2.857) (5.718,2.859) (5.735,2.861) (5.753,2.862) (5.771,2.864) (5.789,2.866) (5.807,2.868) (5.825,2.869) (5.844,2.871) (5.862,2.873) (5.88,2.874) (5.898,2.876) (5.917,2.878) (5.935,2.879) (5.954,2.881) (5.971,2.883) (5.989,2.884) (6.006,2.886) (6.024,2.887) (6.042,2.889) (6.059,2.89) (6.077,2.892) (6.095,2.893) (6.113,2.894) (6.131,2.896) (6.149,2.897) (6.167,2.899) (6.185,2.9) (6.203,2.902) (6.221,2.903) (6.24,2.904) (6.258,2.906) (6.276,2.907) (6.295,2.908) (6.313,2.91) (6.331,2.911) (6.348,2.912) (6.366,2.913) (6.384,2.915) (6.402,2.916) (6.419,2.917) (6.437,2.918) (6.455,2.919) (6.473,2.92) (6.491,2.922) (6.509,2.923) (6.527,2.924) (6.545,2.925) (6.564,2.926) (6.582,2.927) (6.6,2.928) (6.619,2.93) (6.637,2.931) (6.656,2.932) (6.674,2.933) (6.692,2.934) (6.709,2.935) (6.727,2.936) (6.745,2.937) (6.762,2.938) (6.78,2.939) (6.798,2.94) (6.816,2.941) (6.834,2.942) (6.852,2.943) (6.87,2.944) (6.889,2.945) (6.907,2.946) (6.925,2.946) (6.943,2.947) (6.962,2.948) (6.98,2.949) (6.999,2.95) (7.018,2.951) (7.035,2.952) (7.053,2.953) (7.07,2.954) (7.088,2.954) (7.106,2.955) (7.124,2.956) (7.141,2.957) (7.159,2.958) (7.177,2.958) (7.195,2.959) (7.213,2.96) (7.232,2.961) (7.25,2.962) (7.268,2.962) (7.287,2.963) (7.305,2.964) (7.324,2.965) (7.342,2.965) (7.361,2.966) (7.379,2.967) (7.397,2.968) (7.414,2.968) (7.432,2.969) (7.45,2.97) (7.468,2.97) (7.485,2.971) (7.503,2.972) (7.521,2.972) (7.539,2.973) (7.557,2.974) (7.575,2.974) (7.594,2.975) (7.612,2.976) (7.63,2.976) (7.649,2.977) (7.667,2.978) (7.686,2.978) (7.704,2.979) (7.723,2.979) (7.742,2.98) (7.76,2.981) (7.778,2.981) (7.795,2.982) (7.813,2.982) (7.831,2.983) (7.849,2.983) (7.866,2.984) (7.884,2.985) (7.902,2.985) (7.92,2.986) (7.938,2.986) (7.957,2.987) (7.975,2.987) (7.993,2.988) (8.011,2.988) (8.03,2.989) (8.048,2.989) (8.067,2.99) (8.085,2.99) (8.104,2.991) (8.123,2.991) (8.142,2.992) (8.161,2.992) (8.178,2.993) (8.196,2.993) (8.213,2.994) (8.231,2.994) (8.249,2.995) (8.267,2.995) (8.285,2.996) (8.303,2.996) (8.321,2.996) (8.339,2.997) (8.357,2.997) (8.375,2.998) (8.393,2.998) (8.412,2.999) (8.43,2.999) (8.449,2.999) (8.467,3) (8.486,3) (8.505,3.001) (8.523,3.001) (8.542,3.002) (8.561,3.002) (8.58,3.002) (8.599,3.003) (8.617,3.003) (8.634,3.004) (8.652,3.004) (8.67,3.004) (8.688,3.005) (8.706,3.005) (8.724,3.005) (8.742,3.006) (8.76,3.006) (8.778,3.006) (8.796,3.007) (8.814,3.007) (8.833,3.007) (8.851,3.008) (8.869,3.008) (8.888,3.008) (8.907,3.009) (8.925,3.009) (8.944,3.01) (8.963,3.01) (8.982,3.01) (9.001,3.01) (9.02,3.011) (9.039,3.011) (9.058,3.011) (9.076,3.012) (9.093,3.012) (9.111,3.012) (9.129,3.013) (9.146,3.013) (9.164,3.013) (9.182,3.014) (9.2,3.014) (9.218,3.014) (9.236,3.014) (9.254,3.015) (9.272,3.015) (9.291,3.015) (9.309,3.015) (9.328,3.016) (9.346,3.016) (9.365,3.016) (9.383,3.017) (9.402,3.017) (9.421,3.017) (9.44,3.017) (9.459,3.018) (9.478,3.018) (9.497,3.018) (9.516,3.018) (9.535,3.019) (9.554,3.019) (9.574,3.019) (9.593,3.019) (9.611,3.02) (9.628,3.02) (9.646,3.02) (9.664,3.02) (9.682,3.021) (9.699,3.021) (9.717,3.021) (9.735,3.021) (9.753,3.021) (9.772,3.022) (9.79,3.022) (9.808,3.022) (9.826,3.022) (9.845,3.022) (9.863,3.023) (9.882,3.023) (9.9,3.023) (9.919,3.023) (9.938,3.023) (9.957,3.024) (9.975,3.024) (9.994,3.024) (10.013,3.024) (10.032,3.024) (10.052,3.025) (10.071,3.025) (10.09,3.025) (10.11,3.025) (10.129,3.025) (10.148,3.026) (10.168,3.026) (10.188,3.026) (10.205,3.026) (10.223,3.026) (10.241,3.027) (10.258,3.027) (10.276,3.027) (10.294,3.027) (10.312,3.027) (10.33,3.027) (10.348,3.028) (10.366,3.028) (10.384,3.028) (10.402,3.028) (10.421,3.028) (10.439,3.028) (10.457,3.028) (10.476,3.029) (10.494,3.029) (10.513,3.029) (10.532,3.029) (10.55,3.029) (10.569,3.029) (10.588,3.03) (10.607,3.03) (10.626,3.03) (10.645,3.03) (10.664,3.03) (10.683,3.03) (10.703,3.03) (10.722,3.031) (10.741,3.031) (10.761,3.031) (10.78,3.031) (10.8,3.031) (10.82,3.031) (10.84,3.031) (10.859,3.032) (10.879,3.032) (10.899,3.032) (10.917,3.032) (10.934,3.032) (10.952,3.032) (10.97,3.032) (10.987,3.032) (11.005,3.033) (11.023,3.033) (11.041,3.033) (11.059,3.033) (11.077,3.033) (11.095,3.033) (11.113,3.033) (11.131,3.033) (11.15,3.033) (11.168,3.034) (11.186,3.034) (11.205,3.034) (11.223,3.034) (11.242,3.034) (11.26,3.034) (11.279,3.034) (11.298,3.034) (11.316,3.034) (11.335,3.035) (11.354,3.035) (11.373,3.035) (11.392,3.035) (11.411,3.035) (11.43,3.035) (11.449,3.035) (11.469,3.035) (11.488,3.035) (11.507,3.035) (11.527,3.036) (11.546,3.036) (11.566,3.036) (11.585,3.036) (11.605,3.036) (11.625,3.036) (11.644,3.036) (11.664,3.036) (11.684,3.036) (11.704,3.036) (11.724,3.037) (11.744,3.037) (11.764,3.037) (11.784,3.037) (11.804,3.037) (11.825,3.037) (11.845,3.037) (11.865,3.037) (11.883,3.037) (11.9,3.037) (11.918,3.037) (11.936,3.037) (11.953,3.038) (11.971,3.038) (11.989,3.038) (12.006,3.038) (12.024,3.038) (12.042,3.038) (12.06,3.038) (12.078,3.038) (12.096,3.038) (12.114,3.038) (12.132,3.038) (12.15,3.038) (12.168,3.038) (12.186,3.038) (12.205,3.038) (12.223,3.039) (12.241,3.039) (12.259,3.039) (12.278,3.039) (12.296,3.039) (12.315,3.039) (12.333,3.039) (12.351,3.039) (12.37,3.039) (12.389,3.039) (12.407,3.039) (12.426,3.039) (12.444,3.039) (12.463,3.039) (12.482,3.039) (12.5,3.039) (12.519,3.04) (12.538,3.04) (12.557,3.04) (12.576,3.04) (12.594,3.04) (12.613,3.04) (12.632,3.04) (12.651,3.04) (12.67,3.04) (12.689,3.04) (12.708,3.04) (12.727,3.04) (12.746,3.04) (12.765,3.04) (12.784,3.04) (12.803,3.04) (12.822,3.04) (12.841,3.04) (12.86,3.041) (12.879,3.041) (12.898,3.041) (12.917,3.041) (12.936,3.041) (12.955,3.041) (12.974,3.041) (12.993,3.041) (13.012,3.041) (13.032,3.041) (13.051,3.041) (13.07,3.041) (13.089,3.041) (13.108,3.041) (13.127,3.041) (13.146,3.041) (13.165,3.041) (13.184,3.041) (13.203,3.041) (13.221,3.041) (13.24,3.041) (13.259,3.042) (13.278,3.042) (13.297,3.042) (13.316,3.042) (13.334,3.042) (13.353,3.042) (13.372,3.042) (13.39,3.042) (13.409,3.042) (13.427,3.042) (13.446,3.042) (13.464,3.042) (13.482,3.042) (13.501,3.042) (13.519,3.042) (13.537,3.042) (13.555,3.042) (13.573,3.042) (13.591,3.042) (13.609,3.042) (13.626,3.042) (13.644,3.042) (13.661,3.042) (13.682,3.042) (13.702,3.042) (13.722,3.042) (13.742,3.042) (13.762,3.042) (13.781,3.043) (13.801,3.043) (13.82,3.043) (13.839,3.043) (13.858,3.043) (13.877,3.043) (13.895,3.043) (13.914,3.043) (13.932,3.043) (13.95,3.043) (13.967,3.043) (13.987,3.043) (14.007,3.043) (14.026,3.043) (14.045,3.043) (14.064,3.043) (14.082,3.043) (14.1,3.043) (14.118,3.043) (14.137,3.043) (14.156,3.043) (14.174,3.043) (14.192,3.043) (14.211,3.043) (14.229,3.043) (14.247,3.043) (14.266,3.043) (14.284,3.043) (14.302,3.043) (14.321,3.043) (14.339,3.043) (14.357,3.043) (14.375,3.043) (14.394,3.043) (14.411,3.043) (14.429,3.044) (14.431,3.044)};
\draw[BeachOrange,line width=1.1pt] plot coordinates {(0,2.857) (0.027,2.872) (0.044,2.883) (0.062,2.893) (0.079,2.902) (0.097,2.912) (0.114,2.922) (0.132,2.931) (0.15,2.94) (0.167,2.949) (0.185,2.958) (0.203,2.967) (0.22,2.976) (0.238,2.984) (0.256,2.993) (0.273,3.001) (0.291,3.009) (0.309,3.017) (0.326,3.025) (0.344,3.033) (0.361,3.041) (0.379,3.049) (0.397,3.056) (0.415,3.064) (0.432,3.072) (0.45,3.079) (0.468,3.086) (0.486,3.094) (0.503,3.101) (0.521,3.108) (0.538,3.115) (0.556,3.122) (0.574,3.129) (0.592,3.136) (0.61,3.143) (0.627,3.149) (0.645,3.156) (0.663,3.162) (0.68,3.169) (0.698,3.175) (0.716,3.182) (0.734,3.188) (0.751,3.195) (0.769,3.201) (0.787,3.207) (0.805,3.213) (0.822,3.219) (0.84,3.226) (0.858,3.232) (0.876,3.238) (0.894,3.244) (0.912,3.249) (0.93,3.255) (0.948,3.261) (0.965,3.267) (0.983,3.273) (1.001,3.278) (1.019,3.284) (1.036,3.289) (1.054,3.295) (1.072,3.301) (1.09,3.306) (1.108,3.311) (1.126,3.317) (1.143,3.322) (1.161,3.327) (1.179,3.332) (1.197,3.338) (1.215,3.343) (1.233,3.348) (1.251,3.353) (1.268,3.358) (1.286,3.363) (1.304,3.368) (1.322,3.373) (1.34,3.378) (1.357,3.383) (1.375,3.388) (1.393,3.393) (1.412,3.397) (1.429,3.402) (1.447,3.407) (1.465,3.412) (1.484,3.416) (1.501,3.421) (1.519,3.425) (1.537,3.43) (1.555,3.434) (1.574,3.439) (1.591,3.443) (1.609,3.448) (1.627,3.452) (1.645,3.457) (1.664,3.461) (1.681,3.465) (1.699,3.47) (1.717,3.474) (1.735,3.478) (1.754,3.482) (1.772,3.487) (1.79,3.491) (1.808,3.495) (1.826,3.499) (1.844,3.503) (1.863,3.507) (1.881,3.511) (1.899,3.515) (1.917,3.519) (1.935,3.523) (1.953,3.527) (1.972,3.53) (1.99,3.534) (2.008,3.538) (2.026,3.542) (2.044,3.545) (2.062,3.549) (2.08,3.553) (2.099,3.557) (2.117,3.56) (2.135,3.564) (2.153,3.567) (2.171,3.571) (2.189,3.574) (2.207,3.578) (2.225,3.581) (2.244,3.585) (2.263,3.589) (2.281,3.592) (2.298,3.595) (2.316,3.598) (2.334,3.602) (2.353,3.605) (2.371,3.608) (2.39,3.612) (2.409,3.615) (2.428,3.619) (2.445,3.622) (2.463,3.625) (2.481,3.628) (2.499,3.631) (2.517,3.634) (2.535,3.637) (2.553,3.64) (2.572,3.643) (2.619,3.646) (2.637,3.642) (2.659,3.639) (2.681,3.635) (2.699,3.633) (2.719,3.63) (2.74,3.627) (2.758,3.624) (2.778,3.621) (2.797,3.618) (2.816,3.615) (2.834,3.613) (2.854,3.61) (2.872,3.608) (2.891,3.605) (2.91,3.603) (2.928,3.6) (2.946,3.598) (2.964,3.595) (2.984,3.593) (3.002,3.591) (3.02,3.588) (3.039,3.586) (3.057,3.584) (3.075,3.582) (3.094,3.579) (3.112,3.577) (3.13,3.575) (3.148,3.573) (3.166,3.571) (3.184,3.569) (3.202,3.567) (3.221,3.565) (3.239,3.563) (3.257,3.561) (3.275,3.559) (3.293,3.557) (3.311,3.555) (3.329,3.553) (3.347,3.551) (3.365,3.549) (3.383,3.548) (3.401,3.546) (3.419,3.544) (3.437,3.542) (3.455,3.541) (3.473,3.539) (3.49,3.537) (3.509,3.536) (3.527,3.534) (3.545,3.532) (3.563,3.531) (3.581,3.529) (3.598,3.527) (3.616,3.526) (3.634,3.524) (3.652,3.523) (3.671,3.521) (3.689,3.52) (3.706,3.518) (3.724,3.517) (3.742,3.515) (3.76,3.514) (3.777,3.513) (3.796,3.511) (3.814,3.51) (3.832,3.508) (3.85,3.507) (3.868,3.505) (3.886,3.504) (3.904,3.503) (3.922,3.501) (3.94,3.5) (3.958,3.499) (3.975,3.498) (3.993,3.496) (4.011,3.495) (4.028,3.494) (4.046,3.493) (4.064,3.491) (4.082,3.49) (4.101,3.489) (4.119,3.488) (4.137,3.487) (4.155,3.485) (4.173,3.484) (4.192,3.483) (4.21,3.482) (4.228,3.481) (4.246,3.48) (4.264,3.479) (4.282,3.478) (4.3,3.477) (4.318,3.475) (4.336,3.474) (4.354,3.473) (4.371,3.472) (4.389,3.471) (4.407,3.47) (4.425,3.469) (4.443,3.468) (4.461,3.467) (4.479,3.466) (4.497,3.465) (4.515,3.465) (4.532,3.464) (4.55,3.463) (4.568,3.462) (4.586,3.461) (4.604,3.46) (4.622,3.459) (4.64,3.458) (4.658,3.457) (4.676,3.456) (4.693,3.456) (4.711,3.455) (4.729,3.454) (4.747,3.453) (4.765,3.452) (4.783,3.451) (4.801,3.45) (4.819,3.45) (4.837,3.449) (4.855,3.448) (4.873,3.447) (4.891,3.446) (4.91,3.446) (4.928,3.445) (4.946,3.444) (4.964,3.443) (4.982,3.443) (5,3.442) (5.019,3.441) (5.037,3.44) (5.055,3.44) (5.074,3.439) (5.092,3.438) (5.109,3.438) (5.127,3.437) (5.145,3.436) (5.162,3.436) (5.18,3.435) (5.197,3.434) (5.215,3.434) (5.233,3.433) (5.251,3.432) (5.268,3.432) (5.286,3.431) (5.304,3.43) (5.322,3.43) (5.34,3.429) (5.358,3.429) (5.376,3.428) (5.394,3.427) (5.412,3.427) (5.43,3.426) (5.448,3.426) (5.466,3.425) (5.484,3.424) (5.503,3.424) (5.521,3.423) (5.539,3.423) (5.558,3.422) (5.576,3.422) (5.594,3.421) (5.611,3.421) (5.629,3.42) (5.647,3.42) (5.664,3.419) (5.682,3.419) (5.7,3.418) (5.718,3.418) (5.735,3.417) (5.753,3.417) (5.771,3.416) (5.789,3.416) (5.807,3.415) (5.825,3.415) (5.844,3.414) (5.862,3.414) (5.88,3.413) (5.898,3.413) (5.917,3.412) (5.935,3.412) (5.954,3.411) (5.971,3.411) (5.989,3.41) (6.006,3.41) (6.024,3.41) (6.042,3.409) (6.059,3.409) (6.077,3.408) (6.095,3.408) (6.113,3.407) (6.131,3.407) (6.149,3.407) (6.167,3.406) (6.185,3.406) (6.203,3.405) (6.221,3.405) (6.24,3.405) (6.258,3.404) (6.276,3.404) (6.295,3.403) (6.313,3.403) (6.331,3.403) (6.348,3.402) (6.366,3.402) (6.384,3.402) (6.402,3.401) (6.419,3.401) (6.437,3.401) (6.455,3.4) (6.473,3.4) (6.491,3.4) (6.509,3.399) (6.527,3.399) (6.545,3.399) (6.564,3.398) (6.582,3.398) (6.6,3.398) (6.619,3.397) (6.637,3.397) (6.656,3.397) (6.674,3.396) (6.692,3.396) (6.709,3.396) (6.727,3.395) (6.745,3.395) (6.762,3.395) (6.78,3.394) (6.798,3.394) (6.816,3.394) (6.834,3.394) (6.852,3.393) (6.87,3.393) (6.889,3.393) (6.907,3.392) (6.925,3.392) (6.943,3.392) (6.962,3.392) (6.98,3.391) (6.999,3.391) (7.018,3.391) (7.035,3.39) (7.053,3.39) (7.07,3.39) (7.088,3.39) (7.106,3.389) (7.124,3.389) (7.141,3.389) (7.159,3.389) (7.177,3.388) (7.195,3.388) (7.213,3.388) (7.232,3.388) (7.25,3.388) (7.268,3.387) (7.287,3.387) (7.305,3.387) (7.324,3.387) (7.342,3.386) (7.361,3.386) (7.379,3.386) (7.397,3.386) (7.414,3.385) (7.432,3.385) (7.45,3.385) (7.468,3.385) (7.485,3.385) (7.503,3.384) (7.521,3.384) (7.539,3.384) (7.557,3.384) (7.575,3.384) (7.594,3.383) (7.612,3.383) (7.63,3.383) (7.649,3.383) (7.667,3.383) (7.686,3.382) (7.704,3.382) (7.723,3.382) (7.742,3.382) (7.76,3.382) (7.778,3.381) (7.795,3.381) (7.813,3.381) (7.831,3.381) (7.849,3.381) (7.866,3.381) (7.884,3.38) (7.902,3.38) (7.92,3.38) (7.938,3.38) (7.957,3.38) (7.975,3.38) (7.993,3.379) (8.011,3.379) (8.03,3.379) (8.048,3.379) (8.067,3.379) (8.085,3.379) (8.104,3.378) (8.123,3.378) (8.142,3.378) (8.161,3.378) (8.178,3.378) (8.196,3.378) (8.213,3.377) (8.231,3.377) (8.249,3.377) (8.267,3.377) (8.285,3.377) (8.303,3.377) (8.321,3.377) (8.339,3.376) (8.357,3.376) (8.375,3.376) (8.393,3.376) (8.412,3.376) (8.43,3.376) (8.449,3.376) (8.467,3.376) (8.486,3.375) (8.505,3.375) (8.523,3.375) (8.542,3.375) (8.561,3.375) (8.58,3.375) (8.599,3.375) (8.617,3.374) (8.634,3.374) (8.652,3.374) (8.67,3.374) (8.688,3.374) (8.706,3.374) (8.724,3.374) (8.742,3.374) (8.76,3.374) (8.778,3.373) (8.796,3.373) (8.814,3.373) (8.833,3.373) (8.851,3.373) (8.869,3.373) (8.888,3.373) (8.907,3.373) (8.925,3.373) (8.944,3.372) (8.963,3.372) (8.982,3.372) (9.001,3.372) (9.02,3.372) (9.039,3.372) (9.058,3.372) (9.076,3.372) (9.093,3.372) (9.111,3.372) (9.129,3.371) (9.146,3.371) (9.164,3.371) (9.182,3.371) (9.2,3.371) (9.218,3.371) (9.236,3.371) (9.254,3.371) (9.272,3.371) (9.291,3.371) (9.309,3.371) (9.328,3.37) (9.346,3.37) (9.365,3.37) (9.383,3.37) (9.402,3.37) (9.421,3.37) (9.44,3.37) (9.459,3.37) (9.478,3.37) (9.497,3.37) (9.516,3.37) (9.535,3.37) (9.554,3.369) (9.574,3.369) (9.593,3.369) (9.611,3.369) (9.628,3.369) (9.646,3.369) (9.664,3.369) (9.682,3.369) (9.699,3.369) (9.717,3.369) (9.735,3.369) (9.753,3.369) (9.772,3.369) (9.79,3.368) (9.808,3.368) (9.826,3.368) (9.845,3.368) (9.863,3.368) (9.882,3.368) (9.9,3.368) (9.919,3.368) (9.938,3.368) (9.957,3.368) (9.975,3.368) (9.994,3.368) (10.013,3.368) (10.032,3.368) (10.052,3.368) (10.071,3.367) (10.09,3.367) (10.11,3.367) (10.129,3.367) (10.148,3.367) (10.168,3.367) (10.188,3.367) (10.205,3.367) (10.223,3.367) (10.241,3.367) (10.258,3.367) (10.276,3.367) (10.294,3.367) (10.312,3.367) (10.33,3.367) (10.348,3.367) (10.366,3.367) (10.384,3.366) (10.402,3.366) (10.421,3.366) (10.439,3.366) (10.457,3.366) (10.476,3.366) (10.494,3.366) (10.513,3.366) (10.532,3.366) (10.55,3.366) (10.569,3.366) (10.588,3.366) (10.607,3.366) (10.626,3.366) (10.645,3.366) (10.664,3.366) (10.683,3.366) (10.703,3.366) (10.722,3.366) (10.741,3.366) (10.761,3.366) (10.78,3.365) (10.8,3.365) (10.82,3.365) (10.84,3.365) (10.859,3.365) (10.879,3.365) (10.899,3.365) (10.917,3.365) (10.934,3.365) (10.952,3.365) (10.97,3.365) (10.987,3.365) (11.005,3.365) (11.023,3.365) (11.041,3.365) (11.059,3.365) (11.077,3.365) (11.095,3.365) (11.113,3.365) (11.131,3.365) (11.15,3.365) (11.168,3.365) (11.186,3.365) (11.205,3.365) (11.223,3.365) (11.242,3.364) (11.26,3.364) (11.279,3.364) (11.298,3.364) (11.316,3.364) (11.335,3.364) (11.354,3.364) (11.373,3.364) (11.392,3.364) (11.411,3.364) (11.43,3.364) (11.449,3.364) (11.469,3.364) (11.488,3.364) (11.507,3.364) (11.527,3.364) (11.546,3.364) (11.566,3.364) (11.585,3.364) (11.605,3.364) (11.625,3.364) (11.644,3.364) (11.664,3.364) (11.684,3.364) (11.704,3.364) (11.724,3.364) (11.744,3.364) (11.764,3.364) (11.784,3.364) (11.804,3.364) (11.825,3.364) (11.845,3.363) (11.865,3.363) (11.883,3.363) (11.9,3.363) (11.918,3.363) (11.936,3.363) (11.953,3.363) (11.971,3.363) (11.989,3.363) (12.006,3.363) (12.024,3.363) (12.042,3.363) (12.06,3.363) (12.078,3.363) (12.096,3.363) (12.114,3.363) (12.132,3.363) (12.15,3.363) (12.168,3.363) (12.186,3.363) (12.205,3.363) (12.223,3.363) (12.241,3.363) (12.259,3.363) (12.278,3.363) (12.296,3.363) (12.315,3.363) (12.333,3.363) (12.351,3.363) (12.37,3.363) (12.389,3.363) (12.407,3.363) (12.426,3.363) (12.444,3.363) (12.463,3.363) (12.482,3.363) (12.5,3.363) (12.519,3.363) (12.538,3.363) (12.557,3.363) (12.576,3.363) (12.594,3.363) (12.613,3.363) (12.632,3.363) (12.651,3.363) (12.67,3.363) (12.689,3.362) (12.708,3.362) (12.727,3.362) (12.746,3.362) (12.765,3.362) (12.784,3.362) (12.803,3.362) (12.822,3.362) (12.841,3.362) (12.86,3.362) (12.879,3.362) (12.898,3.362) (12.917,3.362) (12.936,3.362) (12.955,3.362) (12.974,3.362) (12.993,3.362) (13.012,3.362) (13.032,3.362) (13.051,3.362) (13.07,3.362) (13.089,3.362) (13.108,3.362) (13.127,3.362) (13.146,3.362) (13.165,3.362) (13.184,3.362) (13.203,3.362) (13.221,3.362) (13.24,3.362) (13.259,3.362) (13.278,3.362) (13.297,3.362) (13.316,3.362) (13.334,3.362) (13.353,3.362) (13.372,3.362) (13.39,3.362) (13.409,3.362) (13.427,3.362) (13.446,3.362) (13.464,3.362) (13.482,3.362) (13.501,3.362) (13.519,3.362) (13.537,3.362) (13.555,3.362) (13.573,3.362) (13.591,3.362) (13.609,3.362) (13.626,3.362) (13.644,3.362) (13.661,3.362) (13.682,3.362) (13.702,3.362) (13.722,3.362) (13.742,3.362) (13.762,3.362) (13.781,3.362) (13.801,3.362) (13.82,3.362) (13.839,3.362) (13.858,3.362) (13.877,3.362) (13.895,3.362) (13.914,3.362) (13.932,3.362) (13.95,3.362) (13.967,3.362) (13.987,3.362) (14.007,3.362) (14.026,3.362) (14.045,3.362) (14.064,3.362) (14.082,3.362) (14.1,3.361) (14.118,3.361) (14.137,3.361) (14.156,3.361) (14.174,3.361) (14.192,3.361) (14.211,3.361) (14.229,3.361) (14.247,3.361) (14.266,3.361) (14.284,3.361) (14.302,3.361) (14.321,3.361) (14.339,3.361) (14.357,3.361) (14.375,3.361) (14.394,3.361) (14.411,3.361) (14.429,3.361) (14.431,3.361)};
\draw[BeachRed,dash dot,line width=1.1pt] plot coordinates {(0,5.524) (0.005,5.509) (0.027,5.439) (0.044,5.386) (0.062,5.335) (0.079,5.286) (0.097,5.239) (0.114,5.194) (0.132,5.15) (0.15,5.107) (0.167,5.067) (0.185,5.027) (0.203,4.99) (0.22,4.953) (0.238,4.917) (0.256,4.882) (0.273,4.848) (0.291,4.815) (0.309,4.783) (0.326,4.752) (0.344,4.722) (0.361,4.693) (0.379,4.664) (0.397,4.635) (0.415,4.607) (0.432,4.58) (0.45,4.553) (0.468,4.528) (0.486,4.502) (0.503,4.477) (0.521,4.452) (0.538,4.428) (0.556,4.404) (0.574,4.381) (0.592,4.358) (0.61,4.335) (0.627,4.313) (0.645,4.291) (0.663,4.27) (0.68,4.248) (0.698,4.227) (0.716,4.206) (0.734,4.186) (0.751,4.165) (0.769,4.145) (0.787,4.125) (0.805,4.106) (0.822,4.086) (0.84,4.067) (0.858,4.047) (0.876,4.028) (0.894,4.01) (0.912,3.991) (0.93,3.972) (0.948,3.953) (0.965,3.935) (0.983,3.917) (1.001,3.899) (1.019,3.881) (1.036,3.863) (1.054,3.845) (1.072,3.827) (1.09,3.809) (1.108,3.792) (1.126,3.774) (1.143,3.757) (1.161,3.739) (1.179,3.721) (1.197,3.704) (1.215,3.686) (1.233,3.669) (1.251,3.651) (1.268,3.634) (1.286,3.616) (1.304,3.599) (1.322,3.581) (1.34,3.564) (1.357,3.546) (1.375,3.528) (1.393,3.51) (1.412,3.492) (1.429,3.474) (1.447,3.456) (1.465,3.438) (1.484,3.419) (1.501,3.401) (1.519,3.383) (1.537,3.364) (1.555,3.345) (1.574,3.326) (1.591,3.307) (1.609,3.288) (1.627,3.269) (1.645,3.249) (1.664,3.229) (1.681,3.209) (1.699,3.189) (1.717,3.169) (1.735,3.148) (1.754,3.127) (1.772,3.105) (1.79,3.084) (1.808,3.062) (1.826,3.04) (1.844,3.017) (1.863,2.994) (1.881,2.97) (1.899,2.947) (1.917,2.923) (1.935,2.898) (1.953,2.873) (1.972,2.847) (1.99,2.82) (2.008,2.793) (2.026,2.766) (2.044,2.739) (2.062,2.71) (2.08,2.68) (2.099,2.648) (2.117,2.616) (2.135,2.584) (2.153,2.551) (2.171,2.517) (2.189,2.481) (2.207,2.443) (2.225,2.404) (2.244,2.362) (2.263,2.318) (2.281,2.274) (2.298,2.228) (2.316,2.179) (2.334,2.127) (2.353,2.071) (2.371,2.01) (2.39,1.944) (2.409,1.871) (2.428,1.791) (2.445,1.709) (2.463,1.617) (2.481,1.512) (2.499,1.389) (2.517,1.244) (2.535,1.064) (2.553,0.831) (2.572,0.473) (2.587,-0.008)};
\end{scope}
\node[font=\scriptsize\sffamily] at (2.603,6.249) {$z_m$};
\draw[BeachInk,line width=1.2pt] (5.3,2.1)--(5.9,2.1);
\node[right,font=\scriptsize\sffamily] at (5.9,2.1) {イオン $n_i$};
\draw[BeachBlue,line width=1.1pt] (5.3,1.72)--(5.9,1.72);
\node[right,font=\scriptsize\sffamily] at (5.9,1.72) {太陽風電子（通過）};
\draw[BeachCyan,dashed,line width=1.1pt] (5.3,1.34)--(5.9,1.34);
\node[right,font=\scriptsize\sffamily] at (5.9,1.34) {太陽風電子（反射）};
\draw[BeachOrange,line width=1.1pt] (5.3,0.96)--(5.9,0.96);
\node[right,font=\scriptsize\sffamily] at (5.9,0.96) {光電子（脱出）};
\draw[BeachRed,dash dot,line width=1.1pt] (5.3,0.58)--(5.9,0.58);
\node[right,font=\scriptsize\sffamily] at (5.9,0.58) {光電子（捕捉）};
\end{tikzpicture}
\caption[Type A の粒子集団ごとの数密度]{Type A（$\alpha=60^\circ$）の粒子集団ごとの数密度。縦軸は対数で、点線は電位極小の位置 $z_m$ である。反射した太陽風電子は極小より上流側だけに、捕捉された光電子は極小より壁側だけに存在する。壁の近くでは捕捉された光電子が負電荷を担い、上流では通過・反射電子と脱出光電子の和がイオン密度と釣り合う。各密度は電位から速度積分で直接求めた値であり、粒子の標本ではない。}
\label{fig:sm_densities}
\end{figure}


図\ref{fig:sm_vdf}は、同じ解の電子の法線速度分布である。分布は境界の Maxwell 分布をエネルギー保存で写したもので、
壁では太陽風電子は障壁を越えた高速の内向き成分だけ、光電子は外向きの成分だけを持つ。
下側では捕捉光電子が低速の両方向を占め、上側では反射電子が同じ速度領域を占める。
分布の切れ目は表\ref{tab:sm_populations}の集団の境界に対応する。

% Generated by tools/sheath_model_figures.py from vdf_A.csv; edit the generator, not this file.
\begin{figure}[htbp]
\centering
\begin{tikzpicture}[beach diagram]
\draw[black!10,line width=0.4pt] (0.683,0)--(0.683,4.4);
\draw[black!10,line width=0.4pt] (1.367,0)--(1.367,4.4);
\draw[black!10,line width=0.4pt] (2.05,0)--(2.05,4.4);
\draw[black!10,line width=0.4pt] (2.733,0)--(2.733,4.4);
\draw[black!10,line width=0.4pt] (3.417,0)--(3.417,4.4);
\draw[black!10,line width=0.4pt] (0,0)--(4.1,0);
\draw[black!10,line width=0.4pt] (0,0.983)--(4.1,0.983);
\draw[black!10,line width=0.4pt] (0,1.966)--(4.1,1.966);
\draw[black!10,line width=0.4pt] (0,2.948)--(4.1,2.948);
\draw[black!10,line width=0.4pt] (0,3.931)--(4.1,3.931);
\draw[line width=0.6pt] (0,0) rectangle (4.1,4.4);
\draw[line width=0.6pt] (0.683,0)--++(0,0.1);
\node[below,font=\scriptsize\sffamily] at (0.683,0) {$-4$};
\draw[line width=0.6pt] (1.367,0)--++(0,0.1);
\node[below,font=\scriptsize\sffamily] at (1.367,0) {$-2$};
\draw[line width=0.6pt] (2.05,0)--++(0,0.1);
\node[below,font=\scriptsize\sffamily] at (2.05,0) {$0$};
\draw[line width=0.6pt] (2.733,0)--++(0,0.1);
\node[below,font=\scriptsize\sffamily] at (2.733,0) {$2$};
\draw[line width=0.6pt] (3.417,0)--++(0,0.1);
\node[below,font=\scriptsize\sffamily] at (3.417,0) {$4$};
\draw[line width=0.6pt] (0,0)--++(0.1,0);
\node[left,font=\scriptsize\sffamily] at (0,0) {$10^{-3}$};
\draw[line width=0.6pt] (0,0.983)--++(0.1,0);
\node[left,font=\scriptsize\sffamily] at (0,0.983) {$10^{-2}$};
\draw[line width=0.6pt] (0,1.966)--++(0.1,0);
\node[left,font=\scriptsize\sffamily] at (0,1.966) {$10^{-1}$};
\draw[line width=0.6pt] (0,2.948)--++(0.1,0);
\node[left,font=\scriptsize\sffamily] at (0,2.948) {$10^{0}$};
\draw[line width=0.6pt] (0,3.931)--++(0.1,0);
\node[left,font=\scriptsize\sffamily] at (0,3.931) {$10^{1}$};
\node[font=\small\sffamily] at (2.05,-0.75) {$v_z$ [$10^6$ m/s]};
\node[rotate=90,font=\small\sffamily] at (-1.05,2.2) {$f(v_z)\,v_{te}/n_i$};
\node[above,font=\small\sffamily] at (2.05,4.4) {壁 $z=0$};
\begin{scope}
\clip (0,0) rectangle (4.1,4.4);
\draw[black!45,dashed,line width=0.5pt] (2.05,0)--(2.05,4.4);
\draw[BeachBlue,line width=1.1pt] plot coordinates {(-2.181,-12.605) (-2.175,-12.567) (-2.17,-12.528) (-2.165,-12.489) (-2.159,-12.451) (-2.154,-12.412) (-2.149,-12.374) (-2.143,-12.335) (-2.138,-12.297) (-2.133,-12.259) (-2.128,-12.22) (-2.122,-12.182) (-2.117,-12.144) (-2.112,-12.106) (-2.106,-12.068) (-2.101,-12.03) (-2.096,-11.992) (-2.091,-11.954) (-2.085,-11.916) (-2.08,-11.878) (-2.075,-11.84) (-2.069,-11.802) (-2.064,-11.764) (-2.059,-11.727) (-2.054,-11.689) (-2.048,-11.652) (-2.043,-11.614) (-2.038,-11.577) (-2.032,-11.539) (-2.027,-11.502) (-2.022,-11.465) (-2.017,-11.427) (-2.011,-11.39) (-2.006,-11.353) (-2.001,-11.316) (-1.995,-11.279) (-1.99,-11.242) (-1.985,-11.205) (-1.98,-11.168) (-1.974,-11.131) (-1.969,-11.094) (-1.964,-11.057) (-1.958,-11.02) (-1.953,-10.984) (-1.948,-10.947) (-1.943,-10.91) (-1.937,-10.874) (-1.932,-10.837) (-1.927,-10.801) (-1.921,-10.765) (-1.916,-10.728) (-1.911,-10.692) (-1.906,-10.656) (-1.9,-10.619) (-1.895,-10.583) (-1.89,-10.547) (-1.884,-10.511) (-1.879,-10.475) (-1.874,-10.439) (-1.869,-10.403) (-1.863,-10.367) (-1.858,-10.331) (-1.853,-10.296) (-1.847,-10.26) (-1.842,-10.224) (-1.837,-10.189) (-1.831,-10.153) (-1.826,-10.118) (-1.821,-10.082) (-1.816,-10.047) (-1.81,-10.011) (-1.805,-9.976) (-1.8,-9.941) (-1.794,-9.905) (-1.789,-9.87) (-1.784,-9.835) (-1.779,-9.8) (-1.773,-9.765) (-1.768,-9.73) (-1.763,-9.695) (-1.757,-9.66) (-1.752,-9.625) (-1.747,-9.59) (-1.742,-9.556) (-1.736,-9.521) (-1.731,-9.486) (-1.726,-9.452) (-1.72,-9.417) (-1.715,-9.383) (-1.71,-9.348) (-1.705,-9.314) (-1.699,-9.279) (-1.694,-9.245) (-1.689,-9.211) (-1.683,-9.176) (-1.678,-9.142) (-1.673,-9.108) (-1.668,-9.074) (-1.662,-9.04) (-1.657,-9.006) (-1.652,-8.972) (-1.646,-8.938) (-1.641,-8.904) (-1.636,-8.871) (-1.631,-8.837) (-1.625,-8.803) (-1.62,-8.77) (-1.615,-8.736) (-1.609,-8.702) (-1.604,-8.669) (-1.599,-8.635) (-1.594,-8.602) (-1.588,-8.569) (-1.583,-8.535) (-1.578,-8.502) (-1.572,-8.469) (-1.567,-8.436) (-1.562,-8.403) (-1.557,-8.37) (-1.551,-8.337) (-1.546,-8.304) (-1.541,-8.271) (-1.535,-8.238) (-1.53,-8.205) (-1.525,-8.172) (-1.519,-8.14) (-1.514,-8.107) (-1.509,-8.074) (-1.504,-8.042) (-1.498,-8.009) (-1.493,-7.977) (-1.488,-7.944) (-1.482,-7.912) (-1.477,-7.879) (-1.472,-7.847) (-1.467,-7.815) (-1.461,-7.783) (-1.456,-7.751) (-1.451,-7.719) (-1.445,-7.686) (-1.44,-7.654) (-1.435,-7.623) (-1.43,-7.591) (-1.424,-7.559) (-1.419,-7.527) (-1.414,-7.495) (-1.408,-7.464) (-1.403,-7.432) (-1.398,-7.4) (-1.393,-7.369) (-1.387,-7.337) (-1.382,-7.306) (-1.377,-7.274) (-1.371,-7.243) (-1.366,-7.212) (-1.361,-7.18) (-1.356,-7.149) (-1.35,-7.118) (-1.345,-7.087) (-1.34,-7.056) (-1.334,-7.025) (-1.329,-6.994) (-1.324,-6.963) (-1.319,-6.932) (-1.313,-6.901) (-1.308,-6.87) (-1.303,-6.84) (-1.297,-6.809) (-1.292,-6.778) (-1.287,-6.748) (-1.282,-6.717) (-1.276,-6.687) (-1.271,-6.656) (-1.266,-6.626) (-1.26,-6.595) (-1.255,-6.565) (-1.25,-6.535) (-1.245,-6.505) (-1.239,-6.474) (-1.234,-6.444) (-1.229,-6.414) (-1.223,-6.384) (-1.218,-6.354) (-1.213,-6.324) (-1.207,-6.295) (-1.202,-6.265) (-1.197,-6.235) (-1.192,-6.205) (-1.186,-6.176) (-1.181,-6.146) (-1.176,-6.116) (-1.17,-6.087) (-1.165,-6.057) (-1.16,-6.028) (-1.155,-5.999) (-1.149,-5.969) (-1.144,-5.94) (-1.139,-5.911) (-1.133,-5.882) (-1.128,-5.852) (-1.123,-5.823) (-1.118,-5.794) (-1.112,-5.765) (-1.107,-5.736) (-1.102,-5.707) (-1.096,-5.679) (-1.091,-5.65) (-1.086,-5.621) (-1.081,-5.592) (-1.075,-5.564) (-1.07,-5.535) (-1.065,-5.507) (-1.059,-5.478) (-1.054,-5.45) (-1.049,-5.421) (-1.044,-5.393) (-1.038,-5.364) (-1.033,-5.336) (-1.028,-5.308) (-1.022,-5.28) (-1.017,-5.252) (-1.012,-5.224) (-1.007,-5.196) (-1.001,-5.168) (-0.996,-5.14) (-0.991,-5.112) (-0.985,-5.084) (-0.98,-5.056) (-0.975,-5.028) (-0.97,-5.001) (-0.964,-4.973) (-0.959,-4.946) (-0.954,-4.918) (-0.948,-4.891) (-0.943,-4.863) (-0.938,-4.836) (-0.933,-4.808) (-0.927,-4.781) (-0.922,-4.754) (-0.917,-4.727) (-0.911,-4.699) (-0.906,-4.672) (-0.901,-4.645) (-0.895,-4.618) (-0.89,-4.591) (-0.885,-4.564) (-0.88,-4.538) (-0.874,-4.511) (-0.869,-4.484) (-0.864,-4.457) (-0.858,-4.431) (-0.853,-4.404) (-0.848,-4.377) (-0.843,-4.351) (-0.837,-4.324) (-0.832,-4.298) (-0.827,-4.272) (-0.821,-4.245) (-0.816,-4.219) (-0.811,-4.193) (-0.806,-4.167) (-0.8,-4.14) (-0.795,-4.114) (-0.79,-4.088) (-0.784,-4.062) (-0.779,-4.036) (-0.774,-4.01) (-0.769,-3.985) (-0.763,-3.959) (-0.758,-3.933) (-0.753,-3.907) (-0.747,-3.882) (-0.742,-3.856) (-0.737,-3.831) (-0.732,-3.805) (-0.726,-3.78) (-0.721,-3.754) (-0.716,-3.729) (-0.71,-3.704) (-0.705,-3.678) (-0.7,-3.653) (-0.695,-3.628) (-0.689,-3.603) (-0.684,-3.578) (-0.679,-3.553) (-0.673,-3.528) (-0.668,-3.503) (-0.663,-3.478) (-0.658,-3.453) (-0.652,-3.428) (-0.647,-3.404) (-0.642,-3.379) (-0.636,-3.354) (-0.631,-3.33) (-0.626,-3.305) (-0.621,-3.281) (-0.615,-3.256) (-0.61,-3.232) (-0.605,-3.207) (-0.599,-3.183) (-0.594,-3.159) (-0.589,-3.135) (-0.583,-3.111) (-0.578,-3.086) (-0.573,-3.062) (-0.568,-3.038) (-0.562,-3.014) (-0.557,-2.991) (-0.552,-2.967) (-0.546,-2.943) (-0.541,-2.919) (-0.536,-2.895) (-0.531,-2.872) (-0.525,-2.848) (-0.52,-2.825) (-0.515,-2.801) (-0.509,-2.778) (-0.504,-2.754) (-0.499,-2.731) (-0.494,-2.707) (-0.488,-2.684) (-0.483,-2.661) (-0.478,-2.638) (-0.472,-2.615) (-0.467,-2.592) (-0.462,-2.569) (-0.457,-2.546) (-0.451,-2.523) (-0.446,-2.5) (-0.441,-2.477) (-0.435,-2.454) (-0.43,-2.431) (-0.425,-2.409) (-0.42,-2.386) (-0.414,-2.363) (-0.409,-2.341) (-0.404,-2.318) (-0.398,-2.296) (-0.393,-2.273) (-0.388,-2.251) (-0.383,-2.229) (-0.377,-2.207) (-0.372,-2.184) (-0.367,-2.162) (-0.361,-2.14) (-0.356,-2.118) (-0.351,-2.096) (-0.346,-2.074) (-0.34,-2.052) (-0.335,-2.03) (-0.33,-2.008) (-0.324,-1.987) (-0.319,-1.965) (-0.314,-1.943) (-0.309,-1.922) (-0.303,-1.9) (-0.298,-1.878) (-0.293,-1.857) (-0.287,-1.836) (-0.282,-1.814) (-0.277,-1.793) (-0.271,-1.772) (-0.266,-1.75) (-0.261,-1.729) (-0.256,-1.708) (-0.25,-1.687) (-0.245,-1.666) (-0.24,-1.645) (-0.234,-1.624) (-0.229,-1.603) (-0.224,-1.582) (-0.219,-1.561) (-0.213,-1.541) (-0.208,-1.52) (-0.203,-1.499) (-0.197,-1.479) (-0.192,-1.458) (-0.187,-1.437) (-0.182,-1.417) (-0.176,-1.397) (-0.171,-1.376) (-0.166,-1.356) (-0.16,-1.336) (-0.155,-1.315) (-0.15,-1.295) (-0.145,-1.275) (-0.139,-1.255) (-0.134,-1.235) (-0.129,-1.215) (-0.123,-1.195) (-0.118,-1.175) (-0.113,-1.155) (-0.108,-1.136) (-0.102,-1.116) (-0.097,-1.096) (-0.092,-1.076) (-0.086,-1.057) (-0.081,-1.037) (-0.076,-1.018) (-0.071,-0.998) (-0.065,-0.979) (-0.06,-0.96) (-0.055,-0.94) (-0.049,-0.921) (-0.044,-0.902) (-0.039,-0.883) (-0.034,-0.864) (-0.028,-0.845) (-0.023,-0.826) (-0.018,-0.807) (-0.012,-0.788) (-0.007,-0.769) (-0.002,-0.75) (0.003,-0.731) (0.009,-0.712) (0.014,-0.694) (0.019,-0.675) (0.025,-0.657) (0.03,-0.638) (0.035,-0.62) (0.041,-0.601) (0.046,-0.583) (0.051,-0.564) (0.056,-0.546) (0.062,-0.528) (0.067,-0.51) (0.072,-0.492) (0.078,-0.473) (0.083,-0.455) (0.088,-0.437) (0.093,-0.419) (0.099,-0.402) (0.104,-0.384) (0.109,-0.366) (0.115,-0.348) (0.12,-0.33) (0.125,-0.313) (0.13,-0.295) (0.136,-0.278) (0.141,-0.26) (0.146,-0.243) (0.152,-0.225) (0.157,-0.208) (0.162,-0.191) (0.167,-0.173) (0.173,-0.156) (0.178,-0.139) (0.183,-0.122) (0.189,-0.105) (0.194,-0.088) (0.199,-0.071) (0.204,-0.054) (0.21,-0.037) (0.215,-0.02) (0.22,-0.003) (0.226,0.014) (0.231,0.03) (0.236,0.047) (0.241,0.063) (0.247,0.08) (0.252,0.096) (0.257,0.113) (0.263,0.129) (0.268,0.146) (0.273,0.162) (0.278,0.178) (0.284,0.194) (0.289,0.211) (0.294,0.227) (0.3,0.243) (0.305,0.259) (0.31,0.275) (0.315,0.291) (0.321,0.306) (0.326,0.322) (0.331,0.338) (0.337,0.354) (0.342,0.369) (0.347,0.385) (0.353,0.401) (0.358,0.416) (0.363,0.432) (0.368,0.447) (0.374,0.462) (0.379,0.478) (0.384,0.493) (0.39,0.508) (0.395,0.524) (0.4,0.539) (0.405,0.554) (0.411,0.569) (0.416,0.584) (0.421,0.599) (0.427,0.614) (0.432,0.628) (0.437,0.643) (0.442,0.658) (0.448,0.673) (0.453,0.687) (0.458,0.702) (0.464,0.717) (0.469,0.731) (0.474,0.746) (0.479,0.76) (0.485,0.774) (0.49,0.789) (0.495,0.803) (0.501,0.817) (0.506,0.831) (0.511,0.845) (0.516,0.859) (0.527,0.887) (0.538,0.915) (0.548,0.943) (0.559,0.97) (0.569,0.998) (0.58,1.025) (0.59,1.051) (0.601,1.078) (0.612,1.104) (0.622,1.131) (0.633,1.157) (0.643,1.183) (0.654,1.208) (0.665,1.234) (0.675,1.259) (0.686,1.284) (0.696,1.309) (0.707,1.334) (0.717,1.358) (0.728,1.383) (0.739,1.407) (0.749,1.431) (0.76,1.454) (0.77,1.478) (0.781,1.501) (0.791,1.524) (0.802,1.547) (0.813,1.57) (0.823,1.593) (0.834,1.615) (0.844,1.637) (0.855,1.659) (0.865,1.681) (0.876,1.703) (0.887,1.724) (0.897,1.745) (0.908,1.766) (0.918,1.787) (0.929,1.808) (0.939,1.828) (0.95,1.849) (0.961,1.869) (0.971,1.888) (0.982,1.908) (0.992,1.928) (1.003,1.947) (1.014,1.966) (1.024,1.985) (1.035,2.004) (1.045,2.022) (1.056,2.04) (1.066,2.059) (1.077,2.076) (1.088,2.094) (1.098,2.112) (1.109,2.129) (1.119,2.146) (1.13,2.163) (1.14,2.18) (1.151,2.196) (1.162,2.213) (1.172,2.229) (1.183,2.245) (1.193,2.261) (1.204,2.276) (1.214,2.292) (1.225,2.307) (1.236,2.322) (1.246,2.337) (1.257,2.352) (1.267,2.366) (1.278,2.38) (1.289,2.394) (1.299,2.408) (1.31,2.422) (1.32,2.435) (1.331,2.448) (1.341,2.462) (1.352,2.474) (1.363,2.487) (1.373,2.5) (1.384,2.512) (1.394,2.524) (1.405,2.536) (1.415,2.548) (1.426,2.559) (1.437,2.571) (1.447,2.582) (1.458,2.593) (1.468,2.603) (1.484,2.619) (1.5,2.635) (1.516,2.649) (1.532,2.664) (1.548,2.678) (1.563,2.691) (1.579,2.705) (1.595,2.717) (1.611,2.73) (1.627,2.741) (1.632,2.745)};
\draw[BeachOrange,line width=1.1pt] plot coordinates {(2.465,3.044) (2.469,3.026) (2.474,3.008) (2.478,2.99) (2.483,2.971) (2.487,2.953) (2.492,2.934) (2.496,2.915) (2.501,2.896) (2.505,2.877) (2.51,2.857) (2.514,2.838) (2.519,2.818) (2.523,2.798) (2.528,2.777) (2.532,2.757) (2.537,2.736) (2.541,2.715) (2.546,2.694) (2.55,2.673) (2.555,2.652) (2.559,2.63) (2.564,2.608) (2.568,2.586) (2.573,2.564) (2.577,2.542) (2.582,2.519) (2.587,2.496) (2.591,2.474) (2.596,2.45) (2.6,2.427) (2.605,2.404) (2.609,2.38) (2.614,2.356) (2.618,2.332) (2.623,2.307) (2.627,2.283) (2.632,2.258) (2.636,2.233) (2.641,2.208) (2.645,2.183) (2.65,2.158) (2.654,2.132) (2.659,2.106) (2.663,2.08) (2.668,2.054) (2.672,2.028) (2.677,2.001) (2.681,1.974) (2.686,1.947) (2.69,1.92) (2.695,1.893) (2.699,1.865) (2.704,1.837) (2.708,1.809) (2.713,1.781) (2.717,1.753) (2.722,1.724) (2.726,1.696) (2.731,1.667) (2.735,1.638) (2.74,1.608) (2.744,1.579) (2.749,1.549) (2.751,1.534) (2.753,1.519) (2.756,1.504) (2.758,1.489) (2.76,1.474) (2.762,1.459) (2.765,1.444) (2.767,1.429) (2.769,1.413) (2.771,1.398) (2.774,1.383) (2.776,1.367) (2.778,1.352) (2.78,1.336) (2.783,1.321) (2.785,1.305) (2.787,1.289) (2.789,1.273) (2.792,1.258) (2.794,1.242) (2.796,1.226) (2.798,1.21) (2.801,1.194) (2.803,1.178) (2.805,1.162) (2.807,1.146) (2.81,1.13) (2.812,1.114) (2.814,1.097) (2.816,1.081) (2.819,1.065) (2.821,1.048) (2.823,1.032) (2.825,1.015) (2.828,0.999) (2.83,0.982) (2.832,0.966) (2.834,0.949) (2.837,0.932) (2.839,0.915) (2.841,0.899) (2.843,0.882) (2.846,0.865) (2.848,0.848) (2.85,0.831) (2.853,0.814) (2.855,0.797) (2.857,0.779) (2.859,0.762) (2.862,0.745) (2.864,0.728) (2.866,0.71) (2.868,0.693) (2.871,0.675) (2.873,0.658) (2.875,0.64) (2.877,0.623) (2.88,0.605) (2.882,0.587) (2.884,0.57) (2.886,0.552) (2.889,0.534) (2.891,0.516) (2.893,0.498) (2.895,0.48) (2.898,0.462) (2.9,0.444) (2.902,0.426) (2.904,0.408) (2.907,0.39) (2.909,0.371) (2.911,0.353) (2.913,0.335) (2.916,0.316) (2.918,0.298) (2.92,0.279) (2.922,0.261) (2.925,0.242) (2.927,0.223) (2.929,0.205) (2.931,0.186) (2.934,0.167) (2.936,0.148) (2.938,0.129) (2.94,0.111) (2.943,0.092) (2.945,0.072) (2.947,0.053) (2.949,0.034) (2.952,0.015) (2.954,-0.004) (2.956,-0.023) (2.958,-0.043) (2.961,-0.062) (2.963,-0.082) (2.965,-0.101) (2.967,-0.121) (2.97,-0.14) (2.972,-0.16) (2.974,-0.179) (2.976,-0.199) (2.979,-0.219) (2.981,-0.239) (2.983,-0.259) (2.986,-0.278) (2.988,-0.298) (2.99,-0.318) (2.992,-0.338) (2.995,-0.359) (2.997,-0.379) (2.999,-0.399) (3.001,-0.419) (3.004,-0.439) (3.006,-0.46) (3.008,-0.48) (3.01,-0.501) (3.013,-0.521) (3.015,-0.542) (3.017,-0.562) (3.019,-0.583) (3.022,-0.603) (3.024,-0.624) (3.026,-0.645) (3.028,-0.666) (3.031,-0.687) (3.033,-0.707) (3.035,-0.728) (3.037,-0.749) (3.04,-0.77) (3.042,-0.792) (3.044,-0.813) (3.046,-0.834) (3.049,-0.855) (3.051,-0.876) (3.053,-0.898) (3.055,-0.919) (3.058,-0.941) (3.06,-0.962) (3.062,-0.984) (3.064,-1.005) (3.067,-1.027) (3.069,-1.049) (3.071,-1.07) (3.073,-1.092) (3.076,-1.114) (3.078,-1.136) (3.08,-1.158) (3.082,-1.18) (3.085,-1.202) (3.087,-1.224) (3.089,-1.246) (3.091,-1.268) (3.094,-1.29) (3.096,-1.312) (3.098,-1.335) (3.1,-1.357) (3.103,-1.38) (3.105,-1.402) (3.107,-1.424) (3.109,-1.447) (3.112,-1.47) (3.114,-1.492) (3.116,-1.515) (3.119,-1.538) (3.121,-1.56) (3.123,-1.583) (3.125,-1.606) (3.128,-1.629) (3.13,-1.652) (3.132,-1.675) (3.134,-1.698) (3.137,-1.721) (3.139,-1.744) (3.141,-1.768) (3.143,-1.791) (3.146,-1.814) (3.148,-1.838) (3.15,-1.861) (3.152,-1.884) (3.155,-1.908) (3.157,-1.932) (3.159,-1.955) (3.161,-1.979) (3.164,-2.002) (3.166,-2.026) (3.168,-2.05) (3.17,-2.074) (3.173,-2.098) (3.175,-2.122) (3.177,-2.146) (3.179,-2.17) (3.182,-2.194) (3.184,-2.218) (3.186,-2.242) (3.188,-2.266) (3.191,-2.291) (3.193,-2.315) (3.195,-2.339) (3.197,-2.364) (3.2,-2.388) (3.202,-2.413) (3.204,-2.437) (3.206,-2.462) (3.209,-2.486) (3.211,-2.511) (3.213,-2.536) (3.215,-2.561) (3.218,-2.586) (3.22,-2.61) (3.222,-2.635) (3.224,-2.66) (3.227,-2.685) (3.229,-2.711) (3.231,-2.736) (3.233,-2.761) (3.236,-2.786) (3.238,-2.811) (3.24,-2.837) (3.242,-2.862) (3.245,-2.888) (3.247,-2.913) (3.249,-2.939) (3.252,-2.964) (3.254,-2.99) (3.256,-3.015) (3.258,-3.041) (3.261,-3.067) (3.263,-3.093) (3.265,-3.119) (3.267,-3.144) (3.27,-3.17) (3.272,-3.196) (3.274,-3.222) (3.276,-3.249) (3.279,-3.275) (3.281,-3.301) (3.283,-3.327) (3.285,-3.353) (3.288,-3.38) (3.29,-3.406) (3.292,-3.433) (3.294,-3.459) (3.297,-3.486) (3.299,-3.512) (3.301,-3.539) (3.303,-3.565) (3.306,-3.592) (3.308,-3.619) (3.31,-3.646) (3.312,-3.673) (3.315,-3.7) (3.317,-3.726) (3.319,-3.753) (3.321,-3.781) (3.324,-3.808) (3.326,-3.835) (3.328,-3.862) (3.33,-3.889) (3.333,-3.917) (3.335,-3.944) (3.337,-3.971) (3.339,-3.999) (3.342,-4.026) (3.344,-4.054) (3.346,-4.081) (3.348,-4.109) (3.351,-4.137) (3.353,-4.164) (3.355,-4.192) (3.357,-4.22) (3.36,-4.248) (3.362,-4.276) (3.364,-4.304) (3.366,-4.332) (3.369,-4.36) (3.371,-4.388) (3.373,-4.416) (3.375,-4.444) (3.378,-4.473) (3.38,-4.501) (3.382,-4.529) (3.385,-4.558) (3.387,-4.586) (3.389,-4.615) (3.391,-4.643) (3.394,-4.672) (3.396,-4.7) (3.398,-4.729) (3.4,-4.758) (3.403,-4.787) (3.405,-4.815) (3.407,-4.844) (3.409,-4.873) (3.412,-4.902) (3.414,-4.931) (3.416,-4.96) (3.418,-4.989) (3.421,-5.019) (3.423,-5.048) (3.425,-5.077) (3.427,-5.106) (3.43,-5.136) (3.432,-5.165) (3.434,-5.195) (3.436,-5.224) (3.439,-5.254) (3.441,-5.283) (3.443,-5.313) (3.445,-5.343) (3.448,-5.372) (3.45,-5.402) (3.452,-5.432) (3.454,-5.462) (3.457,-5.492) (3.459,-5.522) (3.461,-5.552) (3.463,-5.582) (3.466,-5.612) (3.468,-5.642) (3.47,-5.672) (3.472,-5.703) (3.475,-5.733) (3.477,-5.763) (3.479,-5.794) (3.481,-5.824) (3.484,-5.855) (3.486,-5.885) (3.488,-5.916) (3.49,-5.947) (3.493,-5.977) (3.495,-6.008) (3.497,-6.039) (3.499,-6.07) (3.502,-6.101) (3.504,-6.132) (3.506,-6.163) (3.508,-6.194) (3.511,-6.225) (3.513,-6.256) (3.515,-6.287) (3.518,-6.318) (3.52,-6.35) (3.522,-6.381) (3.524,-6.412) (3.527,-6.444) (3.529,-6.475) (3.531,-6.507) (3.533,-6.538) (3.536,-6.57) (3.538,-6.602) (3.54,-6.633) (3.542,-6.665) (3.545,-6.697) (3.547,-6.729) (3.549,-6.761) (3.551,-6.793) (3.554,-6.825) (3.556,-6.857) (3.558,-6.889) (3.56,-6.921) (3.563,-6.953) (3.565,-6.985) (3.567,-7.018) (3.569,-7.05) (3.572,-7.082) (3.574,-7.115) (3.576,-7.147) (3.578,-7.18) (3.581,-7.212) (3.583,-7.245) (3.585,-7.278) (3.587,-7.31) (3.59,-7.343) (3.592,-7.376) (3.594,-7.409) (3.596,-7.442) (3.599,-7.475) (3.601,-7.508) (3.603,-7.541) (3.605,-7.574) (3.608,-7.607) (3.61,-7.64) (3.612,-7.674) (3.614,-7.707) (3.617,-7.74) (3.619,-7.774) (3.621,-7.807) (3.623,-7.841) (3.626,-7.874) (3.628,-7.908) (3.63,-7.941) (3.632,-7.975) (3.635,-8.009) (3.637,-8.042) (3.639,-8.076) (3.641,-8.11) (3.644,-8.144) (3.646,-8.178) (3.648,-8.212) (3.651,-8.246) (3.653,-8.28) (3.655,-8.314) (3.657,-8.349) (3.66,-8.383) (3.662,-8.417) (3.664,-8.452) (3.666,-8.486) (3.669,-8.52) (3.671,-8.555) (3.673,-8.589) (3.675,-8.624) (3.678,-8.659) (3.68,-8.693) (3.682,-8.728) (3.684,-8.763) (3.687,-8.798) (3.689,-8.833) (3.691,-8.868) (3.693,-8.903) (3.696,-8.938) (3.698,-8.973) (3.7,-9.008) (3.702,-9.043) (3.705,-9.078) (3.707,-9.113) (3.709,-9.149) (3.711,-9.184) (3.714,-9.22) (3.716,-9.255) (3.718,-9.291) (3.72,-9.326) (3.723,-9.362) (3.725,-9.397) (3.727,-9.433) (3.729,-9.469) (3.732,-9.505) (3.734,-9.54) (3.736,-9.576) (3.738,-9.612) (3.741,-9.648) (3.743,-9.684) (3.745,-9.72) (3.747,-9.756) (3.75,-9.793) (3.752,-9.829) (3.754,-9.865) (3.756,-9.902) (3.759,-9.938) (3.761,-9.974) (3.763,-10.011) (3.765,-10.047) (3.768,-10.084) (3.77,-10.12) (3.772,-10.157) (3.774,-10.194) (3.777,-10.231) (3.779,-10.267) (3.781,-10.304) (3.784,-10.341) (3.786,-10.378) (3.788,-10.415) (3.79,-10.452) (3.793,-10.489) (3.795,-10.526) (3.797,-10.564) (3.799,-10.601) (3.802,-10.638) (3.804,-10.675) (3.806,-10.713) (3.808,-10.75) (3.811,-10.788) (3.813,-10.825) (3.815,-10.863) (3.817,-10.9) (3.82,-10.938) (3.822,-10.976) (3.824,-11.014) (3.826,-11.051) (3.829,-11.089) (3.831,-11.127) (3.833,-11.165) (3.835,-11.203) (3.838,-11.241) (3.84,-11.279) (3.842,-11.317) (3.844,-11.356) (3.847,-11.394) (3.849,-11.432) (3.851,-11.47) (3.853,-11.509)};
\draw[BeachRed,line width=1.1pt] plot coordinates {(2.05,3.856) (2.066,3.855) (2.082,3.852) (2.097,3.846) (2.111,3.839) (2.124,3.83) (2.138,3.82) (2.149,3.81) (2.16,3.799) (2.172,3.786) (2.183,3.773) (2.194,3.758) (2.203,3.745) (2.212,3.732) (2.221,3.718) (2.23,3.703) (2.239,3.687) (2.248,3.67) (2.257,3.653) (2.264,3.64) (2.271,3.626) (2.278,3.611) (2.284,3.597) (2.291,3.582) (2.298,3.566) (2.305,3.55) (2.311,3.533) (2.318,3.516) (2.325,3.499) (2.332,3.481) (2.339,3.463) (2.345,3.444) (2.352,3.425) (2.359,3.406) (2.366,3.386) (2.372,3.365) (2.379,3.345) (2.386,3.323) (2.39,3.309) (2.395,3.294) (2.399,3.28) (2.404,3.265) (2.408,3.249) (2.413,3.234) (2.417,3.219) (2.422,3.203) (2.426,3.187) (2.431,3.171) (2.435,3.154) (2.44,3.138) (2.444,3.121) (2.449,3.104) (2.454,3.087) (2.458,3.07) (2.463,3.052)};
\draw[BeachRed,line width=1.1pt] plot coordinates {(1.637,3.052) (1.642,3.07) (1.646,3.087) (1.651,3.104) (1.656,3.121) (1.66,3.138) (1.665,3.154) (1.669,3.171) (1.674,3.187) (1.678,3.203) (1.683,3.219) (1.687,3.234) (1.692,3.249) (1.696,3.265) (1.701,3.28) (1.705,3.294) (1.71,3.309) (1.714,3.323) (1.721,3.345) (1.728,3.365) (1.734,3.386) (1.741,3.406) (1.748,3.425) (1.755,3.444) (1.761,3.463) (1.768,3.481) (1.775,3.499) (1.782,3.516) (1.789,3.533) (1.795,3.55) (1.802,3.566) (1.809,3.582) (1.816,3.597) (1.822,3.611) (1.829,3.626) (1.836,3.64) (1.843,3.653) (1.852,3.67) (1.861,3.687) (1.87,3.703) (1.879,3.718) (1.888,3.732) (1.897,3.745) (1.906,3.758) (1.917,3.773) (1.928,3.786) (1.94,3.799) (1.951,3.81) (1.962,3.82) (1.976,3.83) (1.989,3.839) (2.003,3.846) (2.018,3.852) (2.034,3.855) (2.048,3.856)};
\end{scope}
\draw[black!10,line width=0.4pt] (5.283,0)--(5.283,4.4);
\draw[black!10,line width=0.4pt] (5.967,0)--(5.967,4.4);
\draw[black!10,line width=0.4pt] (6.65,0)--(6.65,4.4);
\draw[black!10,line width=0.4pt] (7.333,0)--(7.333,4.4);
\draw[black!10,line width=0.4pt] (8.017,0)--(8.017,4.4);
\draw[black!10,line width=0.4pt] (4.6,0)--(8.7,0);
\draw[black!10,line width=0.4pt] (4.6,0.983)--(8.7,0.983);
\draw[black!10,line width=0.4pt] (4.6,1.966)--(8.7,1.966);
\draw[black!10,line width=0.4pt] (4.6,2.948)--(8.7,2.948);
\draw[black!10,line width=0.4pt] (4.6,3.931)--(8.7,3.931);
\draw[line width=0.6pt] (4.6,0) rectangle (8.7,4.4);
\draw[line width=0.6pt] (5.283,0)--++(0,0.1);
\node[below,font=\scriptsize\sffamily] at (5.283,0) {$-4$};
\draw[line width=0.6pt] (5.967,0)--++(0,0.1);
\node[below,font=\scriptsize\sffamily] at (5.967,0) {$-2$};
\draw[line width=0.6pt] (6.65,0)--++(0,0.1);
\node[below,font=\scriptsize\sffamily] at (6.65,0) {$0$};
\draw[line width=0.6pt] (7.333,0)--++(0,0.1);
\node[below,font=\scriptsize\sffamily] at (7.333,0) {$2$};
\draw[line width=0.6pt] (8.017,0)--++(0,0.1);
\node[below,font=\scriptsize\sffamily] at (8.017,0) {$4$};
\draw[line width=0.6pt] (4.6,0)--++(0.1,0);
\draw[line width=0.6pt] (4.6,0.983)--++(0.1,0);
\draw[line width=0.6pt] (4.6,1.966)--++(0.1,0);
\draw[line width=0.6pt] (4.6,2.948)--++(0.1,0);
\draw[line width=0.6pt] (4.6,3.931)--++(0.1,0);
\node[font=\small\sffamily] at (6.65,-0.75) {$v_z$ [$10^6$ m/s]};
\node[above,font=\small\sffamily] at (6.65,4.4) {下側 $z=7.4$ m};
\begin{scope}
\clip (4.6,0) rectangle (8.7,4.4);
\draw[black!45,dashed,line width=0.5pt] (6.65,0)--(6.65,4.4);
\draw[BeachBlue,line width=1.1pt] plot coordinates {(2.438,-12.605) (2.444,-12.567) (2.449,-12.529) (2.454,-12.49) (2.459,-12.452) (2.465,-12.414) (2.47,-12.376) (2.475,-12.338) (2.48,-12.3) (2.486,-12.262) (2.491,-12.224) (2.496,-12.186) (2.501,-12.148) (2.507,-12.11) (2.512,-12.072) (2.517,-12.035) (2.522,-11.997) (2.528,-11.959) (2.533,-11.922) (2.538,-11.884) (2.544,-11.847) (2.549,-11.809) (2.554,-11.772) (2.559,-11.735) (2.565,-11.697) (2.57,-11.66) (2.575,-11.623) (2.58,-11.586) (2.586,-11.549) (2.591,-11.512) (2.596,-11.475) (2.601,-11.438) (2.607,-11.401) (2.612,-11.364) (2.617,-11.327) (2.623,-11.29) (2.628,-11.254) (2.633,-11.217) (2.638,-11.18) (2.644,-11.144) (2.649,-11.107) (2.654,-11.071) (2.659,-11.034) (2.665,-10.998) (2.67,-10.962) (2.675,-10.925) (2.68,-10.889) (2.686,-10.853) (2.691,-10.817) (2.696,-10.781) (2.701,-10.745) (2.707,-10.709) (2.712,-10.673) (2.717,-10.637) (2.723,-10.601) (2.728,-10.565) (2.733,-10.53) (2.738,-10.494) (2.744,-10.458) (2.749,-10.423) (2.754,-10.387) (2.759,-10.352) (2.765,-10.316) (2.77,-10.281) (2.775,-10.245) (2.78,-10.21) (2.786,-10.175) (2.791,-10.14) (2.796,-10.104) (2.802,-10.069) (2.807,-10.034) (2.812,-9.999) (2.817,-9.964) (2.823,-9.929) (2.828,-9.894) (2.833,-9.86) (2.838,-9.825) (2.844,-9.79) (2.849,-9.755) (2.854,-9.721) (2.859,-9.686) (2.865,-9.652) (2.87,-9.617) (2.875,-9.583) (2.88,-9.548) (2.886,-9.514) (2.891,-9.48) (2.896,-9.445) (2.902,-9.411) (2.907,-9.377) (2.912,-9.343) (2.917,-9.309) (2.923,-9.275) (2.928,-9.241) (2.933,-9.207) (2.938,-9.173) (2.944,-9.139) (2.949,-9.105) (2.954,-9.072) (2.959,-9.038) (2.965,-9.004) (2.97,-8.971) (2.975,-8.937) (2.981,-8.904) (2.986,-8.87) (2.991,-8.837) (2.996,-8.803) (3.002,-8.77) (3.007,-8.737) (3.012,-8.704) (3.017,-8.671) (3.023,-8.637) (3.028,-8.604) (3.033,-8.571) (3.038,-8.538) (3.044,-8.506) (3.049,-8.473) (3.054,-8.44) (3.059,-8.407) (3.065,-8.374) (3.07,-8.342) (3.075,-8.309) (3.081,-8.276) (3.086,-8.244) (3.091,-8.211) (3.096,-8.179) (3.102,-8.147) (3.107,-8.114) (3.112,-8.082) (3.117,-8.05) (3.123,-8.018) (3.128,-7.985) (3.133,-7.953) (3.138,-7.921) (3.144,-7.889) (3.149,-7.857) (3.154,-7.825) (3.16,-7.794) (3.165,-7.762) (3.17,-7.73) (3.175,-7.698) (3.181,-7.667) (3.186,-7.635) (3.191,-7.603) (3.196,-7.572) (3.202,-7.54) (3.207,-7.509) (3.212,-7.478) (3.217,-7.446) (3.223,-7.415) (3.228,-7.384) (3.233,-7.353) (3.238,-7.321) (3.244,-7.29) (3.249,-7.259) (3.254,-7.228) (3.26,-7.197) (3.265,-7.167) (3.27,-7.136) (3.275,-7.105) (3.281,-7.074) (3.286,-7.043) (3.291,-7.013) (3.296,-6.982) (3.302,-6.952) (3.307,-6.921) (3.312,-6.891) (3.317,-6.86) (3.323,-6.83) (3.328,-6.799) (3.333,-6.769) (3.339,-6.739) (3.344,-6.709) (3.349,-6.679) (3.354,-6.649) (3.36,-6.619) (3.365,-6.589) (3.37,-6.559) (3.375,-6.529) (3.381,-6.499) (3.386,-6.469) (3.391,-6.439) (3.396,-6.41) (3.402,-6.38) (3.407,-6.35) (3.412,-6.321) (3.417,-6.291) (3.423,-6.262) (3.428,-6.232) (3.433,-6.203) (3.439,-6.174) (3.444,-6.145) (3.449,-6.115) (3.454,-6.086) (3.46,-6.057) (3.465,-6.028) (3.47,-5.999) (3.475,-5.97) (3.481,-5.941) (3.486,-5.912) (3.491,-5.883) (3.496,-5.855) (3.502,-5.826) (3.507,-5.797) (3.512,-5.768) (3.518,-5.74) (3.523,-5.711) (3.528,-5.683) (3.533,-5.654) (3.539,-5.626) (3.544,-5.598) (3.549,-5.569) (3.554,-5.541) (3.56,-5.513) (3.565,-5.485) (3.57,-5.457) (3.575,-5.429) (3.581,-5.401) (3.586,-5.373) (3.591,-5.345) (3.596,-5.317) (3.602,-5.289) (3.607,-5.261) (3.612,-5.233) (3.618,-5.206) (3.623,-5.178) (3.628,-5.151) (3.633,-5.123) (3.639,-5.096) (3.644,-5.068) (3.649,-5.041) (3.654,-5.013) (3.66,-4.986) (3.665,-4.959) (3.67,-4.932) (3.675,-4.904) (3.681,-4.877) (3.686,-4.85) (3.691,-4.823) (3.697,-4.796) (3.702,-4.769) (3.707,-4.743) (3.712,-4.716) (3.718,-4.689) (3.723,-4.662) (3.728,-4.636) (3.733,-4.609) (3.739,-4.582) (3.744,-4.556) (3.749,-4.529) (3.754,-4.503) (3.76,-4.477) (3.765,-4.45) (3.77,-4.424) (3.775,-4.398) (3.781,-4.371) (3.786,-4.345) (3.791,-4.319) (3.797,-4.293) (3.802,-4.267) (3.807,-4.241) (3.812,-4.215) (3.818,-4.189) (3.823,-4.164) (3.828,-4.138) (3.833,-4.112) (3.839,-4.087) (3.844,-4.061) (3.849,-4.035) (3.854,-4.01) (3.86,-3.984) (3.865,-3.959) (3.87,-3.934) (3.876,-3.908) (3.881,-3.883) (3.886,-3.858) (3.891,-3.833) (3.897,-3.807) (3.902,-3.782) (3.907,-3.757) (3.912,-3.732) (3.918,-3.707) (3.923,-3.682) (3.928,-3.658) (3.933,-3.633) (3.939,-3.608) (3.944,-3.583) (3.949,-3.559) (3.954,-3.534) (3.96,-3.51) (3.965,-3.485) (3.97,-3.461) (3.976,-3.436) (3.981,-3.412) (3.986,-3.387) (3.991,-3.363) (3.997,-3.339) (4.002,-3.315) (4.007,-3.291) (4.012,-3.267) (4.018,-3.243) (4.023,-3.219) (4.028,-3.195) (4.033,-3.171) (4.039,-3.147) (4.044,-3.123) (4.049,-3.099) (4.055,-3.076) (4.06,-3.052) (4.065,-3.028) (4.07,-3.005) (4.076,-2.981) (4.081,-2.958) (4.086,-2.934) (4.091,-2.911) (4.097,-2.888) (4.102,-2.865) (4.107,-2.841) (4.112,-2.818) (4.118,-2.795) (4.123,-2.772) (4.128,-2.749) (4.133,-2.726) (4.139,-2.703) (4.144,-2.68) (4.149,-2.657) (4.155,-2.635) (4.16,-2.612) (4.165,-2.589) (4.17,-2.567) (4.176,-2.544) (4.181,-2.521) (4.186,-2.499) (4.191,-2.476) (4.197,-2.454) (4.202,-2.432) (4.207,-2.409) (4.212,-2.387) (4.218,-2.365) (4.223,-2.343) (4.228,-2.321) (4.234,-2.299) (4.239,-2.277) (4.244,-2.255) (4.249,-2.233) (4.255,-2.211) (4.26,-2.189) (4.265,-2.167) (4.27,-2.146) (4.276,-2.124) (4.281,-2.102) (4.286,-2.081) (4.291,-2.059) (4.297,-2.038) (4.302,-2.016) (4.307,-1.995) (4.312,-1.973) (4.318,-1.952) (4.323,-1.931) (4.328,-1.91) (4.334,-1.889) (4.339,-1.867) (4.344,-1.846) (4.349,-1.825) (4.355,-1.804) (4.36,-1.783) (4.365,-1.763) (4.37,-1.742) (4.376,-1.721) (4.381,-1.7) (4.386,-1.68) (4.391,-1.659) (4.397,-1.638) (4.402,-1.618) (4.407,-1.597) (4.413,-1.577) (4.418,-1.557) (4.423,-1.536) (4.428,-1.516) (4.434,-1.496) (4.439,-1.476) (4.444,-1.455) (4.449,-1.435) (4.455,-1.415) (4.46,-1.395) (4.465,-1.375) (4.47,-1.355) (4.476,-1.336) (4.481,-1.316) (4.486,-1.296) (4.491,-1.276) (4.497,-1.257) (4.502,-1.237) (4.507,-1.217) (4.513,-1.198) (4.518,-1.179) (4.523,-1.159) (4.528,-1.14) (4.534,-1.12) (4.539,-1.101) (4.544,-1.082) (4.549,-1.063) (4.555,-1.044) (4.56,-1.024) (4.565,-1.005) (4.57,-0.986) (4.576,-0.968) (4.581,-0.949) (4.586,-0.93) (4.592,-0.911) (4.597,-0.892) (4.602,-0.874) (4.607,-0.855) (4.613,-0.836) (4.618,-0.818) (4.623,-0.799) (4.628,-0.781) (4.634,-0.762) (4.639,-0.744) (4.644,-0.726) (4.649,-0.707) (4.655,-0.689) (4.66,-0.671) (4.665,-0.653) (4.67,-0.635) (4.676,-0.617) (4.681,-0.599) (4.686,-0.581) (4.692,-0.563) (4.697,-0.545) (4.702,-0.527) (4.707,-0.51) (4.713,-0.492) (4.718,-0.474) (4.723,-0.457) (4.728,-0.439) (4.734,-0.422) (4.739,-0.404) (4.744,-0.387) (4.749,-0.369) (4.755,-0.352) (4.76,-0.335) (4.765,-0.318) (4.771,-0.3) (4.776,-0.283) (4.781,-0.266) (4.786,-0.249) (4.792,-0.232) (4.797,-0.215) (4.802,-0.198) (4.807,-0.182) (4.813,-0.165) (4.818,-0.148) (4.823,-0.131) (4.828,-0.115) (4.834,-0.098) (4.839,-0.082) (4.844,-0.065) (4.849,-0.049) (4.855,-0.032) (4.86,-0.016) (4.865,0) (4.871,0.017) (4.876,0.033) (4.881,0.049) (4.886,0.065) (4.892,0.081) (4.897,0.097) (4.902,0.113) (4.907,0.129) (4.913,0.145) (4.918,0.161) (4.923,0.176) (4.928,0.192) (4.934,0.208) (4.939,0.224) (4.944,0.239) (4.95,0.255) (4.955,0.27) (4.96,0.286) (4.965,0.301) (4.971,0.316) (4.976,0.332) (4.981,0.347) (4.986,0.362) (4.992,0.377) (4.997,0.392) (5.002,0.407) (5.007,0.422) (5.013,0.437) (5.018,0.452) (5.023,0.467) (5.028,0.482) (5.034,0.497) (5.039,0.511) (5.044,0.526) (5.05,0.541) (5.055,0.555) (5.06,0.57) (5.065,0.584) (5.071,0.599) (5.076,0.613) (5.081,0.627) (5.086,0.642) (5.092,0.656) (5.097,0.67) (5.102,0.684) (5.107,0.698) (5.118,0.726) (5.129,0.754) (5.139,0.782) (5.15,0.809) (5.16,0.837) (5.171,0.864) (5.181,0.891) (5.192,0.917) (5.202,0.944) (5.213,0.97) (5.223,0.996) (5.234,1.022) (5.244,1.048) (5.255,1.073) (5.265,1.099) (5.276,1.124) (5.286,1.149) (5.297,1.174) (5.308,1.198) (5.318,1.223) (5.329,1.247) (5.339,1.271) (5.35,1.295) (5.36,1.318) (5.371,1.342) (5.381,1.365) (5.392,1.388) (5.402,1.411) (5.413,1.433) (5.423,1.456) (5.434,1.478) (5.444,1.5) (5.455,1.522) (5.465,1.544) (5.476,1.565) (5.487,1.587) (5.497,1.608) (5.508,1.629) (5.518,1.65) (5.529,1.67) (5.539,1.69) (5.55,1.711) (5.56,1.731) (5.571,1.75) (5.581,1.77) (5.592,1.789) (5.602,1.809) (5.613,1.828) (5.623,1.846) (5.634,1.865) (5.644,1.883) (5.655,1.902) (5.666,1.92) (5.676,1.938) (5.687,1.955) (5.697,1.973) (5.708,1.99) (5.718,2.007) (5.729,2.024) (5.739,2.041) (5.75,2.057) (5.76,2.074) (5.771,2.09) (5.781,2.106) (5.792,2.121) (5.802,2.137) (5.813,2.152) (5.823,2.168) (5.834,2.183) (5.845,2.197) (5.855,2.212) (5.866,2.226) (5.876,2.241) (5.887,2.255) (5.897,2.268) (5.908,2.282) (5.918,2.295) (5.929,2.309) (5.939,2.322) (5.95,2.335) (5.96,2.347) (5.971,2.36) (5.981,2.372) (5.992,2.384) (6.002,2.396) (6.013,2.408) (6.024,2.419) (6.034,2.431) (6.045,2.442) (6.055,2.453) (6.066,2.463) (6.081,2.479) (6.097,2.495) (6.113,2.51) (6.129,2.524) (6.145,2.538) (6.16,2.552) (6.176,2.565) (6.192,2.578) (6.208,2.59) (6.224,2.602) (6.239,2.613) (6.255,2.624) (6.271,2.635) (6.287,2.645) (6.303,2.655) (6.318,2.664) (6.334,2.673) (6.35,2.681) (6.366,2.689) (6.382,2.697) (6.397,2.704) (6.413,2.711) (6.429,2.717) (6.445,2.723) (6.46,2.728) (6.476,2.733) (6.492,2.738) (6.508,2.742) (6.524,2.745) (6.534,2.748)};
\draw[BeachOrange,line width=1.1pt] plot coordinates {(6.765,3.046) (6.776,3.033) (6.788,3.019) (6.797,3.007) (6.806,2.994) (6.815,2.98) (6.824,2.966) (6.833,2.95) (6.842,2.934) (6.851,2.918) (6.86,2.9) (6.866,2.887) (6.873,2.873) (6.88,2.858) (6.887,2.843) (6.893,2.828) (6.9,2.812) (6.907,2.796) (6.914,2.779) (6.921,2.762) (6.927,2.745) (6.934,2.727) (6.941,2.708) (6.948,2.69) (6.954,2.67) (6.961,2.651) (6.968,2.631) (6.975,2.61) (6.981,2.589) (6.988,2.568) (6.993,2.553) (6.997,2.539) (7.002,2.524) (7.006,2.509) (7.011,2.493) (7.015,2.478) (7.02,2.462) (7.024,2.446) (7.029,2.43) (7.033,2.414) (7.038,2.398) (7.042,2.381) (7.047,2.364) (7.051,2.347) (7.056,2.33) (7.06,2.313) (7.065,2.295) (7.069,2.277) (7.074,2.259) (7.078,2.241) (7.083,2.223) (7.087,2.204) (7.092,2.186) (7.096,2.167) (7.101,2.148) (7.105,2.128) (7.11,2.109) (7.114,2.089) (7.119,2.069) (7.123,2.049) (7.128,2.029) (7.132,2.008) (7.137,1.988) (7.141,1.967) (7.146,1.946) (7.15,1.925) (7.155,1.903) (7.159,1.882) (7.164,1.86) (7.168,1.838) (7.173,1.816) (7.177,1.793) (7.182,1.771) (7.187,1.748) (7.191,1.725) (7.196,1.702) (7.2,1.679) (7.205,1.655) (7.209,1.631) (7.214,1.607) (7.218,1.583) (7.223,1.559) (7.227,1.535) (7.232,1.51) (7.236,1.485) (7.241,1.46) (7.245,1.435) (7.25,1.409) (7.254,1.384) (7.259,1.358) (7.263,1.332) (7.268,1.306) (7.272,1.279) (7.277,1.253) (7.281,1.226) (7.286,1.199) (7.29,1.172) (7.295,1.144) (7.299,1.117) (7.304,1.089) (7.308,1.061) (7.313,1.033) (7.317,1.004) (7.322,0.976) (7.326,0.947) (7.331,0.918) (7.335,0.889) (7.34,0.86) (7.344,0.83) (7.349,0.801) (7.351,0.786) (7.353,0.771) (7.356,0.756) (7.358,0.741) (7.36,0.726) (7.362,0.711) (7.365,0.695) (7.367,0.68) (7.369,0.665) (7.371,0.65) (7.374,0.634) (7.376,0.619) (7.378,0.603) (7.38,0.588) (7.383,0.572) (7.385,0.556) (7.387,0.541) (7.389,0.525) (7.392,0.509) (7.394,0.493) (7.396,0.478) (7.398,0.462) (7.401,0.446) (7.403,0.43) (7.405,0.414) (7.407,0.398) (7.41,0.381) (7.412,0.365) (7.414,0.349) (7.416,0.333) (7.419,0.316) (7.421,0.3) (7.423,0.283) (7.425,0.267) (7.428,0.25) (7.43,0.234) (7.432,0.217) (7.434,0.2) (7.437,0.184) (7.439,0.167) (7.441,0.15) (7.443,0.133) (7.446,0.116) (7.448,0.099) (7.45,0.082) (7.453,0.065) (7.455,0.048) (7.457,0.031) (7.459,0.014) (7.462,-0.004) (7.464,-0.021) (7.466,-0.038) (7.468,-0.056) (7.471,-0.073) (7.473,-0.091) (7.475,-0.108) (7.477,-0.126) (7.48,-0.143) (7.482,-0.161) (7.484,-0.179) (7.486,-0.197) (7.489,-0.214) (7.491,-0.232) (7.493,-0.25) (7.495,-0.268) (7.498,-0.286) (7.5,-0.304) (7.502,-0.322) (7.504,-0.341) (7.507,-0.359) (7.509,-0.377) (7.511,-0.395) (7.513,-0.414) (7.516,-0.432) (7.518,-0.451) (7.52,-0.469) (7.522,-0.488) (7.525,-0.506) (7.527,-0.525) (7.529,-0.544) (7.531,-0.562) (7.534,-0.581) (7.536,-0.6) (7.538,-0.619) (7.54,-0.638) (7.543,-0.657) (7.545,-0.676) (7.547,-0.695) (7.549,-0.714) (7.552,-0.733) (7.554,-0.753) (7.556,-0.772) (7.558,-0.791) (7.561,-0.811) (7.563,-0.83) (7.565,-0.849) (7.567,-0.869) (7.57,-0.889) (7.572,-0.908) (7.574,-0.928) (7.576,-0.948) (7.579,-0.967) (7.581,-0.987) (7.583,-1.007) (7.586,-1.027) (7.588,-1.047) (7.59,-1.067) (7.592,-1.087) (7.595,-1.107) (7.597,-1.127) (7.599,-1.147) (7.601,-1.168) (7.604,-1.188) (7.606,-1.208) (7.608,-1.229) (7.61,-1.249) (7.613,-1.269) (7.615,-1.29) (7.617,-1.311) (7.619,-1.331) (7.622,-1.352) (7.624,-1.373) (7.626,-1.393) (7.628,-1.414) (7.631,-1.435) (7.633,-1.456) (7.635,-1.477) (7.637,-1.498) (7.64,-1.519) (7.642,-1.54) (7.644,-1.561) (7.646,-1.582) (7.649,-1.604) (7.651,-1.625) (7.653,-1.646) (7.655,-1.668) (7.658,-1.689) (7.66,-1.711) (7.662,-1.732) (7.664,-1.754) (7.667,-1.775) (7.669,-1.797) (7.671,-1.819) (7.673,-1.841) (7.676,-1.862) (7.678,-1.884) (7.68,-1.906) (7.682,-1.928) (7.685,-1.95) (7.687,-1.972) (7.689,-1.994) (7.691,-2.016) (7.694,-2.039) (7.696,-2.061) (7.698,-2.083) (7.7,-2.106) (7.703,-2.128) (7.705,-2.15) (7.707,-2.173) (7.709,-2.195) (7.712,-2.218) (7.714,-2.241) (7.716,-2.263) (7.719,-2.286) (7.721,-2.309) (7.723,-2.332) (7.725,-2.355) (7.728,-2.378) (7.73,-2.401) (7.732,-2.424) (7.734,-2.447) (7.737,-2.47) (7.739,-2.493) (7.741,-2.516) (7.743,-2.539) (7.746,-2.563) (7.748,-2.586) (7.75,-2.609) (7.752,-2.633) (7.755,-2.656) (7.757,-2.68) (7.759,-2.704) (7.761,-2.727) (7.764,-2.751) (7.766,-2.775) (7.768,-2.798) (7.77,-2.822) (7.773,-2.846) (7.775,-2.87) (7.777,-2.894) (7.779,-2.918) (7.782,-2.942) (7.784,-2.966) (7.786,-2.991) (7.788,-3.015) (7.791,-3.039) (7.793,-3.063) (7.795,-3.088) (7.797,-3.112) (7.8,-3.137) (7.802,-3.161) (7.804,-3.186) (7.806,-3.21) (7.809,-3.235) (7.811,-3.26) (7.813,-3.284) (7.815,-3.309) (7.818,-3.334) (7.82,-3.359) (7.822,-3.384) (7.824,-3.409) (7.827,-3.434) (7.829,-3.459) (7.831,-3.484) (7.833,-3.509) (7.836,-3.535) (7.838,-3.56) (7.84,-3.585) (7.842,-3.611) (7.845,-3.636) (7.847,-3.661) (7.849,-3.687) (7.852,-3.713) (7.854,-3.738) (7.856,-3.764) (7.858,-3.79) (7.861,-3.815) (7.863,-3.841) (7.865,-3.867) (7.867,-3.893) (7.87,-3.919) (7.872,-3.945) (7.874,-3.971) (7.876,-3.997) (7.879,-4.023) (7.881,-4.049) (7.883,-4.076) (7.885,-4.102) (7.888,-4.128) (7.89,-4.155) (7.892,-4.181) (7.894,-4.207) (7.897,-4.234) (7.899,-4.261) (7.901,-4.287) (7.903,-4.314) (7.906,-4.341) (7.908,-4.367) (7.91,-4.394) (7.912,-4.421) (7.915,-4.448) (7.917,-4.475) (7.919,-4.502) (7.921,-4.529) (7.924,-4.556) (7.926,-4.583) (7.928,-4.61) (7.93,-4.638) (7.933,-4.665) (7.935,-4.692) (7.937,-4.72) (7.939,-4.747) (7.942,-4.775) (7.944,-4.802) (7.946,-4.83) (7.948,-4.857) (7.951,-4.885) (7.953,-4.913) (7.955,-4.941) (7.957,-4.968) (7.96,-4.996) (7.962,-5.024) (7.964,-5.052) (7.966,-5.08) (7.969,-5.108) (7.971,-5.136) (7.973,-5.165) (7.975,-5.193) (7.978,-5.221) (7.98,-5.249) (7.982,-5.278) (7.985,-5.306) (7.987,-5.335) (7.989,-5.363) (7.991,-5.392) (7.994,-5.42) (7.996,-5.449) (7.998,-5.477) (8,-5.506) (8.003,-5.535) (8.005,-5.564) (8.007,-5.593) (8.009,-5.622) (8.012,-5.651) (8.014,-5.68) (8.016,-5.709) (8.018,-5.738) (8.021,-5.767) (8.023,-5.796) (8.025,-5.826) (8.027,-5.855) (8.03,-5.884) (8.032,-5.914) (8.034,-5.943) (8.036,-5.973) (8.039,-6.002) (8.041,-6.032) (8.043,-6.061) (8.045,-6.091) (8.048,-6.121) (8.05,-6.151) (8.052,-6.18) (8.054,-6.21) (8.057,-6.24) (8.059,-6.27) (8.061,-6.3) (8.063,-6.33) (8.066,-6.361) (8.068,-6.391) (8.07,-6.421) (8.072,-6.451) (8.075,-6.482) (8.077,-6.512) (8.079,-6.542) (8.081,-6.573) (8.084,-6.603) (8.086,-6.634) (8.088,-6.664) (8.09,-6.695) (8.093,-6.726) (8.095,-6.757) (8.097,-6.787) (8.099,-6.818) (8.102,-6.849) (8.104,-6.88) (8.106,-6.911) (8.108,-6.942) (8.111,-6.973) (8.113,-7.004) (8.115,-7.036) (8.118,-7.067) (8.12,-7.098) (8.122,-7.129) (8.124,-7.161) (8.127,-7.192) (8.129,-7.224) (8.131,-7.255) (8.133,-7.287) (8.136,-7.318) (8.138,-7.35) (8.14,-7.382) (8.142,-7.414) (8.145,-7.445) (8.147,-7.477) (8.149,-7.509) (8.151,-7.541) (8.154,-7.573) (8.156,-7.605) (8.158,-7.637) (8.16,-7.669) (8.163,-7.702) (8.165,-7.734) (8.167,-7.766) (8.169,-7.798) (8.172,-7.831) (8.174,-7.863) (8.176,-7.896) (8.178,-7.928) (8.181,-7.961) (8.183,-7.993) (8.185,-8.026) (8.187,-8.059) (8.19,-8.092) (8.192,-8.124) (8.194,-8.157) (8.196,-8.19) (8.199,-8.223) (8.201,-8.256) (8.203,-8.289) (8.205,-8.322) (8.208,-8.356) (8.21,-8.389) (8.212,-8.422) (8.214,-8.455) (8.217,-8.489) (8.219,-8.522) (8.221,-8.555) (8.223,-8.589) (8.226,-8.623) (8.228,-8.656) (8.23,-8.69) (8.232,-8.723) (8.235,-8.757) (8.237,-8.791) (8.239,-8.825) (8.241,-8.859) (8.244,-8.893) (8.246,-8.927) (8.248,-8.961) (8.251,-8.995) (8.253,-9.029) (8.255,-9.063) (8.257,-9.097) (8.26,-9.131) (8.262,-9.166) (8.264,-9.2) (8.266,-9.234) (8.269,-9.269) (8.271,-9.303) (8.273,-9.338) (8.275,-9.373) (8.278,-9.407) (8.28,-9.442) (8.282,-9.477) (8.284,-9.511) (8.287,-9.546) (8.289,-9.581) (8.291,-9.616) (8.293,-9.651) (8.296,-9.686) (8.298,-9.721) (8.3,-9.756) (8.302,-9.791) (8.305,-9.827) (8.307,-9.862) (8.309,-9.897) (8.311,-9.933) (8.314,-9.968) (8.316,-10.003) (8.318,-10.039) (8.32,-10.075) (8.323,-10.11) (8.325,-10.146) (8.327,-10.181) (8.329,-10.217) (8.332,-10.253) (8.334,-10.289) (8.336,-10.325) (8.338,-10.361) (8.341,-10.397) (8.343,-10.433) (8.345,-10.469) (8.347,-10.505) (8.35,-10.541) (8.352,-10.577) (8.354,-10.614) (8.356,-10.65) (8.359,-10.686) (8.361,-10.723) (8.363,-10.759) (8.365,-10.796) (8.368,-10.832) (8.37,-10.869) (8.372,-10.906) (8.374,-10.942) (8.377,-10.979) (8.379,-11.016) (8.381,-11.053) (8.384,-11.09) (8.386,-11.127) (8.388,-11.164) (8.39,-11.201) (8.393,-11.238) (8.395,-11.275) (8.397,-11.312) (8.399,-11.349) (8.402,-11.387) (8.404,-11.424) (8.406,-11.461) (8.408,-11.499) (8.411,-11.536) (8.413,-11.574) (8.415,-11.611) (8.417,-11.649) (8.42,-11.687) (8.422,-11.724) (8.424,-11.762) (8.426,-11.8) (8.429,-11.838) (8.431,-11.876) (8.433,-11.914) (8.435,-11.952) (8.438,-11.99) (8.44,-12.028) (8.442,-12.066) (8.444,-12.104) (8.447,-12.142) (8.449,-12.181) (8.451,-12.219) (8.453,-12.257)};
\draw[BeachRed,line width=1.1pt] plot coordinates {(6.65,3.108) (6.666,3.107) (6.682,3.103) (6.697,3.097) (6.711,3.09) (6.724,3.082) (6.738,3.071) (6.749,3.061) (6.76,3.05) (6.763,3.048)};
\draw[BeachRed,line width=1.1pt] plot coordinates {(6.537,3.048) (6.549,3.059) (6.56,3.07) (6.573,3.08) (6.587,3.089) (6.6,3.096) (6.616,3.103) (6.632,3.106) (6.648,3.108)};
\end{scope}
\draw[black!10,line width=0.4pt] (9.883,0)--(9.883,4.4);
\draw[black!10,line width=0.4pt] (10.567,0)--(10.567,4.4);
\draw[black!10,line width=0.4pt] (11.25,0)--(11.25,4.4);
\draw[black!10,line width=0.4pt] (11.933,0)--(11.933,4.4);
\draw[black!10,line width=0.4pt] (12.617,0)--(12.617,4.4);
\draw[black!10,line width=0.4pt] (9.2,0)--(13.3,0);
\draw[black!10,line width=0.4pt] (9.2,0.983)--(13.3,0.983);
\draw[black!10,line width=0.4pt] (9.2,1.966)--(13.3,1.966);
\draw[black!10,line width=0.4pt] (9.2,2.948)--(13.3,2.948);
\draw[black!10,line width=0.4pt] (9.2,3.931)--(13.3,3.931);
\draw[line width=0.6pt] (9.2,0) rectangle (13.3,4.4);
\draw[line width=0.6pt] (9.883,0)--++(0,0.1);
\node[below,font=\scriptsize\sffamily] at (9.883,0) {$-4$};
\draw[line width=0.6pt] (10.567,0)--++(0,0.1);
\node[below,font=\scriptsize\sffamily] at (10.567,0) {$-2$};
\draw[line width=0.6pt] (11.25,0)--++(0,0.1);
\node[below,font=\scriptsize\sffamily] at (11.25,0) {$0$};
\draw[line width=0.6pt] (11.933,0)--++(0,0.1);
\node[below,font=\scriptsize\sffamily] at (11.933,0) {$2$};
\draw[line width=0.6pt] (12.617,0)--++(0,0.1);
\node[below,font=\scriptsize\sffamily] at (12.617,0) {$4$};
\draw[line width=0.6pt] (9.2,0)--++(0.1,0);
\draw[line width=0.6pt] (9.2,0.983)--++(0.1,0);
\draw[line width=0.6pt] (9.2,1.966)--++(0.1,0);
\draw[line width=0.6pt] (9.2,2.948)--++(0.1,0);
\draw[line width=0.6pt] (9.2,3.931)--++(0.1,0);
\node[font=\small\sffamily] at (11.25,-0.75) {$v_z$ [$10^6$ m/s]};
\node[above,font=\small\sffamily] at (11.25,4.4) {上側 $z=40$ m};
\begin{scope}
\clip (9.2,0) rectangle (13.3,4.4);
\draw[black!45,dashed,line width=0.5pt] (11.25,0)--(11.25,4.4);
\draw[BeachBlue,line width=1.1pt] plot coordinates {(7.039,-12.605) (7.044,-12.567) (7.049,-12.529) (7.054,-12.49) (7.06,-12.452) (7.065,-12.414) (7.07,-12.376) (7.075,-12.338) (7.081,-12.3) (7.086,-12.262) (7.091,-12.224) (7.096,-12.186) (7.102,-12.148) (7.107,-12.11) (7.112,-12.072) (7.118,-12.035) (7.123,-11.997) (7.128,-11.959) (7.133,-11.922) (7.139,-11.884) (7.144,-11.847) (7.149,-11.809) (7.154,-11.772) (7.16,-11.735) (7.165,-11.697) (7.17,-11.66) (7.175,-11.623) (7.181,-11.586) (7.186,-11.549) (7.191,-11.512) (7.196,-11.475) (7.202,-11.438) (7.207,-11.401) (7.212,-11.364) (7.218,-11.327) (7.223,-11.291) (7.228,-11.254) (7.233,-11.217) (7.239,-11.181) (7.244,-11.144) (7.249,-11.108) (7.254,-11.071) (7.26,-11.035) (7.265,-10.998) (7.27,-10.962) (7.275,-10.926) (7.281,-10.889) (7.286,-10.853) (7.291,-10.817) (7.296,-10.781) (7.302,-10.745) (7.307,-10.709) (7.312,-10.673) (7.318,-10.637) (7.323,-10.601) (7.328,-10.566) (7.333,-10.53) (7.339,-10.494) (7.344,-10.459) (7.349,-10.423) (7.354,-10.387) (7.36,-10.352) (7.365,-10.316) (7.37,-10.281) (7.375,-10.246) (7.381,-10.21) (7.386,-10.175) (7.391,-10.14) (7.397,-10.105) (7.402,-10.07) (7.407,-10.035) (7.412,-10) (7.418,-9.965) (7.423,-9.93) (7.428,-9.895) (7.433,-9.86) (7.439,-9.825) (7.444,-9.79) (7.449,-9.756) (7.454,-9.721) (7.46,-9.686) (7.465,-9.652) (7.47,-9.617) (7.475,-9.583) (7.481,-9.549) (7.486,-9.514) (7.491,-9.48) (7.497,-9.446) (7.502,-9.411) (7.507,-9.377) (7.512,-9.343) (7.518,-9.309) (7.523,-9.275) (7.528,-9.241) (7.533,-9.207) (7.539,-9.173) (7.544,-9.14) (7.549,-9.106) (7.554,-9.072) (7.56,-9.038) (7.565,-9.005) (7.57,-8.971) (7.575,-8.938) (7.581,-8.904) (7.586,-8.871) (7.591,-8.837) (7.597,-8.804) (7.602,-8.771) (7.607,-8.737) (7.612,-8.704) (7.618,-8.671) (7.623,-8.638) (7.628,-8.605) (7.633,-8.572) (7.639,-8.539) (7.644,-8.506) (7.649,-8.473) (7.654,-8.44) (7.66,-8.408) (7.665,-8.375) (7.67,-8.342) (7.676,-8.31) (7.681,-8.277) (7.686,-8.245) (7.691,-8.212) (7.697,-8.18) (7.702,-8.147) (7.707,-8.115) (7.712,-8.083) (7.718,-8.05) (7.723,-8.018) (7.728,-7.986) (7.733,-7.954) (7.739,-7.922) (7.744,-7.89) (7.749,-7.858) (7.754,-7.826) (7.76,-7.794) (7.765,-7.762) (7.77,-7.731) (7.776,-7.699) (7.781,-7.667) (7.786,-7.636) (7.791,-7.604) (7.797,-7.573) (7.802,-7.541) (7.807,-7.51) (7.812,-7.478) (7.818,-7.447) (7.823,-7.416) (7.828,-7.384) (7.833,-7.353) (7.839,-7.322) (7.844,-7.291) (7.849,-7.26) (7.855,-7.229) (7.86,-7.198) (7.865,-7.167) (7.87,-7.136) (7.876,-7.106) (7.881,-7.075) (7.886,-7.044) (7.891,-7.013) (7.897,-6.983) (7.902,-6.952) (7.907,-6.922) (7.912,-6.891) (7.918,-6.861) (7.923,-6.831) (7.928,-6.8) (7.933,-6.77) (7.939,-6.74) (7.944,-6.71) (7.949,-6.679) (7.955,-6.649) (7.96,-6.619) (7.965,-6.589) (7.97,-6.559) (7.976,-6.53) (7.981,-6.5) (7.986,-6.47) (7.991,-6.44) (7.997,-6.41) (8.002,-6.381) (8.007,-6.351) (8.012,-6.322) (8.018,-6.292) (8.023,-6.263) (8.028,-6.233) (8.033,-6.204) (8.039,-6.175) (8.044,-6.145) (8.049,-6.116) (8.055,-6.087) (8.06,-6.058) (8.065,-6.029) (8.07,-6) (8.076,-5.971) (8.081,-5.942) (8.086,-5.913) (8.091,-5.884) (8.097,-5.855) (8.102,-5.827) (8.107,-5.798) (8.112,-5.769) (8.118,-5.741) (8.123,-5.712) (8.128,-5.684) (8.134,-5.655) (8.139,-5.627) (8.144,-5.599) (8.149,-5.57) (8.155,-5.542) (8.16,-5.514) (8.165,-5.486) (8.17,-5.458) (8.176,-5.429) (8.181,-5.401) (8.186,-5.373) (8.191,-5.346) (8.197,-5.318) (8.202,-5.29) (8.207,-5.262) (8.212,-5.234) (8.218,-5.207) (8.223,-5.179) (8.228,-5.152) (8.234,-5.124) (8.239,-5.096) (8.244,-5.069) (8.249,-5.042) (8.255,-5.014) (8.26,-4.987) (8.265,-4.96) (8.27,-4.933) (8.276,-4.905) (8.281,-4.878) (8.286,-4.851) (8.291,-4.824) (8.297,-4.797) (8.302,-4.77) (8.307,-4.744) (8.313,-4.717) (8.318,-4.69) (8.323,-4.663) (8.328,-4.637) (8.334,-4.61) (8.339,-4.583) (8.344,-4.557) (8.349,-4.53) (8.355,-4.504) (8.36,-4.478) (8.365,-4.451) (8.37,-4.425) (8.376,-4.399) (8.381,-4.373) (8.386,-4.346) (8.391,-4.32) (8.397,-4.294) (8.402,-4.268) (8.407,-4.242) (8.413,-4.216) (8.418,-4.191) (8.423,-4.165) (8.428,-4.139) (8.434,-4.113) (8.439,-4.088) (8.444,-4.062) (8.449,-4.036) (8.455,-4.011) (8.46,-3.985) (8.465,-3.96) (8.47,-3.935) (8.476,-3.909) (8.481,-3.884) (8.486,-3.859) (8.491,-3.834) (8.497,-3.809) (8.502,-3.783) (8.507,-3.758) (8.513,-3.733) (8.518,-3.708) (8.523,-3.684) (8.528,-3.659) (8.534,-3.634) (8.539,-3.609) (8.544,-3.585) (8.549,-3.56) (8.555,-3.535) (8.56,-3.511) (8.565,-3.486) (8.57,-3.462) (8.576,-3.437) (8.581,-3.413) (8.586,-3.389) (8.592,-3.364) (8.597,-3.34) (8.602,-3.316) (8.607,-3.292) (8.613,-3.268) (8.618,-3.244) (8.623,-3.22) (8.628,-3.196) (8.634,-3.172) (8.639,-3.148) (8.644,-3.124) (8.649,-3.101) (8.655,-3.077) (8.66,-3.053) (8.665,-3.03) (8.67,-3.006) (8.676,-2.983) (8.681,-2.959) (8.686,-2.936) (8.692,-2.912) (8.697,-2.889) (8.702,-2.866) (8.707,-2.843) (8.713,-2.819) (8.718,-2.796) (8.723,-2.773) (8.728,-2.75) (8.734,-2.727) (8.739,-2.704) (8.744,-2.681) (8.749,-2.659) (8.755,-2.636) (8.76,-2.613) (8.765,-2.59) (8.771,-2.568) (8.776,-2.545) (8.781,-2.523) (8.786,-2.5) (8.792,-2.478) (8.797,-2.455) (8.802,-2.433) (8.807,-2.411) (8.813,-2.388) (8.818,-2.366) (8.823,-2.344) (8.828,-2.322) (8.834,-2.3) (8.839,-2.278) (8.844,-2.256) (8.849,-2.234) (8.855,-2.212) (8.86,-2.19) (8.865,-2.169) (8.871,-2.147) (8.876,-2.125) (8.881,-2.104) (8.886,-2.082) (8.892,-2.06) (8.897,-2.039) (8.902,-2.018) (8.907,-1.996) (8.913,-1.975) (8.918,-1.954) (8.923,-1.932) (8.928,-1.911) (8.934,-1.89) (8.939,-1.869) (8.944,-1.848) (8.949,-1.827) (8.955,-1.806) (8.96,-1.785) (8.965,-1.764) (8.971,-1.743) (8.976,-1.722) (8.981,-1.702) (8.986,-1.681) (8.992,-1.66) (8.997,-1.64) (9.002,-1.619) (9.007,-1.599) (9.013,-1.578) (9.018,-1.558) (9.023,-1.538) (9.028,-1.517) (9.034,-1.497) (9.039,-1.477) (9.044,-1.457) (9.05,-1.437) (9.055,-1.417) (9.06,-1.397) (9.065,-1.377) (9.071,-1.357) (9.076,-1.337) (9.081,-1.317) (9.086,-1.298) (9.092,-1.278) (9.097,-1.258) (9.102,-1.239) (9.107,-1.219) (9.113,-1.199) (9.118,-1.18) (9.123,-1.161) (9.128,-1.141) (9.134,-1.122) (9.139,-1.103) (9.144,-1.083) (9.15,-1.064) (9.155,-1.045) (9.16,-1.026) (9.165,-1.007) (9.171,-0.988) (9.176,-0.969) (9.181,-0.95) (9.186,-0.931) (9.192,-0.912) (9.197,-0.894) (9.202,-0.875) (9.207,-0.856) (9.213,-0.838) (9.218,-0.819) (9.223,-0.801) (9.228,-0.782) (9.234,-0.764) (9.239,-0.745) (9.244,-0.727) (9.25,-0.709) (9.255,-0.691) (9.26,-0.673) (9.265,-0.654) (9.271,-0.636) (9.276,-0.618) (9.281,-0.6) (9.286,-0.582) (9.292,-0.564) (9.297,-0.547) (9.302,-0.529) (9.307,-0.511) (9.313,-0.493) (9.318,-0.476) (9.323,-0.458) (9.329,-0.441) (9.334,-0.423) (9.339,-0.406) (9.344,-0.388) (9.35,-0.371) (9.355,-0.354) (9.36,-0.336) (9.365,-0.319) (9.371,-0.302) (9.376,-0.285) (9.381,-0.268) (9.386,-0.251) (9.392,-0.234) (9.397,-0.217) (9.402,-0.2) (9.407,-0.183) (9.413,-0.166) (9.418,-0.15) (9.423,-0.133) (9.429,-0.116) (9.434,-0.1) (9.439,-0.083) (9.444,-0.067) (9.45,-0.05) (9.455,-0.034) (9.46,-0.018) (9.465,-0.001) (9.471,0.015) (9.476,0.031) (9.481,0.047) (9.486,0.063) (9.492,0.079) (9.497,0.095) (9.502,0.111) (9.508,0.127) (9.513,0.143) (9.518,0.159) (9.523,0.175) (9.529,0.191) (9.534,0.206) (9.539,0.222) (9.544,0.237) (9.55,0.253) (9.555,0.268) (9.56,0.284) (9.565,0.299) (9.571,0.315) (9.576,0.33) (9.581,0.345) (9.586,0.36) (9.592,0.375) (9.597,0.391) (9.602,0.406) (9.608,0.421) (9.613,0.436) (9.618,0.45) (9.623,0.465) (9.629,0.48) (9.634,0.495) (9.639,0.51) (9.644,0.524) (9.65,0.539) (9.655,0.553) (9.66,0.568) (9.665,0.582) (9.671,0.597) (9.676,0.611) (9.681,0.626) (9.686,0.64) (9.692,0.654) (9.697,0.668) (9.702,0.682) (9.708,0.697) (9.718,0.725) (9.729,0.752) (9.739,0.78) (9.75,0.808) (9.76,0.835) (9.771,0.862) (9.781,0.889) (9.792,0.915) (9.802,0.942) (9.813,0.968) (9.823,0.994) (9.834,1.02) (9.844,1.046) (9.855,1.072) (9.865,1.097) (9.876,1.122) (9.887,1.147) (9.897,1.172) (9.908,1.196) (9.918,1.221) (9.929,1.245) (9.939,1.269) (9.95,1.293) (9.96,1.316) (9.971,1.34) (9.981,1.363) (9.992,1.386) (10.002,1.409) (10.013,1.432) (10.023,1.454) (10.034,1.476) (10.044,1.498) (10.055,1.52) (10.066,1.542) (10.076,1.564) (10.087,1.585) (10.097,1.606) (10.108,1.627) (10.118,1.648) (10.129,1.668) (10.139,1.689) (10.15,1.709) (10.16,1.729) (10.171,1.749) (10.181,1.768) (10.192,1.788) (10.202,1.807) (10.213,1.826) (10.223,1.845) (10.234,1.863) (10.245,1.882) (10.255,1.9) (10.266,1.918) (10.276,1.936) (10.287,1.953) (10.297,1.971) (10.308,1.988) (10.318,2.005) (10.329,2.022) (10.339,2.039) (10.35,2.055) (10.36,2.072) (10.371,2.088) (10.381,2.104) (10.392,2.12) (10.402,2.135) (10.413,2.15) (10.424,2.166) (10.434,2.181) (10.445,2.195) (10.455,2.21) (10.466,2.224) (10.476,2.239) (10.487,2.253) (10.497,2.266) (10.508,2.28) (10.518,2.294) (10.529,2.307) (10.539,2.32) (10.55,2.333) (10.56,2.345) (10.571,2.358) (10.581,2.37) (10.592,2.382) (10.602,2.394) (10.613,2.406) (10.624,2.417) (10.634,2.429) (10.645,2.44) (10.655,2.451) (10.666,2.462) (10.681,2.477) (10.697,2.493) (10.713,2.508) (10.729,2.522) (10.745,2.536) (10.76,2.55) (10.776,2.563) (10.792,2.576) (10.808,2.588) (10.824,2.6) (10.839,2.611) (10.855,2.622) (10.871,2.633) (10.887,2.643) (10.903,2.653) (10.918,2.662) (10.934,2.671) (10.95,2.679) (10.966,2.687) (10.982,2.695) (10.997,2.702) (11.013,2.709) (11.029,2.715) (11.045,2.721) (11.06,2.726) (11.076,2.731) (11.092,2.736) (11.108,2.74) (11.124,2.743) (11.139,2.747) (11.145,2.748)};
\draw[BeachCyan,line width=1.1pt] plot coordinates {(11.15,2.749) (11.166,2.751) (11.182,2.753) (11.197,2.755) (11.213,2.756) (11.229,2.757) (11.245,2.757) (11.25,2.757)};
\draw[BeachCyan,line width=1.1pt] plot coordinates {(11.255,2.757) (11.271,2.757) (11.287,2.756) (11.303,2.755) (11.318,2.753) (11.334,2.751) (11.35,2.749)};
\draw[BeachOrange,line width=1.1pt] plot coordinates {(11.354,3.046) (11.365,3.035) (11.376,3.022) (11.388,3.008) (11.397,2.996) (11.406,2.983) (11.415,2.969) (11.424,2.955) (11.433,2.94) (11.442,2.924) (11.451,2.907) (11.46,2.889) (11.466,2.876) (11.473,2.862) (11.48,2.847) (11.487,2.832) (11.493,2.817) (11.5,2.801) (11.507,2.785) (11.514,2.768) (11.521,2.751) (11.527,2.734) (11.534,2.716) (11.541,2.698) (11.548,2.679) (11.554,2.66) (11.561,2.64) (11.568,2.62) (11.575,2.599) (11.581,2.578) (11.588,2.557) (11.593,2.542) (11.597,2.528) (11.602,2.513) (11.606,2.498) (11.611,2.482) (11.615,2.467) (11.62,2.451) (11.624,2.436) (11.629,2.419) (11.633,2.403) (11.638,2.387) (11.642,2.37) (11.647,2.353) (11.651,2.336) (11.656,2.319) (11.66,2.302) (11.665,2.284) (11.669,2.267) (11.674,2.249) (11.678,2.23) (11.683,2.212) (11.687,2.194) (11.692,2.175) (11.696,2.156) (11.701,2.137) (11.705,2.117) (11.71,2.098) (11.714,2.078) (11.719,2.058) (11.723,2.038) (11.728,2.018) (11.732,1.998) (11.737,1.977) (11.741,1.956) (11.746,1.935) (11.75,1.914) (11.755,1.892) (11.759,1.871) (11.764,1.849) (11.768,1.827) (11.773,1.805) (11.777,1.783) (11.782,1.76) (11.787,1.737) (11.791,1.714) (11.796,1.691) (11.8,1.668) (11.805,1.644) (11.809,1.62) (11.814,1.597) (11.818,1.572) (11.823,1.548) (11.827,1.524) (11.832,1.499) (11.836,1.474) (11.841,1.449) (11.845,1.424) (11.85,1.398) (11.854,1.373) (11.859,1.347) (11.863,1.321) (11.868,1.295) (11.872,1.268) (11.877,1.242) (11.881,1.215) (11.886,1.188) (11.89,1.161) (11.895,1.133) (11.899,1.106) (11.904,1.078) (11.908,1.05) (11.913,1.022) (11.917,0.994) (11.922,0.965) (11.926,0.936) (11.931,0.907) (11.935,0.878) (11.94,0.849) (11.944,0.82) (11.949,0.79) (11.951,0.775) (11.953,0.76) (11.956,0.745) (11.958,0.73) (11.96,0.715) (11.962,0.7) (11.965,0.685) (11.967,0.669) (11.969,0.654) (11.971,0.639) (11.974,0.623) (11.976,0.608) (11.978,0.592) (11.98,0.577) (11.983,0.561) (11.985,0.546) (11.987,0.53) (11.989,0.514) (11.992,0.498) (11.994,0.483) (11.996,0.467) (11.998,0.451) (12.001,0.435) (12.003,0.419) (12.005,0.403) (12.007,0.387) (12.01,0.371) (12.012,0.354) (12.014,0.338) (12.016,0.322) (12.019,0.305) (12.021,0.289) (12.023,0.273) (12.025,0.256) (12.028,0.24) (12.03,0.223) (12.032,0.206) (12.034,0.19) (12.037,0.173) (12.039,0.156) (12.041,0.139) (12.043,0.122) (12.046,0.105) (12.048,0.088) (12.05,0.071) (12.053,0.054) (12.055,0.037) (12.057,0.02) (12.059,0.003) (12.062,-0.014) (12.064,-0.032) (12.066,-0.049) (12.068,-0.066) (12.071,-0.084) (12.073,-0.101) (12.075,-0.119) (12.077,-0.137) (12.08,-0.154) (12.082,-0.172) (12.084,-0.19) (12.086,-0.207) (12.089,-0.225) (12.091,-0.243) (12.093,-0.261) (12.095,-0.279) (12.098,-0.297) (12.1,-0.315) (12.102,-0.333) (12.104,-0.351) (12.107,-0.37) (12.109,-0.388) (12.111,-0.406) (12.113,-0.425) (12.116,-0.443) (12.118,-0.462) (12.12,-0.48) (12.122,-0.499) (12.125,-0.517) (12.127,-0.536) (12.129,-0.555) (12.131,-0.573) (12.134,-0.592) (12.136,-0.611) (12.138,-0.63) (12.14,-0.649) (12.143,-0.668) (12.145,-0.687) (12.147,-0.706) (12.149,-0.725) (12.152,-0.744) (12.154,-0.763) (12.156,-0.783) (12.158,-0.802) (12.161,-0.821) (12.163,-0.841) (12.165,-0.86) (12.167,-0.88) (12.17,-0.899) (12.172,-0.919) (12.174,-0.939) (12.176,-0.958) (12.179,-0.978) (12.181,-0.998) (12.183,-1.018) (12.186,-1.038) (12.188,-1.058) (12.19,-1.078) (12.192,-1.098) (12.195,-1.118) (12.197,-1.138) (12.199,-1.158) (12.201,-1.178) (12.204,-1.199) (12.206,-1.219) (12.208,-1.239) (12.21,-1.26) (12.213,-1.28) (12.215,-1.301) (12.217,-1.321) (12.219,-1.342) (12.222,-1.363) (12.224,-1.383) (12.226,-1.404) (12.228,-1.425) (12.231,-1.446) (12.233,-1.467) (12.235,-1.488) (12.237,-1.509) (12.24,-1.53) (12.242,-1.551) (12.244,-1.572) (12.246,-1.593) (12.249,-1.614) (12.251,-1.636) (12.253,-1.657) (12.255,-1.679) (12.258,-1.7) (12.26,-1.721) (12.262,-1.743) (12.264,-1.765) (12.267,-1.786) (12.269,-1.808) (12.271,-1.83) (12.273,-1.851) (12.276,-1.873) (12.278,-1.895) (12.28,-1.917) (12.282,-1.939) (12.285,-1.961) (12.287,-1.983) (12.289,-2.005) (12.291,-2.027) (12.294,-2.05) (12.296,-2.072) (12.298,-2.094) (12.3,-2.116) (12.303,-2.139) (12.305,-2.161) (12.307,-2.184) (12.309,-2.206) (12.312,-2.229) (12.314,-2.252) (12.316,-2.274) (12.319,-2.297) (12.321,-2.32) (12.323,-2.343) (12.325,-2.365) (12.328,-2.388) (12.33,-2.411) (12.332,-2.434) (12.334,-2.457) (12.337,-2.481) (12.339,-2.504) (12.341,-2.527) (12.343,-2.55) (12.346,-2.574) (12.348,-2.597) (12.35,-2.62) (12.352,-2.644) (12.355,-2.667) (12.357,-2.691) (12.359,-2.714) (12.361,-2.738) (12.364,-2.762) (12.366,-2.786) (12.368,-2.809) (12.37,-2.833) (12.373,-2.857) (12.375,-2.881) (12.377,-2.905) (12.379,-2.929) (12.382,-2.953) (12.384,-2.977) (12.386,-3.001) (12.388,-3.026) (12.391,-3.05) (12.393,-3.074) (12.395,-3.099) (12.397,-3.123) (12.4,-3.147) (12.402,-3.172) (12.404,-3.197) (12.406,-3.221) (12.409,-3.246) (12.411,-3.27) (12.413,-3.295) (12.415,-3.32) (12.418,-3.345) (12.42,-3.37) (12.422,-3.395) (12.424,-3.42) (12.427,-3.445) (12.429,-3.47) (12.431,-3.495) (12.433,-3.52) (12.436,-3.545) (12.438,-3.571) (12.44,-3.596) (12.442,-3.621) (12.445,-3.647) (12.447,-3.672) (12.449,-3.698) (12.452,-3.723) (12.454,-3.749) (12.456,-3.775) (12.458,-3.8) (12.461,-3.826) (12.463,-3.852) (12.465,-3.878) (12.467,-3.904) (12.47,-3.93) (12.472,-3.956) (12.474,-3.982) (12.476,-4.008) (12.479,-4.034) (12.481,-4.06) (12.483,-4.086) (12.485,-4.113) (12.488,-4.139) (12.49,-4.165) (12.492,-4.192) (12.494,-4.218) (12.497,-4.245) (12.499,-4.271) (12.501,-4.298) (12.503,-4.325) (12.506,-4.351) (12.508,-4.378) (12.51,-4.405) (12.512,-4.432) (12.515,-4.459) (12.517,-4.486) (12.519,-4.513) (12.521,-4.54) (12.524,-4.567) (12.526,-4.594) (12.528,-4.621) (12.53,-4.649) (12.533,-4.676) (12.535,-4.703) (12.537,-4.731) (12.539,-4.758) (12.542,-4.786) (12.544,-4.813) (12.546,-4.841) (12.548,-4.868) (12.551,-4.896) (12.553,-4.924) (12.555,-4.951) (12.557,-4.979) (12.56,-5.007) (12.562,-5.035) (12.564,-5.063) (12.566,-5.091) (12.569,-5.119) (12.571,-5.147) (12.573,-5.175) (12.575,-5.204) (12.578,-5.232) (12.58,-5.26) (12.582,-5.289) (12.585,-5.317) (12.587,-5.345) (12.589,-5.374) (12.591,-5.402) (12.594,-5.431) (12.596,-5.46) (12.598,-5.488) (12.6,-5.517) (12.603,-5.546) (12.605,-5.575) (12.607,-5.604) (12.609,-5.633) (12.612,-5.661) (12.614,-5.691) (12.616,-5.72) (12.618,-5.749) (12.621,-5.778) (12.623,-5.807) (12.625,-5.836) (12.627,-5.866) (12.63,-5.895) (12.632,-5.924) (12.634,-5.954) (12.636,-5.983) (12.639,-6.013) (12.641,-6.043) (12.643,-6.072) (12.645,-6.102) (12.648,-6.132) (12.65,-6.161) (12.652,-6.191) (12.654,-6.221) (12.657,-6.251) (12.659,-6.281) (12.661,-6.311) (12.663,-6.341) (12.666,-6.371) (12.668,-6.402) (12.67,-6.432) (12.672,-6.462) (12.675,-6.492) (12.677,-6.523) (12.679,-6.553) (12.681,-6.584) (12.684,-6.614) (12.686,-6.645) (12.688,-6.675) (12.69,-6.706) (12.693,-6.737) (12.695,-6.767) (12.697,-6.798) (12.699,-6.829) (12.702,-6.86) (12.704,-6.891) (12.706,-6.922) (12.708,-6.953) (12.711,-6.984) (12.713,-7.015) (12.715,-7.046) (12.718,-7.078) (12.72,-7.109) (12.722,-7.14) (12.724,-7.172) (12.727,-7.203) (12.729,-7.235) (12.731,-7.266) (12.733,-7.298) (12.736,-7.329) (12.738,-7.361) (12.74,-7.393) (12.742,-7.424) (12.745,-7.456) (12.747,-7.488) (12.749,-7.52) (12.751,-7.552) (12.754,-7.584) (12.756,-7.616) (12.758,-7.648) (12.76,-7.68) (12.763,-7.712) (12.765,-7.745) (12.767,-7.777) (12.769,-7.809) (12.772,-7.842) (12.774,-7.874) (12.776,-7.907) (12.778,-7.939) (12.781,-7.972) (12.783,-8.004) (12.785,-8.037) (12.787,-8.07) (12.79,-8.102) (12.792,-8.135) (12.794,-8.168) (12.796,-8.201) (12.799,-8.234) (12.801,-8.267) (12.803,-8.3) (12.805,-8.333) (12.808,-8.366) (12.81,-8.4) (12.812,-8.433) (12.814,-8.466) (12.817,-8.499) (12.819,-8.533) (12.821,-8.566) (12.823,-8.6) (12.826,-8.633) (12.828,-8.667) (12.83,-8.701) (12.832,-8.734) (12.835,-8.768) (12.837,-8.802) (12.839,-8.836) (12.841,-8.869) (12.844,-8.903) (12.846,-8.937) (12.848,-8.971) (12.851,-9.005) (12.853,-9.04) (12.855,-9.074) (12.857,-9.108) (12.86,-9.142) (12.862,-9.177) (12.864,-9.211) (12.866,-9.245) (12.869,-9.28) (12.871,-9.314) (12.873,-9.349) (12.875,-9.383) (12.878,-9.418) (12.88,-9.453) (12.882,-9.487) (12.884,-9.522) (12.887,-9.557) (12.889,-9.592) (12.891,-9.627) (12.893,-9.662) (12.896,-9.697) (12.898,-9.732) (12.9,-9.767) (12.902,-9.802) (12.905,-9.837) (12.907,-9.873) (12.909,-9.908) (12.911,-9.943) (12.914,-9.979) (12.916,-10.014) (12.918,-10.05) (12.92,-10.085) (12.923,-10.121) (12.925,-10.157) (12.927,-10.192) (12.929,-10.228) (12.932,-10.264) (12.934,-10.3) (12.936,-10.336) (12.938,-10.372) (12.941,-10.408) (12.943,-10.444) (12.945,-10.48) (12.947,-10.516) (12.95,-10.552) (12.952,-10.588) (12.954,-10.624) (12.956,-10.661) (12.959,-10.697) (12.961,-10.734) (12.963,-10.77) (12.965,-10.807) (12.968,-10.843) (12.97,-10.88) (12.972,-10.916) (12.974,-10.953) (12.977,-10.99) (12.979,-11.027) (12.981,-11.064) (12.984,-11.1) (12.986,-11.137) (12.988,-11.174) (12.99,-11.211) (12.993,-11.249) (12.995,-11.286) (12.997,-11.323) (12.999,-11.36) (13.002,-11.397) (13.004,-11.435) (13.006,-11.472) (13.008,-11.51) (13.011,-11.547) (13.013,-11.585) (13.015,-11.622) (13.017,-11.66) (13.02,-11.697) (13.022,-11.735) (13.024,-11.773) (13.026,-11.811) (13.029,-11.849) (13.031,-11.886) (13.033,-11.924) (13.035,-11.962) (13.038,-12) (13.04,-12.039) (13.042,-12.077) (13.044,-12.115) (13.047,-12.153) (13.049,-12.191) (13.051,-12.23) (13.053,-12.268)};
\end{scope}
\draw[BeachBlue,line width=1.1pt] (0.3,-1.25)--(0.9,-1.25);
\node[right,font=\scriptsize\sffamily] at (0.9,-1.25) {太陽風電子（通過、壁へ向かう）};
\draw[BeachCyan,line width=1.1pt] (0.3,-1.7)--(0.9,-1.7);
\node[right,font=\scriptsize\sffamily] at (0.9,-1.7) {太陽風電子（反射、往復）};
\draw[BeachOrange,line width=1.1pt] (7,-1.25)--(7.6,-1.25);
\node[right,font=\scriptsize\sffamily] at (7.6,-1.25) {光電子（脱出、上流へ向かう）};
\draw[BeachRed,line width=1.1pt] (7,-1.7)--(7.6,-1.7);
\node[right,font=\scriptsize\sffamily] at (7.6,-1.7) {光電子（捕捉、往復）};
\end{tikzpicture}
\caption[Type A の電子の速度分布]{Type A（$\alpha=60^\circ$）の電子の法線速度分布 $f(v_z)$。$v_z>0$ は上流向きで、縦軸は対数、$v_{te}=\sqrt{2eT_e/m_e}$ と上流イオン密度 $n_i$ で規格化した。分布は上流の Maxwell 分布と壁の半 Maxwell 分布をエネルギー保存で写したもので、到達できない速度で 0 になる。壁では太陽風電子は $v_z<0$ だけ、光電子は $v_z>0$ だけを持つ。下側では捕捉光電子が $|v_z|$ の小さい領域を両方向に占め、上側では反射電子が同じ役割を担う。}
\label{fig:sm_vdf}
\end{figure}


\subsection{太陽高度への応答}
図\ref{fig:sm_alpha_sweep}は、太陽高度を $1^\circ$ から $89^\circ$ まで変えたときの零電流根である。
高度が下がると光電子の放出 $N_{\mathrm{pe}}\propto\sin\alpha$ とイオンの流束 $\propto\sin\alpha$ がともに減り、壁電位が下がる。
この条件では、Type C は $\alpha=5$--$11^\circ$、Type A と B は $\alpha\ge12^\circ$ と $\alpha\ge11^\circ$ で見つかった。
Type A の極小は高度が下がると深くなり、$\alpha=12^\circ$ では $\phi_0=-1.73$ V、$\phi_m=-1.79$ V と壁電位に近づいて Type C へ移る。
同じ高度で A と B が両方見つかる範囲では、B の壁電位が A より高い。
$\alpha\le4^\circ$ で根が見つからないのは、イオンの法線速度 $u_i=v_{\mathrm{sw}}\sin\alpha$ が Bohm 速度の目安
$\sqrt{eT_e/m_i}\approx34$ km/s を下回り、上流の漸近条件を満たせないためと考えられる。
下段の数流束は、太陽風電子の吸収がイオンの吸収と光電子の脱出の和と釣り合うことを示す。

% Generated by tools/sheath_model_figures.py from alpha_sweep.csv; edit the generator, not this file.
\begin{figure}[htbp]
\centering
\begin{tikzpicture}[beach diagram]
\draw[black!10,line width=0.4pt] (0,4.6)--(0,8.6);
\draw[black!10,line width=0.4pt] (1.583,4.6)--(1.583,8.6);
\draw[black!10,line width=0.4pt] (3.167,4.6)--(3.167,8.6);
\draw[black!10,line width=0.4pt] (4.75,4.6)--(4.75,8.6);
\draw[black!10,line width=0.4pt] (6.333,4.6)--(6.333,8.6);
\draw[black!10,line width=0.4pt] (7.917,4.6)--(7.917,8.6);
\draw[black!10,line width=0.4pt] (9.5,4.6)--(9.5,8.6);
\draw[black!10,line width=0.4pt] (0,4.782)--(9.5,4.782);
\draw[black!10,line width=0.4pt] (0,5.691)--(9.5,5.691);
\draw[black!10,line width=0.4pt] (0,6.6)--(9.5,6.6);
\draw[black!10,line width=0.4pt] (0,7.509)--(9.5,7.509);
\draw[black!10,line width=0.4pt] (0,8.418)--(9.5,8.418);
\draw[line width=0.6pt] (0,4.6) rectangle (9.5,8.6);
\draw[line width=0.6pt] (0,4.6)--++(0,0.1);
\draw[line width=0.6pt] (1.583,4.6)--++(0,0.1);
\draw[line width=0.6pt] (3.167,4.6)--++(0,0.1);
\draw[line width=0.6pt] (4.75,4.6)--++(0,0.1);
\draw[line width=0.6pt] (6.333,4.6)--++(0,0.1);
\draw[line width=0.6pt] (7.917,4.6)--++(0,0.1);
\draw[line width=0.6pt] (9.5,4.6)--++(0,0.1);
\draw[line width=0.6pt] (0,4.782)--++(0.1,0);
\node[left,font=\scriptsize\sffamily] at (0,4.782) {$-15$};
\draw[line width=0.6pt] (0,5.691)--++(0.1,0);
\node[left,font=\scriptsize\sffamily] at (0,5.691) {$-10$};
\draw[line width=0.6pt] (0,6.6)--++(0.1,0);
\node[left,font=\scriptsize\sffamily] at (0,6.6) {$-5$};
\draw[line width=0.6pt] (0,7.509)--++(0.1,0);
\node[left,font=\scriptsize\sffamily] at (0,7.509) {$0$};
\draw[line width=0.6pt] (0,8.418)--++(0.1,0);
\node[left,font=\scriptsize\sffamily] at (0,8.418) {$5$};
\node[rotate=90,font=\small\sffamily] at (-1.05,6.6) {電位 [V]};
\begin{scope}
\clip (0,4.6) rectangle (9.5,8.6);
\draw[black!45,dashed,line width=0.5pt] (0,7.509)--(9.5,7.509);
\draw[BeachBlue,solid,line width=1.2pt] plot coordinates {(1.267,7.195) (1.372,7.245) (1.478,7.291) (1.583,7.334) (1.689,7.374) (1.794,7.411) (1.9,7.447) (2.006,7.481) (2.111,7.513) (2.217,7.544) (2.322,7.573) (2.428,7.602) (2.533,7.629) (2.639,7.655) (2.744,7.681) (2.85,7.705) (2.956,7.729) (3.061,7.752) (3.167,7.774) (3.272,7.796) (3.378,7.817) (3.483,7.837) (3.589,7.857) (3.694,7.876) (3.8,7.894) (3.906,7.912) (4.011,7.93) (4.117,7.947) (4.222,7.963) (4.328,7.979) (4.433,7.995) (4.539,8.01) (4.644,8.025) (4.75,8.039) (4.856,8.053) (4.961,8.066) (5.067,8.079) (5.172,8.092) (5.278,8.104) (5.383,8.116) (5.489,8.128) (5.594,8.139) (5.7,8.15) (5.806,8.16) (5.911,8.17) (6.017,8.18) (6.122,8.189) (6.228,8.199) (6.333,8.207) (6.439,8.216) (6.544,8.224) (6.65,8.232) (6.756,8.239) (6.861,8.246) (6.967,8.253) (7.072,8.26) (7.178,8.266) (7.283,8.272) (7.389,8.278) (7.494,8.283) (7.6,8.288) (7.706,8.293) (7.811,8.298) (7.917,8.302) (8.022,8.306) (8.128,8.31) (8.233,8.313) (8.339,8.316) (8.444,8.319) (8.55,8.322) (8.656,8.324) (8.761,8.326) (8.867,8.328) (8.972,8.329) (9.078,8.33) (9.183,8.331) (9.289,8.332) (9.394,8.332)};
\draw[BeachBlue,dotted,line width=1.2pt] plot coordinates {(1.267,7.184) (1.372,7.199) (1.478,7.212) (1.583,7.224) (1.689,7.235) (1.794,7.245) (1.9,7.255) (2.006,7.263) (2.111,7.272) (2.217,7.28) (2.322,7.288) (2.428,7.295) (2.533,7.302) (2.639,7.309) (2.744,7.316) (2.85,7.322) (2.956,7.329) (3.061,7.335) (3.167,7.341) (3.272,7.346) (3.378,7.352) (3.483,7.357) (3.589,7.362) (3.694,7.367) (3.8,7.372) (3.906,7.377) (4.011,7.381) (4.117,7.386) (4.222,7.39) (4.328,7.394) (4.433,7.398) (4.539,7.402) (4.644,7.405) (4.75,7.409) (4.856,7.412) (4.961,7.416) (5.067,7.419) (5.172,7.422) (5.278,7.425) (5.383,7.428) (5.489,7.43) (5.594,7.433) (5.7,7.436) (5.806,7.438) (5.911,7.44) (6.017,7.443) (6.122,7.445) (6.228,7.447) (6.333,7.449) (6.439,7.451) (6.544,7.452) (6.65,7.454) (6.756,7.456) (6.861,7.457) (6.967,7.459) (7.072,7.46) (7.178,7.461) (7.283,7.463) (7.389,7.464) (7.494,7.465) (7.6,7.466) (7.706,7.467) (7.811,7.468) (7.917,7.469) (8.022,7.469) (8.128,7.47) (8.233,7.471) (8.339,7.471) (8.444,7.472) (8.55,7.472) (8.656,7.473) (8.761,7.473) (8.867,7.474) (8.972,7.474) (9.078,7.474) (9.183,7.474) (9.289,7.474) (9.394,7.475)};
\draw[BeachOrange,dashed,line width=1.2pt] plot coordinates {(1.161,7.542) (1.267,7.579) (1.372,7.614) (1.478,7.646) (1.583,7.676) (1.689,7.704) (1.794,7.731) (1.9,7.756) (2.006,7.78) (2.111,7.803) (2.217,7.825) (2.322,7.845) (2.428,7.865) (2.533,7.885) (2.639,7.903) (2.744,7.921) (2.85,7.938) (2.956,7.954) (3.061,7.97) (3.167,7.986) (3.272,8) (3.378,8.015) (3.483,8.029) (3.589,8.042) (3.694,8.056) (3.8,8.068) (3.906,8.081) (4.011,8.093) (4.117,8.105) (4.222,8.116) (4.328,8.127) (4.433,8.138) (4.539,8.148) (4.644,8.159) (4.75,8.168) (4.856,8.178) (4.961,8.187) (5.067,8.197) (5.172,8.205) (5.278,8.214) (5.383,8.222) (5.489,8.23) (5.594,8.238) (5.7,8.246) (5.806,8.253) (5.911,8.261) (6.017,8.268) (6.122,8.274) (6.228,8.281) (6.333,8.287) (6.439,8.293) (6.544,8.299) (6.65,8.305) (6.756,8.31) (6.861,8.315) (6.967,8.321) (7.072,8.325) (7.178,8.33) (7.283,8.334) (7.389,8.339) (7.494,8.343) (7.6,8.346) (7.706,8.35) (7.811,8.353) (7.917,8.356) (8.022,8.359) (8.128,8.362) (8.233,8.365) (8.339,8.367) (8.444,8.369) (8.55,8.371) (8.656,8.373) (8.761,8.374) (8.867,8.376) (8.972,8.377) (9.078,8.378) (9.183,8.379) (9.289,8.379) (9.394,8.379)};
\draw[BeachTeal,dash dot,line width=1.2pt] plot coordinates {(0.528,5.181) (0.633,5.58) (0.739,5.921) (0.844,6.219) (0.95,6.488) (1.056,6.739) (1.161,6.986)};
\end{scope}
\draw[black!10,line width=0.4pt] (0,0)--(0,3.6);
\draw[black!10,line width=0.4pt] (1.583,0)--(1.583,3.6);
\draw[black!10,line width=0.4pt] (3.167,0)--(3.167,3.6);
\draw[black!10,line width=0.4pt] (4.75,0)--(4.75,3.6);
\draw[black!10,line width=0.4pt] (6.333,0)--(6.333,3.6);
\draw[black!10,line width=0.4pt] (7.917,0)--(7.917,3.6);
\draw[black!10,line width=0.4pt] (9.5,0)--(9.5,3.6);
\draw[black!10,line width=0.4pt] (0,0)--(9.5,0);
\draw[black!10,line width=0.4pt] (0,1.011)--(9.5,1.011);
\draw[black!10,line width=0.4pt] (0,2.119)--(9.5,2.119);
\draw[black!10,line width=0.4pt] (0,3.13)--(9.5,3.13);
\draw[line width=0.6pt] (0,0) rectangle (9.5,3.6);
\draw[line width=0.6pt] (0,0)--++(0,0.1);
\node[below,font=\scriptsize\sffamily] at (0,0) {$0$};
\draw[line width=0.6pt] (1.583,0)--++(0,0.1);
\node[below,font=\scriptsize\sffamily] at (1.583,0) {$15$};
\draw[line width=0.6pt] (3.167,0)--++(0,0.1);
\node[below,font=\scriptsize\sffamily] at (3.167,0) {$30$};
\draw[line width=0.6pt] (4.75,0)--++(0,0.1);
\node[below,font=\scriptsize\sffamily] at (4.75,0) {$45$};
\draw[line width=0.6pt] (6.333,0)--++(0,0.1);
\node[below,font=\scriptsize\sffamily] at (6.333,0) {$60$};
\draw[line width=0.6pt] (7.917,0)--++(0,0.1);
\node[below,font=\scriptsize\sffamily] at (7.917,0) {$75$};
\draw[line width=0.6pt] (9.5,0)--++(0,0.1);
\node[below,font=\scriptsize\sffamily] at (9.5,0) {$90$};
\draw[line width=0.6pt] (0,0)--++(0.1,0);
\node[left,font=\scriptsize\sffamily] at (0,0) {$0.1$};
\draw[line width=0.6pt] (0,1.011)--++(0.1,0);
\node[left,font=\scriptsize\sffamily] at (0,1.011) {$0.3$};
\draw[line width=0.6pt] (0,2.119)--++(0.1,0);
\node[left,font=\scriptsize\sffamily] at (0,2.119) {$1$};
\draw[line width=0.6pt] (0,3.13)--++(0.1,0);
\node[left,font=\scriptsize\sffamily] at (0,3.13) {$3$};
\node[font=\small\sffamily] at (4.75,-0.75) {太陽高度 $\alpha$ [deg]};
\node[rotate=90,font=\small\sffamily] at (-1,1.8) {数流束 $/(n_i v_{\mathrm{sw}})$};
\begin{scope}
\clip (0,0) rectangle (9.5,3.6);
\draw[BeachBlue,line width=1.2pt] plot coordinates {(1.267,2.116) (1.372,2.125) (1.478,2.134) (1.583,2.143) (1.689,2.151) (1.794,2.16) (1.9,2.168) (2.006,2.175) (2.111,2.183) (2.217,2.191) (2.322,2.198) (2.428,2.206) (2.533,2.213) (2.639,2.22) (2.744,2.227) (2.85,2.234) (2.956,2.241) (3.061,2.248) (3.167,2.255) (3.272,2.262) (3.378,2.268) (3.483,2.275) (3.589,2.281) (3.694,2.287) (3.8,2.293) (3.906,2.299) (4.011,2.305) (4.117,2.311) (4.222,2.317) (4.328,2.323) (4.433,2.328) (4.539,2.334) (4.644,2.339) (4.75,2.345) (4.856,2.35) (4.961,2.355) (5.067,2.36) (5.172,2.365) (5.278,2.369) (5.383,2.374) (5.489,2.379) (5.594,2.383) (5.7,2.387) (5.806,2.391) (5.911,2.396) (6.017,2.4) (6.122,2.403) (6.228,2.407) (6.333,2.411) (6.439,2.414) (6.544,2.418) (6.65,2.421) (6.756,2.424) (6.861,2.427) (6.967,2.43) (7.072,2.433) (7.178,2.436) (7.283,2.438) (7.389,2.441) (7.494,2.443) (7.6,2.445) (7.706,2.447) (7.811,2.449) (7.917,2.451) (8.022,2.453) (8.128,2.455) (8.233,2.456) (8.339,2.458) (8.444,2.459) (8.55,2.46) (8.656,2.461) (8.761,2.462) (8.867,2.463) (8.972,2.464) (9.078,2.464) (9.183,2.465) (9.289,2.465) (9.394,2.465)};
\draw[BeachInk,line width=1.2pt] plot coordinates {(1.267,0.674) (1.372,0.746) (1.478,0.813) (1.583,0.875) (1.689,0.933) (1.794,0.987) (1.9,1.038) (2.006,1.086) (2.111,1.132) (2.217,1.175) (2.322,1.215) (2.428,1.254) (2.533,1.291) (2.639,1.326) (2.744,1.36) (2.85,1.392) (2.956,1.423) (3.061,1.453) (3.167,1.481) (3.272,1.508) (3.378,1.535) (3.483,1.56) (3.589,1.584) (3.694,1.607) (3.8,1.63) (3.906,1.652) (4.011,1.673) (4.117,1.693) (4.222,1.712) (4.328,1.731) (4.433,1.749) (4.539,1.767) (4.644,1.784) (4.75,1.8) (4.856,1.816) (4.961,1.831) (5.067,1.846) (5.172,1.86) (5.278,1.874) (5.383,1.887) (5.489,1.9) (5.594,1.912) (5.7,1.924) (5.806,1.935) (5.911,1.946) (6.017,1.957) (6.122,1.967) (6.228,1.977) (6.333,1.987) (6.439,1.996) (6.544,2.004) (6.65,2.013) (6.756,2.021) (6.861,2.028) (6.967,2.036) (7.072,2.043) (7.178,2.049) (7.283,2.056) (7.389,2.062) (7.494,2.067) (7.6,2.073) (7.706,2.078) (7.811,2.083) (7.917,2.087) (8.022,2.091) (8.128,2.095) (8.233,2.099) (8.339,2.102) (8.444,2.105) (8.55,2.108) (8.656,2.11) (8.761,2.112) (8.867,2.114) (8.972,2.115) (9.078,2.117) (9.183,2.118) (9.289,2.118) (9.394,2.119)};
\draw[BeachOrange,line width=1.2pt] plot coordinates {(1.267,1.9) (1.372,1.893) (1.478,1.884) (1.583,1.876) (1.689,1.867) (1.794,1.858) (1.9,1.848) (2.006,1.839) (2.111,1.83) (2.217,1.82) (2.322,1.811) (2.428,1.801) (2.533,1.792) (2.639,1.783) (2.744,1.773) (2.85,1.764) (2.956,1.754) (3.061,1.745) (3.167,1.736) (3.272,1.726) (3.378,1.717) (3.483,1.708) (3.589,1.699) (3.694,1.69) (3.8,1.68) (3.906,1.672) (4.011,1.663) (4.117,1.654) (4.222,1.645) (4.328,1.636) (4.433,1.628) (4.539,1.619) (4.644,1.611) (4.75,1.603) (4.856,1.595) (4.961,1.587) (5.067,1.579) (5.172,1.571) (5.278,1.563) (5.383,1.556) (5.489,1.548) (5.594,1.541) (5.7,1.534) (5.806,1.527) (5.911,1.52) (6.017,1.513) (6.122,1.507) (6.228,1.5) (6.333,1.494) (6.439,1.488) (6.544,1.482) (6.65,1.476) (6.756,1.471) (6.861,1.466) (6.967,1.46) (7.072,1.455) (7.178,1.451) (7.283,1.446) (7.389,1.442) (7.494,1.437) (7.6,1.433) (7.706,1.43) (7.811,1.426) (7.917,1.423) (8.022,1.419) (8.128,1.416) (8.233,1.414) (8.339,1.411) (8.444,1.409) (8.55,1.407) (8.656,1.405) (8.761,1.403) (8.867,1.402) (8.972,1.4) (9.078,1.399) (9.183,1.399) (9.289,1.398) (9.394,1.398)};
\draw[BeachOrange,dashed,line width=0.9pt] plot coordinates {(1.267,1.926) (1.372,1.999) (1.478,2.066) (1.583,2.128) (1.689,2.186) (1.794,2.24) (1.9,2.291) (2.006,2.339) (2.111,2.384) (2.217,2.427) (2.322,2.468) (2.428,2.507) (2.533,2.544) (2.639,2.579) (2.744,2.613) (2.85,2.645) (2.956,2.676) (3.061,2.705) (3.167,2.734) (3.272,2.761) (3.378,2.787) (3.483,2.812) (3.589,2.837) (3.694,2.86) (3.8,2.883) (3.906,2.904) (4.011,2.925) (4.117,2.945) (4.222,2.965) (4.328,2.984) (4.433,3.002) (4.539,3.019) (4.644,3.036) (4.75,3.053) (4.856,3.068) (4.961,3.084) (5.067,3.098) (5.172,3.113) (5.278,3.126) (5.383,3.14) (5.489,3.152) (5.594,3.165) (5.7,3.176) (5.806,3.188) (5.911,3.199) (6.017,3.21) (6.122,3.22) (6.228,3.23) (6.333,3.239) (6.439,3.248) (6.544,3.257) (6.65,3.265) (6.756,3.273) (6.861,3.281) (6.967,3.288) (7.072,3.295) (7.178,3.302) (7.283,3.308) (7.389,3.314) (7.494,3.32) (7.6,3.325) (7.706,3.33) (7.811,3.335) (7.917,3.34) (8.022,3.344) (8.128,3.348) (8.233,3.351) (8.339,3.354) (8.444,3.357) (8.55,3.36) (8.656,3.363) (8.761,3.365) (8.867,3.366) (8.972,3.368) (9.078,3.369) (9.183,3.37) (9.289,3.371) (9.394,3.371)};
\end{scope}
\draw[BeachBlue,solid,line width=1.2pt] (9.9,8.3)--(10.5,8.3);
\node[right,font=\scriptsize\sffamily] at (10.5,8.3) {Type A の $\phi_0$};
\draw[BeachOrange,dashed,line width=1.2pt] (9.9,7.84)--(10.5,7.84);
\node[right,font=\scriptsize\sffamily] at (10.5,7.84) {Type B の $\phi_0$};
\draw[BeachTeal,dash dot,line width=1.2pt] (9.9,7.38)--(10.5,7.38);
\node[right,font=\scriptsize\sffamily] at (10.5,7.38) {Type C の $\phi_0$};
\draw[BeachBlue,dotted,line width=1.2pt] (9.9,6.92)--(10.5,6.92);
\node[right,font=\scriptsize\sffamily] at (10.5,6.92) {Type A の $\phi_m$};
\draw[BeachBlue,line width=1.2pt] (9.9,3.3)--(10.5,3.3);
\node[right,font=\scriptsize\sffamily] at (10.5,3.3) {太陽風電子 $\Gamma_e$};
\draw[BeachInk,line width=1.2pt] (9.9,2.84)--(10.5,2.84);
\node[right,font=\scriptsize\sffamily] at (10.5,2.84) {イオン $\Gamma_i$};
\draw[BeachOrange,line width=1.2pt] (9.9,2.38)--(10.5,2.38);
\node[right,font=\scriptsize\sffamily] at (10.5,2.38) {脱出光電子 $\Gamma_{\mathrm{esc}}$};
\draw[BeachOrange,dashed,line width=0.9pt] (9.9,1.92)--(10.5,1.92);
\node[right,font=\scriptsize\sffamily] at (10.5,1.92) {放出光電子 $\Gamma_{\mathrm{emit}}$};
\node[right,font=\scriptsize\sffamily] at (9.9,1.2) {下段は Type A の根};
\end{tikzpicture}
\caption[太陽高度に対する零電流根]{太陽高度 $\alpha$ を $1^\circ$ 刻みで変えたときの零電流根。上段は壁電位 $\phi_0$ と Type A の極小 $\phi_m$、下段は Type A の根での壁への数流束である。各型を個別に指定し、直前の高度の根を初期値に加えて解いた。曲線の切れ目は、有限の探索で物理解が見つからなかった高度である（不存在の証明ではない）。下段では $\Gamma_e=\Gamma_i+\Gamma_{\mathrm{esc}}$ が零電流条件を表し、放出の大部分は壁へ戻る。}
\label{fig:sm_alpha_sweep}
\end{figure}


\subsection{型の分布}
図\ref{fig:sm_type_map}は、太陽高度と光電子源の強さの平面で、零電流の物理解が見つかった型である。
光電子源が弱いか太陽高度が低いと Type C だけが見つかり、壁は負に帯電する。光電子源が強くなると、
B と C が同時に見つかる狭い帯を経て、A と B が同時に見つかる広い領域に移る。この計算の範囲では、A だけ、または B だけが見つかる点はなかった。
$\alpha\le4^\circ$ では光電子源の強さによらず根が見つからず、境界は前節の Bohm 速度の目安と一致する。
複数の型が見つかる条件では、本モデルはどれが実現するかを決めない。
研究で外部シースを使うときは、型を明示して結果の依存性を確かめる必要がある。

% Generated by tools/sheath_model_figures.py from type_map.csv; edit the generator, not this file.
\begin{figure}[htbp]
\centering
\begin{tikzpicture}[beach diagram]
\fill[black!8] (0.056,0) rectangle (0.513,0.078);
\fill[BeachTeal!70] (0.5,0) rectangle (9.958,0.078);
\fill[black!8] (0.056,0.063) rectangle (0.513,0.203);
\fill[BeachTeal!70] (0.5,0.063) rectangle (9.958,0.203);
\fill[black!8] (0.056,0.187) rectangle (0.513,0.327);
\fill[BeachTeal!70] (0.5,0.187) rectangle (9.958,0.327);
\fill[black!8] (0.056,0.312) rectangle (0.513,0.452);
\fill[BeachTeal!70] (0.5,0.312) rectangle (9.958,0.452);
\fill[black!8] (0.056,0.437) rectangle (0.513,0.577);
\fill[BeachTeal!70] (0.5,0.437) rectangle (9.958,0.577);
\fill[black!8] (0.056,0.562) rectangle (0.513,0.702);
\fill[BeachTeal!70] (0.5,0.562) rectangle (9.958,0.702);
\fill[black!8] (0.056,0.688) rectangle (0.513,0.828);
\fill[BeachTeal!70] (0.5,0.688) rectangle (9.958,0.828);
\fill[black!8] (0.056,0.812) rectangle (0.513,0.952);
\fill[BeachTeal!70] (0.5,0.812) rectangle (9.958,0.952);
\fill[black!8] (0.056,0.937) rectangle (0.513,1.077);
\fill[BeachTeal!70] (0.5,0.937) rectangle (8.291,1.077);
\fill[BeachRed!65] (8.278,0.937) rectangle (8.736,1.077);
\fill[BeachMagenta!65] (8.722,0.937) rectangle (9.958,1.077);
\fill[black!8] (0.056,1.063) rectangle (0.513,1.203);
\fill[BeachTeal!70] (0.5,1.063) rectangle (7.402,1.203);
\fill[BeachRed!65] (7.389,1.063) rectangle (7.847,1.203);
\fill[BeachMagenta!65] (7.833,1.063) rectangle (9.958,1.203);
\fill[black!8] (0.056,1.188) rectangle (0.513,1.328);
\fill[BeachTeal!70] (0.5,1.188) rectangle (6.847,1.328);
\fill[BeachRed!65] (6.833,1.188) rectangle (7.18,1.328);
\fill[BeachMagenta!65] (7.167,1.188) rectangle (9.958,1.328);
\fill[black!8] (0.056,1.312) rectangle (0.513,1.452);
\fill[BeachTeal!70] (0.5,1.312) rectangle (6.291,1.452);
\fill[BeachRed!65] (6.278,1.312) rectangle (6.624,1.452);
\fill[BeachMagenta!65] (6.611,1.312) rectangle (9.958,1.452);
\fill[black!8] (0.056,1.438) rectangle (0.513,1.578);
\fill[BeachTeal!70] (0.5,1.438) rectangle (5.847,1.578);
\fill[BeachRed!65] (5.833,1.438) rectangle (6.18,1.578);
\fill[BeachMagenta!65] (6.167,1.438) rectangle (9.958,1.578);
\fill[black!8] (0.056,1.562) rectangle (0.513,1.702);
\fill[BeachTeal!70] (0.5,1.562) rectangle (5.513,1.702);
\fill[BeachRed!65] (5.5,1.562) rectangle (5.847,1.702);
\fill[BeachMagenta!65] (5.833,1.562) rectangle (9.958,1.702);
\fill[black!8] (0.056,1.688) rectangle (0.513,1.828);
\fill[BeachTeal!70] (0.5,1.688) rectangle (5.18,1.828);
\fill[BeachRed!65] (5.167,1.688) rectangle (5.513,1.828);
\fill[BeachMagenta!65] (5.5,1.688) rectangle (9.958,1.828);
\fill[black!8] (0.056,1.813) rectangle (0.513,1.953);
\fill[BeachTeal!70] (0.5,1.813) rectangle (4.847,1.953);
\fill[BeachRed!65] (4.833,1.813) rectangle (5.18,1.953);
\fill[BeachMagenta!65] (5.167,1.813) rectangle (9.958,1.953);
\fill[black!8] (0.056,1.937) rectangle (0.513,2.077);
\fill[BeachTeal!70] (0.5,1.937) rectangle (4.513,2.077);
\fill[BeachRed!65] (4.5,1.937) rectangle (4.847,2.077);
\fill[BeachMagenta!65] (4.833,1.937) rectangle (9.958,2.077);
\fill[black!8] (0.056,2.062) rectangle (0.513,2.202);
\fill[BeachTeal!70] (0.5,2.062) rectangle (4.291,2.202);
\fill[BeachRed!65] (4.278,2.062) rectangle (4.513,2.202);
\fill[BeachMagenta!65] (4.5,2.062) rectangle (9.958,2.202);
\fill[black!8] (0.056,2.188) rectangle (0.513,2.328);
\fill[BeachTeal!70] (0.5,2.188) rectangle (3.958,2.328);
\fill[BeachRed!65] (3.944,2.188) rectangle (4.291,2.328);
\fill[BeachMagenta!65] (4.278,2.188) rectangle (9.958,2.328);
\fill[black!8] (0.056,2.313) rectangle (0.513,2.453);
\fill[BeachTeal!70] (0.5,2.313) rectangle (3.736,2.453);
\fill[BeachRed!65] (3.722,2.313) rectangle (4.069,2.453);
\fill[BeachMagenta!65] (4.056,2.313) rectangle (9.958,2.453);
\fill[black!8] (0.056,2.438) rectangle (0.513,2.578);
\fill[BeachTeal!70] (0.5,2.438) rectangle (3.513,2.578);
\fill[BeachRed!65] (3.5,2.438) rectangle (3.847,2.578);
\fill[BeachMagenta!65] (3.833,2.438) rectangle (9.958,2.578);
\fill[black!8] (0.056,2.563) rectangle (0.513,2.703);
\fill[BeachTeal!70] (0.5,2.563) rectangle (3.291,2.703);
\fill[BeachRed!65] (3.278,2.563) rectangle (3.624,2.703);
\fill[BeachMagenta!65] (3.611,2.563) rectangle (9.958,2.703);
\fill[black!8] (0.056,2.688) rectangle (0.513,2.828);
\fill[BeachTeal!70] (0.5,2.688) rectangle (3.069,2.828);
\fill[BeachRed!65] (3.056,2.688) rectangle (3.402,2.828);
\fill[BeachMagenta!65] (3.389,2.688) rectangle (9.958,2.828);
\fill[black!8] (0.056,2.812) rectangle (0.513,2.952);
\fill[BeachTeal!70] (0.5,2.812) rectangle (2.847,2.952);
\fill[BeachRed!65] (2.833,2.812) rectangle (3.18,2.952);
\fill[BeachMagenta!65] (3.167,2.812) rectangle (9.958,2.952);
\fill[black!8] (0.056,2.937) rectangle (0.513,3.077);
\fill[BeachTeal!70] (0.5,2.937) rectangle (2.736,3.077);
\fill[BeachRed!65] (2.722,2.937) rectangle (2.958,3.077);
\fill[BeachMagenta!65] (2.944,2.937) rectangle (9.958,3.077);
\fill[black!8] (0.056,3.062) rectangle (0.513,3.202);
\fill[BeachTeal!70] (0.5,3.062) rectangle (2.513,3.202);
\fill[BeachRed!65] (2.5,3.062) rectangle (2.847,3.202);
\fill[BeachMagenta!65] (2.833,3.062) rectangle (9.958,3.202);
\fill[black!8] (0.056,3.188) rectangle (0.513,3.328);
\fill[BeachTeal!70] (0.5,3.188) rectangle (2.402,3.328);
\fill[BeachRed!65] (2.389,3.188) rectangle (2.624,3.328);
\fill[BeachMagenta!65] (2.611,3.188) rectangle (9.958,3.328);
\fill[black!8] (0.056,3.312) rectangle (0.513,3.452);
\fill[BeachTeal!70] (0.5,3.312) rectangle (2.291,3.452);
\fill[BeachRed!65] (2.278,3.312) rectangle (2.513,3.452);
\fill[BeachMagenta!65] (2.5,3.312) rectangle (9.958,3.452);
\fill[black!8] (0.056,3.437) rectangle (0.513,3.577);
\fill[BeachTeal!70] (0.5,3.437) rectangle (2.069,3.577);
\fill[BeachRed!65] (2.056,3.437) rectangle (2.402,3.577);
\fill[BeachMagenta!65] (2.389,3.437) rectangle (9.958,3.577);
\fill[black!8] (0.056,3.563) rectangle (0.513,3.703);
\fill[BeachTeal!70] (0.5,3.563) rectangle (1.958,3.703);
\fill[BeachRed!65] (1.944,3.563) rectangle (2.18,3.703);
\fill[BeachMagenta!65] (2.167,3.563) rectangle (9.958,3.703);
\fill[black!8] (0.056,3.687) rectangle (0.513,3.827);
\fill[BeachTeal!70] (0.5,3.687) rectangle (1.847,3.827);
\fill[BeachRed!65] (1.833,3.687) rectangle (2.069,3.827);
\fill[BeachMagenta!65] (2.056,3.687) rectangle (9.958,3.827);
\fill[black!8] (0.056,3.813) rectangle (0.513,3.953);
\fill[BeachTeal!70] (0.5,3.813) rectangle (1.736,3.953);
\fill[BeachRed!65] (1.722,3.813) rectangle (1.958,3.953);
\fill[BeachMagenta!65] (1.944,3.813) rectangle (9.958,3.953);
\fill[black!8] (0.056,3.938) rectangle (0.513,4.078);
\fill[BeachTeal!70] (0.5,3.938) rectangle (1.624,4.078);
\fill[BeachRed!65] (1.611,3.938) rectangle (1.847,4.078);
\fill[BeachMagenta!65] (1.833,3.938) rectangle (9.958,4.078);
\fill[black!8] (0.056,4.063) rectangle (0.513,4.203);
\fill[BeachTeal!70] (0.5,4.063) rectangle (1.513,4.203);
\fill[BeachRed!65] (1.5,4.063) rectangle (1.736,4.203);
\fill[BeachMagenta!65] (1.722,4.063) rectangle (9.958,4.203);
\fill[black!8] (0.056,4.187) rectangle (0.513,4.327);
\fill[BeachTeal!70] (0.5,4.187) rectangle (1.402,4.327);
\fill[BeachRed!65] (1.389,4.187) rectangle (1.624,4.327);
\fill[BeachMagenta!65] (1.611,4.187) rectangle (9.958,4.327);
\fill[black!8] (0.056,4.312) rectangle (0.513,4.452);
\fill[BeachTeal!70] (0.5,4.312) rectangle (1.291,4.452);
\fill[BeachRed!65] (1.278,4.312) rectangle (1.513,4.452);
\fill[BeachMagenta!65] (1.5,4.312) rectangle (9.958,4.452);
\fill[black!8] (0.056,4.437) rectangle (0.513,4.577);
\fill[BeachTeal!70] (0.5,4.437) rectangle (1.291,4.577);
\fill[BeachRed!65] (1.278,4.437) rectangle (1.402,4.577);
\fill[BeachMagenta!65] (1.389,4.437) rectangle (9.958,4.577);
\fill[black!8] (0.056,4.563) rectangle (0.513,4.703);
\fill[BeachTeal!70] (0.5,4.563) rectangle (1.18,4.703);
\fill[BeachRed!65] (1.167,4.563) rectangle (1.291,4.703);
\fill[BeachMagenta!65] (1.278,4.563) rectangle (9.958,4.703);
\fill[black!8] (0.056,4.688) rectangle (0.513,4.828);
\fill[BeachTeal!70] (0.5,4.688) rectangle (1.069,4.828);
\fill[BeachRed!65] (1.056,4.688) rectangle (1.291,4.828);
\fill[BeachMagenta!65] (1.278,4.688) rectangle (9.958,4.828);
\fill[black!8] (0.056,4.813) rectangle (0.513,4.953);
\fill[BeachTeal!70] (0.5,4.813) rectangle (1.069,4.953);
\fill[BeachRed!65] (1.056,4.813) rectangle (1.18,4.953);
\fill[BeachMagenta!65] (1.167,4.813) rectangle (9.958,4.953);
\fill[black!8] (0.056,4.938) rectangle (0.513,5.078);
\fill[BeachTeal!70] (0.5,4.938) rectangle (0.958,5.078);
\fill[BeachRed!65] (0.944,4.938) rectangle (1.069,5.078);
\fill[BeachMagenta!65] (1.056,4.938) rectangle (9.958,5.078);
\fill[black!8] (0.056,5.063) rectangle (0.513,5.203);
\fill[BeachTeal!70] (0.5,5.063) rectangle (0.847,5.203);
\fill[BeachRed!65] (0.833,5.063) rectangle (1.069,5.203);
\fill[BeachMagenta!65] (1.056,5.063) rectangle (9.958,5.203);
\fill[black!8] (0.056,5.187) rectangle (0.513,5.327);
\fill[BeachTeal!70] (0.5,5.187) rectangle (0.847,5.327);
\fill[BeachRed!65] (0.833,5.187) rectangle (0.958,5.327);
\fill[BeachMagenta!65] (0.944,5.187) rectangle (9.958,5.327);
\fill[black!8] (0.056,5.312) rectangle (0.513,5.452);
\fill[BeachTeal!70] (0.5,5.312) rectangle (0.736,5.452);
\fill[BeachRed!65] (0.722,5.312) rectangle (0.847,5.452);
\fill[BeachMagenta!65] (0.833,5.312) rectangle (9.958,5.452);
\fill[black!8] (0.056,5.437) rectangle (0.513,5.577);
\fill[BeachTeal!70] (0.5,5.437) rectangle (0.736,5.577);
\fill[BeachRed!65] (0.722,5.437) rectangle (0.847,5.577);
\fill[BeachMagenta!65] (0.833,5.437) rectangle (9.958,5.577);
\fill[black!8] (0.056,5.563) rectangle (0.513,5.703);
\fill[BeachTeal!70] (0.5,5.563) rectangle (0.624,5.703);
\fill[BeachRed!65] (0.611,5.563) rectangle (0.736,5.703);
\fill[BeachMagenta!65] (0.722,5.563) rectangle (9.958,5.703);
\fill[black!8] (0.056,5.688) rectangle (0.513,5.828);
\fill[BeachTeal!70] (0.5,5.688) rectangle (0.624,5.828);
\fill[BeachRed!65] (0.611,5.688) rectangle (0.736,5.828);
\fill[BeachMagenta!65] (0.722,5.688) rectangle (9.958,5.828);
\fill[black!8] (0.056,5.813) rectangle (0.513,5.953);
\fill[BeachTeal!70] (0.5,5.813) rectangle (0.624,5.953);
\fill[BeachMagenta!65] (0.611,5.813) rectangle (9.958,5.953);
\fill[black!8] (0.056,5.937) rectangle (0.513,6);
\fill[BeachMagenta!65] (0.5,5.937) rectangle (9.958,6);
\draw[line width=0.6pt] (0,0) rectangle (10,6);
\draw[line width=0.6pt] (0,0)--++(0,0.1);
\node[below,font=\scriptsize\sffamily] at (0,0) {$0$};
\draw[line width=0.6pt] (1.667,0)--++(0,0.1);
\node[below,font=\scriptsize\sffamily] at (1.667,0) {$15$};
\draw[line width=0.6pt] (3.333,0)--++(0,0.1);
\node[below,font=\scriptsize\sffamily] at (3.333,0) {$30$};
\draw[line width=0.6pt] (5,0)--++(0,0.1);
\node[below,font=\scriptsize\sffamily] at (5,0) {$45$};
\draw[line width=0.6pt] (6.667,0)--++(0,0.1);
\node[below,font=\scriptsize\sffamily] at (6.667,0) {$60$};
\draw[line width=0.6pt] (8.333,0)--++(0,0.1);
\node[below,font=\scriptsize\sffamily] at (8.333,0) {$75$};
\draw[line width=0.6pt] (10,0)--++(0,0.1);
\node[below,font=\scriptsize\sffamily] at (10,0) {$90$};
\draw[line width=0.6pt] (0,0)--++(0.1,0);
\node[left,font=\scriptsize\sffamily] at (0,0) {$0.5$};
\draw[line width=0.6pt] (0,1.2)--++(0.1,0);
\node[left,font=\scriptsize\sffamily] at (0,1.2) {$1$};
\draw[line width=0.6pt] (0,2.4)--++(0.1,0);
\node[left,font=\scriptsize\sffamily] at (0,2.4) {$2$};
\draw[line width=0.6pt] (0,3.6)--++(0.1,0);
\node[left,font=\scriptsize\sffamily] at (0,3.6) {$4$};
\draw[line width=0.6pt] (0,4.8)--++(0.1,0);
\node[left,font=\scriptsize\sffamily] at (0,4.8) {$8$};
\draw[line width=0.6pt] (0,6)--++(0.1,0);
\node[left,font=\scriptsize\sffamily] at (0,6) {$16$};
\node[font=\small\sffamily] at (5,-0.75) {太陽高度 $\alpha$ [deg]};
\node[rotate=90,font=\small\sffamily] at (-1.05,3) {光電子源の強さ $n_{\mathrm{pe,ref}}/n_i$};
\node[draw=white,fill=BeachInk,circle,inner sep=1.6pt,line width=0.6pt] at (6.667,4.655) {};
\node[font=\scriptsize\sffamily\bfseries,text=white] at (6.667,4.999) {A};
\node[draw=white,fill=BeachInk,circle,inner sep=1.6pt,line width=0.6pt] at (2.222,4.655) {};
\node[font=\scriptsize\sffamily\bfseries,text=white] at (2.222,4.999) {B};
\node[draw=white,fill=BeachInk,circle,inner sep=1.6pt,line width=0.6pt] at (1.111,4.655) {};
\node[font=\scriptsize\sffamily\bfseries,text=white] at (1.111,4.999) {C};
\draw[white,dashed,line width=0.5pt] (0,4.655)--(10,4.655);
\fill[BeachTeal!70] (10.5,5.65) rectangle (10.9,5.95);
\node[right,font=\scriptsize\sffamily] at (10.95,5.8) {C};
\fill[BeachMagenta!65] (10.5,5.2) rectangle (10.9,5.5);
\node[right,font=\scriptsize\sffamily] at (10.95,5.35) {A+B};
\fill[BeachRed!65] (10.5,4.75) rectangle (10.9,5.05);
\node[right,font=\scriptsize\sffamily] at (10.95,4.9) {B+C};
\fill[black!8] (10.5,4.3) rectangle (10.9,4.6);
\node[right,font=\scriptsize\sffamily] at (10.95,4.45) {なし};
\end{tikzpicture}
\caption[零電流の物理解が見つかる型の分布]{太陽高度 $\alpha$ と光電子源の強さ $n_{\mathrm{pe,ref}}/n_i$ の平面で、零電流の物理解が見つかった型。格子は 89$\times$49 点で、各点で A、B、C を個別に指定し、既定の初期値だけから解いた。他の条件は表\ref{tab:sm_example_inputs}と同じで、白い破線と点は図\ref{fig:sm_profiles}の例の位置である。「A+B」などは同じ条件で複数の型が見つかったことを表し、色は解の安定性による選択を表さない。「なし」は有限の探索で物理解が見つからなかった点で、解の不存在を保証しない。}
\label{fig:sm_type_map}
\end{figure}


\section{適用範囲}
\label{sec:sm_limits}
本モデルは平面・定常・無衝突・非磁化のシースであり、曲率、時間変化、衝突、磁場による旋回、波動や不安定性による緩和を含まない。
壁は完全に吸収し、二次電子放出、後方散乱、表面伝導を含まない。イオンは一定係数の圧力を持つ流体で、イオンの速度分布は定めない。
上流は半無限で準中性とし、presheath や上流の加速機構とは結合しない。
入射状態と電子の drift は固定した物理条件であり、根が見つからないときにそれらを自動で変更しない。

代数的な根の収束、物理的な電位分布の存在、動的な安定性は別の概念である。本モデルは前の二つを判定するが、
複数の根の中から安定なものを選ぶ判定は含まない。根の探索と物理解の判定はどちらも有限の計算であり、
見つからないことは解の不存在の証明ではない。drift のある電子で受理した根は、上流の帯の内側で理想化した分布に従わない。
これらが結論を左右するケースでは、独立の運動論的計算や PIC と比較する。
